% Leçon 20 — Grammaire : 比较字句:……跟/和…. (不)一样et ……跟/和…. (不)一样+ 形容词 ; 一方面… ,另一方面… ; 如果…, 但是/可是… — Chinois 1ère A — Cours signé OnBuch+
% Programme officiel MINESEC. Compiler avec : tectonic ZHA20-grammaire-et.tex
\newcommand{\DOCMATIERE}{Chinois}
\newcommand{\DOCNIVEAU}{1ère A}
\newcommand{\DOCMODULE}{Module 3 : Citoyenneté et ouverture au monde}
\newcommand{\DOCLECON}{ZHA20}
\newcommand{\DOCTITRE}{Leçon 20 — Grammaire : 比较字句:……跟/和…. (不)一样et ……跟/和…. (不)一样+ 形容词 ; 一方面… ,另一方面… ; 如果…, 但是/可是…}
% OnBuch+ — préambule « cours signés » ENRICHI (compilé avec tectonic / XeLaTeX)
% Système visuel « Pop » d'OnBuch+ : fond crème (jamais blanc), bordures encre
% épaisses, ombres « dures » décalées, titres Archivo Black en capitales.
% Version MATHÉMATIQUES : graphes (pgfplots/tikz), boîtes pédagogiques
% (définition, propriété, méthode, attention, le savais-tu, exercice résolu,
% exercice, TP, bilan), page de garde et signature OnBuch+ ancrée en bas.
%
% Macros attendues dans main.tex AVANT \input{preamble} :
%   \DOCMATIERE  \DOCNIVEAU  \DOCMODULE  \DOCLECON (numéro)  \DOCTITRE
\documentclass[11pt]{article}

\usepackage{fontspec}
\usepackage[french]{babel} % français = langue principale ; le chinois via \zh{..} / textebox (xeCJK)
\usepackage{xeCJK}
\usepackage{pifont} % \ding{51} \ding{55} utilisés par les modèles
\usepackage[a4paper,margin=2cm,top=3.4cm,bottom=2.8cm,headheight=1.4cm,footskip=1.3cm]{geometry}
\usepackage{amsmath,amssymb,amsfonts}
\usepackage{mathtools} % \xmapsto, \coloneqq... souvent utilisés par les modèles
\usepackage{cancel} % simplifications barrées (bilans de forces)
% \abs{z} (module, valeur absolue) et \norm{u} (norme), utilisés par les modèles
\providecommand{\abs}{}\let\abs\relax\DeclarePairedDelimiter{\abs}{\lvert}{\rvert}
\providecommand{\norm}{}\let\norm\relax\DeclarePairedDelimiter{\norm}{\lVert}{\rVert}
\usepackage[table]{xcolor} % [table] : fournit aussi \rowcolors (alternance de lignes)
\usepackage{pagecolor}
\usepackage{tikz}
\usetikzlibrary{arrows.meta,positioning,calc,shapes.geometric,shapes.misc,shapes.arrows,decorations.pathmorphing,decorations.pathreplacing,decorations.markings,patterns,shadows,mindmap,trees,fit,backgrounds,matrix,intersections,angles,quotes}
\usepackage{pgfplots}
\usepackage[version=4]{mhchem} % équations nucléaires \ce{^{235}_{92}U}
\usepackage[european,siunitx]{circuitikz} % schémas électriques
\pgfplotsset{compat=1.18}
\usepgfplotslibrary{fillbetween,groupplots} % aires entre courbes (calcul intégral)
\usepackage{enumitem}
\usepackage{fancyhdr}
% [most] charge toutes les bibliothèques tcolorbox (listings, minted, raster,
% documentation...) et gonfle inutilement la mémoire TeX sur un document long
% (~300 boîtes) : on ne charge que ce qui est réellement utilisé.
\usepackage[skins,breakable]{tcolorbox}
\usepackage{graphicx}
\usepackage{colortbl}
\usepackage{booktabs}
\usepackage{tabularx}
\usepackage{diagbox} % case d'en-tête diagonale des tableaux à double entrée
% Colonnes extensibles centrée / gauche / droite (C, L, R), souvent utilisées par les modèles
\newcolumntype{C}{>{\centering\arraybackslash}X}
\newcolumntype{L}{>{\raggedright\arraybackslash}X}
\newcolumntype{R}{>{\raggedleft\arraybackslash}X}
\usepackage{array}
\usepackage{multicol}
\usepackage{multirow}
\usepackage{siunitx}
% parse-numbers=false : accepte aussi \SI{2,5 \times 10^{-3}}{...}
\sisetup{locale=FR,output-decimal-marker={,},group-minimum-digits=4,parse-numbers=false}
\DeclareSIUnit\ohm{\ensuremath{\symup{\Omega}}} % U+2126 absent de la police
\DeclareSIUnit\division{div} % réglages d'oscilloscope (ms/div, V/div)
\DeclareSIUnit\tr{tr} % tours (vitesse de rotation, ex. \SI{1500}{\tr\per\minute})
\DeclareSIUnit\molar{M} % concentration molaire (ex. \SI{3}{\molar}, \SI{1}{\micro\molar})
\DeclareSIUnit\an{an} % année (le modèle confond parfois avec \year, un compteur TeX) — ex. \SI{5}{\kilo\meter\per\an}
\DeclareSIUnit\FCFA{FCFA} % franc CFA (ex. \SI{500}{\FCFA\per\kilo\gram})
\DeclareSIUnit\degreeBrix{\textdegree Brix} % teneur en sucre d'un sirop/jus (réfractométrie)
\DeclareSIUnit\month{mois} % le modèle utilise parfois \month, un compteur TeX, comme unité de durée
\usepackage{qrcode}
% \degree : commande fréquemment utilisée par les modèles pour un angle ou une
% température en dehors de \SI{}{} (non fournie par les paquets chargés).
\providecommand{\degree}{^{\circ}}
% \qed (fin de démonstration), utilisé par les modèles sans amsthm
\providecommand{\qed}{\hfill$\square$}
% \barre : double barre des valeurs interdites dans un tableau de variations (tkz-tab)
\providecommand{\barre}{\ensuremath{\|}}
% environnement proof (amsthm non chargé) utilisé par les modèles
\ifdefined\proof\else\newenvironment{proof}[1][Démonstration]{\par\noindent\textit{#1.}\ }{\qed\par}\fi
\usepackage{lastpage}
\usepackage{hyperref}
\hypersetup{hidelinks}

% ============================================================
% Palette « Pop » OnBuch+ — mêmes hex que la classe P (plus_ui.dart)
% ============================================================
\definecolor{poporange}{HTML}{F59321}
\definecolor{popdark}{HTML}{E07A0C}
\definecolor{popcream}{HTML}{FFF6E8}
\definecolor{popink}{HTML}{1C1714}
\definecolor{poppurple}{HTML}{7A5AE0}
\definecolor{popgreen}{HTML}{2E9E6A}
\definecolor{poppink}{HTML}{E0457B}
\definecolor{popblue}{HTML}{2F7DE1}
\definecolor{popgold}{HTML}{C98A12}
\definecolor{popmuted}{HTML}{8A7F76}
\definecolor{popline}{HTML}{EADFD2}
\definecolor{popgray}{HTML}{8A7F76}
\colorlet{popgrey}{popgray}
% Alias tolérants (le modèle nomme parfois les couleurs différemment)
\colorlet{popwhite}{white}
\colorlet{popblack}{popink}
\colorlet{popred}{poppink}
\colorlet{popyellow}{popgold}
% Teintes claires (fonds de tableaux/encadrés) — le modèle nomme parfois
% les couleurs avec un suffixe L (ex. popblueL) sans les avoir définies.
\colorlet{poporangeL}{poporange!20}
\colorlet{popdarkL}{popdark!20}
\colorlet{poppurpleL}{poppurple!20}
\colorlet{popgreenL}{popgreen!20}
\colorlet{poppinkL}{poppink!20}
\colorlet{popblueL}{popblue!20}
\colorlet{popgoldL}{popgold!20}
\colorlet{popredL}{popred!20}
\colorlet{popgrayL}{popgray!20}
% Teintes claires pour fonds de tableaux / zones
\colorlet{popblueL}{popblue!10!white}
\colorlet{poporangeL}{poporange!14!white}
\colorlet{popgreenL}{popgreen!12!white}
\colorlet{poppurpleL}{poppurple!12!white}
\colorlet{poppinkL}{poppink!12!white}
\colorlet{popgoldL}{popgold!14!white}

\pagecolor{popcream}

% ============================================================
% Typographie
% ============================================================
\defaultfontfeatures{Ligatures=TeX}
\setmainfont{PlusJakartaSans-Medium.ttf}[Path=../../../fonts/,BoldFont=PlusJakartaSans-Bold.ttf]
% Symboles, API et emojis absents de la police principale : repli sur DejaVu Sans (caractères actifs)
\newfontface\symfont{DejaVuSans.ttf}[Path=../../../fonts/]
\makeatletter
\newcommand\symactive[1]{\catcode#1=13 \begingroup\lccode`\~=#1 \lowercase{\endgroup\def~}{{\symfont\char#1}}}
\newcommand\symblank[1]{\catcode#1=13 \begingroup\lccode`\~=#1 \lowercase{\endgroup\def~}{}}
\newcommand\symhyph[1]{\catcode#1=13 \begingroup\lccode`\~=#1 \lowercase{\endgroup\def~}{\mbox{-}}}
\newcount\symc
\newcommand\symrange[2]{\symc=#1 \@whilenum\symc<#2 \do{\expandafter\symactive\expandafter{\the\symc}\advance\symc by 1 }}
\newcommand\symblankrange[2]{\symc=#1 \@whilenum\symc<#2 \do{\expandafter\symblank\expandafter{\the\symc}\advance\symc by 1 }}
\makeatother
\symrange{"0250}{"02FF}\symactive{"2713}\symactive{"2714}\symactive{"2717}\symactive{"2718}\symactive{"2610}\symactive{"2611}\symactive{"2612}
\symrange{"2600}{"27BF}
\catcode"2705=13 \begingroup\lccode`\~="2705 \lowercase{\endgroup\def~}{{\symfont\char"2713}}\symblank{"26BD}\symblank{"26A1}
\symrange{"2460}{"257F}\symactive{"223C}\symactive{"2248}\symactive{"2260}
\symactive{"01F8}\symactive{"01F9}\symactive{"1E3E}\symactive{"1E3F}
\symhyph{"2011}\symhyph{"2E1A}\symblankrange{"1F300}{"1FAFF}\symblank{"FE0F}
% Caractères chinois : WenQuanYi Zen Hei (sous-ensemble GB2312 + pinyin), gras/italique simulés
\setCJKmainfont{WenQuanYiZenHei-subset.ttf}[Path=../../../fonts/,AutoFakeBold=3,AutoFakeSlant=0.2]
\xeCJKsetup{CJKspace=false}
% Pinyin : voyelles à caron (ǎ ǐ ǒ ǔ ǖ ǘ ǚ ǜ) absentes de la police principale -> caractères actifs
\newfontface\pyfont{WenQuanYiZenHei-subset.ttf}[Path=../../../fonts/]
\newcommand\pyactive[1]{\catcode#1=13 \begingroup\lccode`\~=#1 \lowercase{\endgroup\def~}{{\pyfont\char#1}}}
\pyactive{"01CD}\pyactive{"01CE}\pyactive{"01CF}\pyactive{"01D0}\pyactive{"01D1}\pyactive{"01D2}\pyactive{"01D3}\pyactive{"01D4}\pyactive{"01D5}\pyactive{"01D6}\pyactive{"01D7}\pyactive{"01D8}\pyactive{"01D9}\pyactive{"01DA}\pyactive{"01DB}\pyactive{"01DC}
% Maths en sans-serif (Fira Math) avec chiffres et lettres droites de Plus
% Jakarta, pour que les formules se fondent dans le texte.
\usepackage[math-style=ISO,bold-style=ISO]{unicode-math}
\setmathfont{FiraMath-Regular.otf}[Path=../../../fonts/]
\setmathfont{PlusJakartaSans-Medium.ttf}[Path=../../../fonts/,range={up/num,up/{Latin,latin},\mathup}]
\usepackage{icomma} % virgule décimale sans espace en mode math
% Caractères mathématiques Unicode absents des polices de texte (Plus Jakarta,
% Archivo Black) : rendus par leur équivalent mathématique, en texte comme en
% formule — sinon ils s'affichent en carré vide (ex. « Arithmétique dans ℤ »).
\usepackage{newunicodechar}
\newunicodechar{ℝ}{\ensuremath{\mathbb{R}}}
\newunicodechar{ℤ}{\ensuremath{\mathbb{Z}}}
\newunicodechar{ℂ}{\ensuremath{\mathbb{C}}}
\newunicodechar{ℕ}{\ensuremath{\mathbb{N}}}
\newunicodechar{ℚ}{\ensuremath{\mathbb{Q}}}
\newunicodechar{𝔻}{\ensuremath{\mathbb{D}}}
\newunicodechar{α}{\ensuremath{\alpha}}
\newunicodechar{β}{\ensuremath{\beta}}
\newunicodechar{γ}{\ensuremath{\gamma}}
\newunicodechar{δ}{\ensuremath{\delta}}
\newunicodechar{ε}{\ensuremath{\varepsilon}}
\newunicodechar{θ}{\ensuremath{\theta}}
\newunicodechar{λ}{\ensuremath{\lambda}}
\newunicodechar{μ}{\ensuremath{\mu}}
\newunicodechar{σ}{\ensuremath{\sigma}}
\newunicodechar{τ}{\ensuremath{\tau}}
\newunicodechar{φ}{\ensuremath{\varphi}}
\newunicodechar{ψ}{\ensuremath{\psi}}
\newunicodechar{ω}{\ensuremath{\omega}}
\newunicodechar{Σ}{\ensuremath{\Sigma}}
% majuscules produites par \MakeUppercase dans les titres (α-aminés -> Α)
\newunicodechar{Α}{\ensuremath{\alpha}}
\newunicodechar{Β}{\ensuremath{\beta}}
\newunicodechar{Γ}{\ensuremath{\Gamma}}
\newunicodechar{Δ}{\ensuremath{\Delta}}
\newunicodechar{Ε}{\ensuremath{\varepsilon}}
\newunicodechar{Θ}{\ensuremath{\Theta}}
\newunicodechar{Λ}{\ensuremath{\Lambda}}
\newunicodechar{Μ}{\ensuremath{\mu}}
\newunicodechar{Τ}{\ensuremath{\tau}}
\newunicodechar{Φ}{\ensuremath{\Phi}}
\newunicodechar{Ψ}{\ensuremath{\Psi}}
\newunicodechar{Ω}{\ensuremath{\Omega}}
\newunicodechar{↦}{\ensuremath{\mapsto}}
\newunicodechar{∈}{\ensuremath{\in}}
\newunicodechar{∉}{\ensuremath{\notin}}
\newunicodechar{∧}{\ensuremath{\wedge}}
\newunicodechar{⇔}{\ensuremath{\Leftrightarrow}}
\newunicodechar{⟺}{\ensuremath{\Longleftrightarrow}}
\newunicodechar{⇒}{\ensuremath{\Rightarrow}}
\newunicodechar{⟹}{\ensuremath{\Longrightarrow}}
\newunicodechar{∘}{\ensuremath{\circ}}
\newunicodechar{∪}{\ensuremath{\cup}}
\newunicodechar{∩}{\ensuremath{\cap}}
\newunicodechar{≡}{\ensuremath{\equiv}}
\newunicodechar{⊂}{\ensuremath{\subset}}
\newunicodechar{⊥}{\ensuremath{\perp}}
\newunicodechar{∃}{\ensuremath{\exists}}
\newunicodechar{∀}{\ensuremath{\forall}}
\newunicodechar{∛}{\ensuremath{\sqrt[3]{\,}}}
\newunicodechar{∜}{\ensuremath{\sqrt[4]{\,}}}
\newunicodechar{⃗}{\ensuremath{{}^{\to}}}
\newunicodechar{ⁿ}{\ifmmode^{n}\else\textsuperscript{\ensuremath{n}}\fi}
\newunicodechar{ˣ}{\ifmmode^{x}\else\textsuperscript{\ensuremath{x}}\fi}
\newunicodechar{ᵇ}{\ifmmode^{b}\else\textsuperscript{\ensuremath{b}}\fi}
\newunicodechar{ʳ}{\ifmmode^{r}\else\textsuperscript{\ensuremath{r}}\fi}
\newunicodechar{ᵘ}{\ifmmode^{u}\else\textsuperscript{\ensuremath{u}}\fi}
\newunicodechar{ʸ}{\ifmmode^{y}\else\textsuperscript{\ensuremath{y}}\fi}
\newunicodechar{ᵃ}{\ifmmode^{a}\else\textsuperscript{\ensuremath{a}}\fi}
\newunicodechar{ᵉ}{\ifmmode^{e}\else\textsuperscript{\ensuremath{e}}\fi}
\newunicodechar{ᵏ}{\ifmmode^{k}\else\textsuperscript{\ensuremath{k}}\fi}
\newunicodechar{ᵖ}{\ifmmode^{p}\else\textsuperscript{\ensuremath{p}}\fi}
\newunicodechar{ᴺ}{\ifmmode^{N}\else\textsuperscript{\ensuremath{N}}\fi}
\newunicodechar{ᶜ}{\ifmmode^{c}\else\textsuperscript{\ensuremath{c}}\fi}
\newunicodechar{ʰ}{\ifmmode^{h}\else\textsuperscript{\ensuremath{h}}\fi}
\newunicodechar{ᵠ}{\ifmmode^{q}\else\textsuperscript{\ensuremath{q}}\fi}
\newunicodechar{ₙ}{\ifmmode_{n}\else\textsubscript{\ensuremath{n}}\fi}
\newunicodechar{ₐ}{\ifmmode_{a}\else\textsubscript{\ensuremath{a}}\fi}
\newunicodechar{ᵢ}{\ifmmode_{i}\else\textsubscript{\ensuremath{i}}\fi}
\newunicodechar{ᵣ}{\ifmmode_{r}\else\textsubscript{\ensuremath{r}}\fi}
\newunicodechar{ₖ}{\ifmmode_{k}\else\textsubscript{\ensuremath{k}}\fi}
\newunicodechar{ᵦ}{\ifmmode_{\beta}\else\textsubscript{\ensuremath{\beta}}\fi}
\newunicodechar{ₓ}{\ifmmode_{x}\else\textsubscript{\ensuremath{x}}\fi}
\newfontfamily\popdisplay{ArchivoBlack-Regular.ttf}[Path=../../../fonts/]
\newfontfamily\poplogo{SpaceGrotesk-Bold.ttf}[Path=../../../fonts/]
\newfontfamily\popsemi{PlusJakartaSans-SemiBold.ttf}[Path=../../../fonts/]
\newfontfamily\popxbold{PlusJakartaSans-ExtraBold.ttf}[Path=../../../fonts/]
\newfontfamily\popxboldls{PlusJakartaSans-ExtraBold.ttf}[Path=../../../fonts/,LetterSpace=3.0]

\setlist[enumerate]{leftmargin=1.7em,itemsep=5pt,topsep=4pt}
\setlist[itemize]{leftmargin=1.7em,itemsep=4pt,topsep=4pt,label={\color{poporange}\textbullet}}
\setlength{\parindent}{0pt}
\setlength{\parskip}{5pt}
\renewcommand{\arraystretch}{1.25}

% pgfplots : style Pop par défaut
\pgfplotsset{
  every axis/.append style={axis line style={popink,line width=1pt},tick style={popink},
    grid style={popline,line width=0.5pt},label style={font=\small},tick label style={font=\footnotesize}},
  popcurve/.style={poporange,line width=1.8pt,no markers,smooth},
}
% Même style disponible dans un \draw TikZ simple (hors pgfplots)
\tikzset{popcurve/.style={poporange,line width=1.8pt}}
\tikzset{popgrid/.style={popline,very thin}}
\colorlet{popcurve}{poporange} % le nom du style est parfois utilisé comme couleur

% ============================================================
% Pill / squiggle
% ============================================================
\newcommand{\poppill}[3]{%
  \tikz[baseline=-0.55ex]\node[draw=popink,line width=1.2pt,rounded corners=2.6pt,%
    fill=#2,inner xsep=6.5pt,inner ysep=3pt]{%
    \popxboldls\fontsize{7}{7}\selectfont\color{#3}\MakeUppercase{#1}};%
}
\newcommand{\popsquiggle}[1]{%
  \tikz[baseline=-0.4ex]{
    \draw[#1,line width=2.6pt,line cap=round]
      plot [smooth] coordinates {(0,0.09) (0.34,0.01) (0.68,0.21) (1.02,0.01) (1.36,0.10)};
  }%
}
% Mot-clé surligné (marqueur orange sous le texte)
\newcommand{\cle}[1]{{\popxbold\color{popdark}#1}}

% ============================================================
% En-tête / pied
% ============================================================
\pagestyle{fancy}
\fancyhf{}
\renewcommand{\headrulewidth}{0pt}
\renewcommand{\footrulewidth}{0pt}
\lhead{\raisebox{-0.15ex}{\poplogo\fontsize{13}{13}\selectfont\color{popink}OnBuch\color{poporange}+}}
\rhead{\poppill{\DOCMATIERE\ \textbullet\ \DOCNIVEAU\ \textbullet\ Leçon \DOCLECON}{white}{popink}}
\lfoot{\popxbold\fontsize{7.6}{7.6}\selectfont\color{popmuted}\MakeUppercase{Cours signé OnBuch+}\ \textbullet\ {\popsemi\DOCTITRE}}
\rfoot{\tikz[baseline=-0.6ex]\node[rounded corners=8.5pt,fill=popink,minimum height=17pt,minimum width=17pt,inner xsep=5pt,inner sep=0]{\popxbold\fontsize{8}{8}\selectfont\color{popcream}\thepage\,/\,\pageref*{LastPage}};}

% ============================================================
% Titres
% ============================================================
\newcommand{\chaptitre}[2]{%
  \begin{tcolorbox}[enhanced,colback=poporange,colframe=popink,boxrule=2.6pt,arc=11pt,
      drop shadow=popink,left=15pt,right=15pt,top=12pt,bottom=11pt,before skip=3pt,after skip=18pt]
    \poppill{#2}{popink}{popcream}\par\vspace{9pt}
    {\raggedright\hyphenpenalty=10000\exhyphenpenalty=10000\popdisplay\fontsize{22}{24}\selectfont\color{popcream}\MakeUppercase{#1}\par}
  \end{tcolorbox}
}
% Section numérotée : pastille orange avec le numéro + titre Archivo
\newcounter{popsec}
\newcommand{\coursec}[1]{%
  \stepcounter{popsec}\setcounter{popsub}{0}%
  \par\vspace{16pt}\noindent
  \begin{minipage}{\linewidth}
  \tikz[baseline=-0.7ex]\node[draw=popink,line width=1.6pt,fill=poporange,rounded corners=4pt,minimum size=22pt,inner sep=0]{\popdisplay\fontsize{12}{12}\selectfont\color{popcream}\thepopsec};%
  \hspace{8pt}{\popdisplay\fontsize{15}{17}\selectfont\color{popink}\MakeUppercase{#1}}\par
  \vspace{4pt}\noindent{\color{poporange}\rule{36pt}{3.6pt}}
  \end{minipage}\par\nopagebreak\vspace{8pt}
}
\newcounter{popsub}
\newcommand{\courssub}[1]{%
  \stepcounter{popsub}%
  \par\vspace{9pt}\noindent{\popxbold\fontsize{11.5}{13}\selectfont\color{popink}{\color{poporange}\thepopsec.\thepopsub}\hspace{6pt}#1}\par\nopagebreak\vspace{4pt}
}

% ============================================================
% Boîtes pédagogiques — même recette : fond blanc, bordure épaisse colorée,
% ombre dure encre, badge plein. Titre optionnel [#1].
% ============================================================
\tcbset{popbase/.style={breakable,enhanced,colback=white,boxrule=2.3pt,arc=9pt,
  shadow={3pt}{-3pt}{0pt}{popink},left=14pt,right=14pt,top=11pt,bottom=11pt,before skip=11pt,after skip=15pt,
  fontupper=\popsemi\fontsize{10.6}{14.5}\selectfont\color{popink},
  fontlower=\popsemi\fontsize{10.6}{14.5}\selectfont\color{popink}}}
% Coche dessinée (le glyphe ✓ n'existe pas dans les polices du document)
\newcommand{\popcheck}[1]{\tikz[baseline=-0.5ex]\draw[#1,line width=1.6pt,line cap=round,line join=round](-0.13,0.0)--(-0.04,-0.09)--(0.13,0.1);}
\providecommand{\coche}{\popcheck{popgreen}}
\providecommand{\resultat}[1]{\par\smallskip\noindent\fcolorbox{popgreen}{popgreenL}{\parbox{\dimexpr\linewidth-2\fboxsep-2\fboxrule\relax}{\textbf{Résultat.} #1}}\par\smallskip}
\newcommand{\popboxhead}[3]{\poppill{#1}{#2}{white}\ifx\relax#3\relax\else\hspace{6pt}{\popxbold\fontsize{10.6}{12}\selectfont\color{popink}#3}\fi\par\vspace{7pt}}

\newtcolorbox{aretenir}[1][]{popbase,colframe=popgreen,before upper={\popboxhead{À retenir}{popgreen}{#1}}}
% Texte en chinois (document de lecture, dialogue, extrait) : cadre
% d'étude. Un titre optionnel (source, auteur) s'affiche dans l'en-tête.
\newtcolorbox{textebox}[1][]{popbase,colframe=popblue,before upper={\popboxhead{文本 · Texte}{popblue}{#1}},after upper={}}
% Marques ✓ / ✗ (forme correcte / fautive) souvent utilisées par les modèles
\providecommand{\cross}{\textcolor{poppink}{\ensuremath{\times}}}
\providecommand{\xmark}{\cross}\providecommand{\crossmark}{\cross}\providecommand{\cmark}{\ensuremath{\checkmark}}
% Ligne de réponse / justification à compléter (inventée par les modèles)
\providecommand{\justification}[1]{\par\noindent\textit{Justification~:} \dotfill\par}
% \textipa (package tipa non chargé) : la police ne couvre pas l'API -> simple passage
\providecommand{\textipa}[1]{#1}
% Soulignement (ulem non chargé) : \uline, \ul
\providecommand{\uline}[1]{\underline{#1}}\providecommand{\ul}[1]{\underline{#1}}
% \glossaire{\item ...} inventée par les modèles : liste de glossaire
\providecommand{\glossaire}[1]{\begin{itemize}#1\end{itemize}}
% \alert (beamer) inventée par les modèles : texte mis en évidence
\providecommand{\alert}[1]{\textbf{\textcolor{poppink}{#1}}}
% variantes ASCII d'ordinaux écrites en macro
\providecommand{\iere}{\textsuperscript{re}}\providecommand{\ieme}{\textsuperscript{e}}\providecommand{\ere}{\textsuperscript{re}}\providecommand{\eme}{\textsuperscript{e}}\providecommand{\er}{\textsuperscript{er}}\providecommand{\ie}{\textsuperscript{e}}
% Ordinaux français écrits en macro par les modèles : 1\ière, 3\ième
\newcommand{\ière}{\textsuperscript{re}}\newcommand{\ième}{\textsuperscript{e}}
% \exemple (inventée par les modèles) : étiquette « Exemple : »
\providecommand{\exemple}{\textbf{Exemple~:}\ }
% \square utilisable aussi hors mode math (cases à cocher)
\AtBeginDocument{\let\squareMath\square\renewcommand{\square}{\ensuremath{\squareMath}}}
% Mot ou phrase en chinois à l'intérieur d'un paragraphe français
\newcommand{\zh}[1]{\ifmmode\text{#1}\else{#1}\fi}
\let\eng\zh \let\esp\zh \let\all\zh % alias tolérés
% flèches d'intonation inventées par les modèles
\providecommand{\textsearrow}{}\renewcommand{\textsearrow}{\ensuremath{\searrow}}\providecommand{\textnearrow}{}\renewcommand{\textnearrow}{\ensuremath{\nearrow}}
% \fr{..} : mot français (inventé par les modèles)
\providecommand{\fr}[1]{\textit{#1}}
% zhbox : encadré de texte chinois inventé par les modèles
\newenvironment{zhbox}{\begin{quote}}{\end{quote}}
% \courssubsub : titre de niveau 3 inventé par les modèles
\providecommand{\courssubsub}[1]{\par\medskip\noindent\textbf{#1}\par\smallskip}
% \vs : « versus » inventé par les modèles
\providecommand{\vs}{\textit{vs}\ }
% \blank : case à compléter inventée par les modèles
\providecommand{\blank}[1][]{\underline{\hspace{1.6em}}}
\newtcolorbox{exemplebox}[1][]{popbase,colframe=poporange,before upper={\popboxhead{Exemple}{poporange}{#1}}}
\newtcolorbox{definition}[1][]{popbase,colframe=popblue,before upper={\popboxhead{Définition}{popblue}{#1}}}
\newtcolorbox{propriete}[1][]{popbase,colframe=poppurple,before upper={\popboxhead{Propriété}{poppurple}{#1}}}
\newtcolorbox{methode}[1][]{popbase,colframe=popgold,before upper={\popboxhead{Méthode}{popgold}{#1}}}
\newtcolorbox{attention}[1][]{popbase,colframe=poppink,before upper={\popboxhead{Attention}{poppink}{#1}}}
\newtcolorbox{experience}[1][]{popbase,colframe=popblue,colback=popblueL,before upper={\popboxhead{Expérience}{popblue}{#1}}}
% « Le savais-tu ? » : carte violette pleine (contexte, culture, Cameroun)
\newtcolorbox{savaistu}[1][]{popbase,colback=poppurple,colframe=popink,
  fontupper=\popsemi\fontsize{10.4}{14.2}\selectfont\color{white},
  before upper={\poppill{Le savais-tu\,?}{white}{poppurple}\ifx\relax#1\relax\else\hspace{6pt}{\popxbold\color{white}#1}\fi\par\vspace{7pt}}}
% Exercice résolu : énoncé en haut, \tcblower puis la solution sur fond crème
\newcounter{popexo}\newcounter{popexr}
\newtcolorbox{exoresolu}[1][]{popbase,colframe=poporange,colbacklower=popcream,
  segmentation style={popink,line width=1.2pt,dashed},
  before upper={\stepcounter{popexr}\popboxhead{Exercice résolu \thepopexr}{poporange}{#1}},
  before lower={\poppill{Solution}{popink}{popcream}\par\vspace{6pt}}}
% Exercice (non résolu en place) — la correction est dans la partie Corrigés
\newtcolorbox{exercice}[1][]{popbase,colframe=popink,
  before upper={\stepcounter{popexo}\popboxhead{Exercice \thepopexo}{popink}{#1}}}
% Corrigé
\newtcolorbox{corrige}[1][]{popbase,colframe=popgreen,colback=popgreenL,
  before upper={\popboxhead{Corrigé}{popgreen}{#1}}}
% Objectifs / prérequis / bilan
\newtcolorbox{objectifs}{popbase,colframe=popink,colback=poporangeL,
  before upper={\popboxhead{Objectifs de la leçon}{popink}{}}}
\newtcolorbox{prerequis}{popbase,colframe=popink,colback=popblueL,
  before upper={\popboxhead{Prérequis}{popblue}{}}}
\newtcolorbox{bilan}{popbase,colframe=popink,colback=white,boxrule=2.8pt,
  before upper={\popboxhead{Fiche bilan}{poporange}{L'essentiel à connaître pour l'examen}}}
% Figure encadrée avec légende
\newtcolorbox{popfigure}[1][]{enhanced,breakable=false,colback=white,colframe=popline,boxrule=1.4pt,arc=8pt,
  left=10pt,right=10pt,top=10pt,bottom=8pt,before skip=10pt,after skip=12pt,halign=center,
  lower separated=false,halign lower=center,
  fontlower=\popsemi\fontsize{9}{11.5}\selectfont\color{popmuted},
  #1}
\newcommand{\legende}[1]{\par\vspace{6pt}{\popsemi\fontsize{9}{11.5}\selectfont\color{popmuted}#1}}

% ============================================================
% Signature OnBuch+
% ============================================================
\newcommand{\poplogoinline}[1]{{\poplogo\fontsize{#1}{#1}\selectfont\color{popink}OnBuch\color{poporange}+}}
\newcommand{\signatureonbuch}{%
  \begin{tcolorbox}[enhanced,colback=popink,colframe=popink,boxrule=2.2pt,arc=10pt,
      drop shadow=poporange,left=14pt,right=14pt,top=10pt,bottom=10pt,before skip=0pt,after skip=0pt]
    \tikz[baseline=-0.7ex]\node[circle,fill=popgreen,draw=popcream,line width=1.3pt,minimum size=22pt,inner sep=0]{\popcheck{white}};%
    \hspace{9pt}\begin{minipage}[c]{0.62\linewidth}
      {\popxbold\fontsize{12}{13}\selectfont\color{popcream}Signé\ \textbullet\ {\poplogo OnBuch\color{poporange}+}}\par\vspace{2pt}
      {\raggedright\popsemi\fontsize{8}{10.5}\selectfont\color{popline}\MakeUppercase{Équipe pédagogique OnBuch+}\\\MakeUppercase{Conforme au programme officiel MINESEC}\par}
    \end{minipage}\hfill
    {\poplogo\fontsize{22}{22}\selectfont\color{popcream}OnBuch\color{poporange}+}
  \end{tcolorbox}%
}

% ============================================================
% Page de garde
% #1 titre · #2 matière · #3 niveau · #4 module · #5 n° leçon · #6 durée ·
% #7 description / objectifs (texte ou itemize)
% ============================================================
\newcommand{\pagegarde}[7]{%
  \thispagestyle{empty}
  \newgeometry{a4paper,margin=1.7cm,top=1.6cm,bottom=1.4cm}%
  \begin{tikzpicture}[remember picture,overlay]
    \fill[poporange!18] ([xshift=-3.2cm,yshift=-3.6cm]current page.north east) circle (4.6cm);
    \fill[poppurple!14] ([xshift=2.4cm,yshift=7.6cm]current page.south west) circle (3.8cm);
  \end{tikzpicture}%
  {\poplogo\fontsize{30}{30}\selectfont\color{popink}OnBuch\color{poporange}+}\hfill
  \raisebox{0.6ex}{\poppill{Cours signé \textbullet\ Premium}{popink}{popcream}}\par
  \vspace{1pt}\popsquiggle{poppurple}\par
  \vspace{8pt}
  \begin{tcolorbox}[enhanced,colback=poporange,colframe=popink,boxrule=2.8pt,arc=13pt,
      drop shadow=popink,left=18pt,right=18pt,top=12pt,bottom=13pt,after skip=10pt]
    \poppill{#2 \textbullet\ #3}{popink}{popcream}\hspace{5pt}\poppill{#4}{white}{popink}\par\vspace{8pt}
    {\popdisplay\fontsize{14}{14}\selectfont\color{popink}LEÇON #5}\par\vspace{4pt}
    {\raggedright\hyphenpenalty=10000\exhyphenpenalty=10000\popdisplay\fontsize{25}{27}\selectfont\color{popcream}\MakeUppercase{#1}\par}
    \vspace{8pt}
    {\popxbold\fontsize{9.5}{9.5}\selectfont\color{popink}\MakeUppercase{Durée conseillée : #6}}
  \end{tcolorbox}
  \begin{tcolorbox}[enhanced,colback=white,colframe=popink,boxrule=2.2pt,arc=10pt,
      drop shadow=popink,left=16pt,right=16pt,top=10pt,bottom=10pt,after skip=8pt]
    \poppill{Dans cette leçon}{poppurple}{white}\par\vspace{5pt}
    \popsemi\fontsize{9.3}{12.6}\selectfont\color{popink}#7
  \end{tcolorbox}
  \noindent\begin{minipage}[t]{0.487\linewidth}
    \begin{tcolorbox}[enhanced,colback=popgreen,colframe=popink,boxrule=2.2pt,arc=10pt,
        drop shadow=popink,left=11pt,right=11pt,top=10pt,bottom=10pt,height=3.1cm,valign=center]
      \tcbox[colback=white,colframe=popink,boxrule=1.3pt,arc=5pt,boxsep=0pt,nobeforeafter,
        left=3pt,right=3pt,top=3pt,bottom=3pt,valign=center]{\qrcode[height=1.9cm]{https://play.google.com/store/apps/details?id=com.luvvix.onbuch&hl=fr}}%
      \hspace{8pt}\begin{minipage}[c]{0.5\linewidth}
        {\popdisplay\fontsize{11}{12.5}\selectfont\color{white}\MakeUppercase{Télécharge OnBuch}}\par\vspace{4pt}
        {\raggedright\popsemi\fontsize{8.4}{11}\selectfont\color{white}Gratuit \textbullet\ Google Play\\Tuteur IA, annales, concours, résultats\par}
      \end{minipage}
    \end{tcolorbox}
  \end{minipage}\hfill
  \begin{minipage}[t]{0.487\linewidth}
    \begin{tcolorbox}[enhanced,colback=poppurple,colframe=popink,boxrule=2.2pt,arc=10pt,
        drop shadow=popink,left=11pt,right=11pt,top=10pt,bottom=10pt,height=3.1cm,valign=center]
      \tikz[baseline=-0.7ex]\node[circle,fill=white,draw=popink,line width=1.3pt,minimum size=1.6cm,inner sep=0]{\popdisplay\fontsize{24}{24}\selectfont\color{poppurple}+};%
      \hspace{8pt}\begin{minipage}[c]{0.6\linewidth}
        {\popdisplay\fontsize{11}{12.5}\selectfont\color{white}\MakeUppercase{Passe à OnBuch+}}\par\vspace{4pt}
        {\raggedright\popsemi\fontsize{8.4}{11}\selectfont\color{white}Tous les cours signés, Tuteur IA illimité, fiches PDF — onglet OnBuch+ de l'app\par}
      \end{minipage}
    \end{tcolorbox}
  \end{minipage}
  \vspace{8pt}
  \popcontenu
  \vfill
  \signatureonbuch
  \restoregeometry
  \newpage
}

% Rangée « Ce cours contient » (page de garde)
\newcommand{\poptile}[3]{% #1 couleur #2 gros texte #3 légende
  \begin{tcolorbox}[enhanced,colback=#1,colframe=popink,boxrule=2pt,arc=9pt,shadow={2.5pt}{-2.5pt}{0pt}{popink},
      left=6pt,right=6pt,top=9pt,bottom=9pt,halign=center,valign=center,height=1.75cm,width=\linewidth,nobeforeafter]
    {\popdisplay\fontsize{12.5}{13}\selectfont\color{white}\MakeUppercase{#2}}\par\vspace{3pt}
    {\centering\popsemi\fontsize{7.8}{9.5}\selectfont\color{white}#3\par}
  \end{tcolorbox}}
\newcommand{\popcontenu}{%
  \par\noindent{\popxboldls\fontsize{7.5}{7.5}\selectfont\color{popmuted}CE COURS SIGNÉ CONTIENT}\par\vspace{7pt}
  \noindent\begin{minipage}[t]{0.235\linewidth}\poptile{popblue}{Cours}{complet, illustré, pas à pas}\end{minipage}\hfill
  \begin{minipage}[t]{0.235\linewidth}\poptile{poporange}{Méthodes}{et exercices résolus}\end{minipage}\hfill
  \begin{minipage}[t]{0.235\linewidth}\poptile{poppink}{Exercices}{type examen + corrigés}\end{minipage}\hfill
  \begin{minipage}[t]{0.235\linewidth}\poptile{popgreen}{Fiche bilan}{l'essentiel à retenir}\end{minipage}\par}

% Sommaire
\newtcolorbox{sommaire}{enhanced,colback=white,colframe=popink,boxrule=2.2pt,arc=10pt,
  drop shadow=popink,left=18pt,right=18pt,top=13pt,bottom=13pt,before skip=6pt,after skip=20pt,
  before upper={\poppill{Sommaire}{poppurple}{white}\par\vspace{9pt}\popsemi\fontsize{10.6}{14.5}\selectfont\color{popink}}}

% Signature de fin de cours — poussée tout en bas de la dernière page
\newcommand{\findecours}{%
  \par\vfill\nopagebreak
  \begin{center}\popsquiggle{poporange}\end{center}\vspace{4pt}
  \signatureonbuch
}
\AtBeginDocument{\def\llbracket{\mathopen{[\mkern-3.5mu[}}\def\rrbracket{\mathclose{]\mkern-3.5mu]}}} % intervalles d'entiers (glyphe ⟦ absent de la police)
% Glyphes absents de Fira Math : redéfinis à partir de glyphes disponibles
\AtBeginDocument{%
  \def\setminus{\mathbin{\backslash}}\let\smallsetminus\setminus
  \def\vdots{\mathord{\vbox{\baselineskip3.2pt\lineskiplimit0pt\kern4pt\hbox{.}\hbox{.}\hbox{.}}}}%
  \def\ddots{\mathinner{\mkern1mu\raise6.5pt\hbox{.}\mkern2mu\raise3.6pt\hbox{.}\mkern2mu\raise0.7pt\hbox{.}\mkern1mu}}%
  \def\triangle{\mathord{\scalebox{0.9}{$\symup{\Delta}$}}}%
  \def\nabla{\mathord{\raisebox{\depth}{\scalebox{1}[-1]{$\symup{\Delta}$}}}}%
  \def\checkmark{\popcheck{popdark}}%
  \def\star{\mathbin{\ast}}\def\bigstar{\mathord{\ast}}%
  \def\diamond{\mathbin{\scalebox{0.6}{$\Diamond$}}}%
  \def\bigcup{\mathop{\mathchoice{\raisebox{-0.3em}{\scalebox{1.7}{$\cup$}}}{\scalebox{1.25}{$\cup$}}{\cup}{\cup}}\displaylimits}%
  \def\bigcap{\mathop{\mathchoice{\raisebox{-0.3em}{\scalebox{1.7}{$\cap$}}}{\scalebox{1.25}{$\cap$}}{\cap}{\cap}}\displaylimits}%
  \def\lesseqqgtr{\mathrel{\lessgtr}}\def\bigoplus{\mathop{\oplus}\displaylimits}%
}
\providecommand{\ssi}{si et seulement si} % abréviation fréquente
\AtBeginDocument{\providecommand{\wideparen}{\overparen}} % arc de cercle

%% FIGFIX-BEGIN v5 — correctifs globaux des figures (docs/onbuch-generation/tools/figfix)
\usepackage{adjustbox}\usepackage{etoolbox}\tcbuselibrary{raster}
% toute figure TikZ/pgfplots plus large que la ligne est réduite à la largeur disponible
\BeforeBeginEnvironment{tikzpicture}{\begin{adjustbox}{max width=\linewidth}}
\AfterEndEnvironment{tikzpicture}{\end{adjustbox}}
% légendes pgfplots lisibles par-dessus les courbes
\pgfplotsset{every axis/.append style={legend style={fill=white,fill opacity=0.92,text opacity=1,draw=black!15,inner sep=3pt}}}
% étiquettes d'arêtes (syntaxe edge["..."]) avec fond blanc
\tikzset{every edge quotes/.append style={fill=white,inner sep=1.2pt}}
%% FIGFIX-END
\begin{document}

\pagegarde{Leçon 20 — Grammaire : 比较字句:……跟/和…. (不)一样et ……跟/和…. (不)一样+ 形容词 ; 一方面… ,另一方面… ; 如果…, 但是/可是…}{Chinois}{1ère A}{Module 3 : Citoyenneté et ouverture au monde}{ZHA20}{2 h}{Cette leçon de grammaire présente trois structures essentielles du chinois niveau 1ère : la phrase comparative d'égalité avec 跟/和……(不)一样 (avec ou sans adjectif), le double connecteur d'argumentation 一方面……另一方面, et la construction hypothético-adversative 如果……但是/可是. Chaque structure est observée, schématisée, comparée au français et entraînée par des exercices d'application et de transformation.
\vspace{4pt}
\begin{multicols}{2}\small
\begin{itemize}
\item À la fin de cette leçon, je sais construire une comparaison d'égalité affirmative avec A 跟/和 B 一样
\item À la fin de cette leçon, je sais construire la négation avec A 跟/和 B 不一样 et placer correctement 不
\item À la fin de cette leçon, je sais ajouter un adjectif après 一样 pour préciser le point de comparaison (一样高, 一样大)
\item À la fin de cette leçon, je sais distinguer 跟 et 和 dans cette structure et connaître leur interchangeabilité
\item À la fin de cette leçon, je sais présenter deux aspects d'une situation avec 一方面……,另一方面……
\item À la fin de cette leçon, je sais exprimer une hypothèse nuancée avec 如果……,但是/可是……
\end{itemize}\end{multicols}}

\begin{sommaire}
\begin{multicols}{2}
\begin{enumerate}
\item Observation et découverte des exemples
\item Structure 1 : A 跟/和 B (不)一样
\item Structure 2 : 一方面……,另一方面……
\item Structure 3 : 如果……,但是/可是……
\item Comparaison avec le français et pièges des francophones
\item Exercices d'application et de transformation
\item Activité d'intégration
\item Méthodes et exercices résolus
\item Exercices
\item Corrigés des exercices
\item Fiche bilan
\end{enumerate}
\end{multicols}
\end{sommaire}

\begin{objectifs}
\begin{itemize}
\item À la fin de cette leçon, je sais construire une comparaison d'égalité affirmative avec A 跟/和 B 一样
\item À la fin de cette leçon, je sais construire la négation avec A 跟/和 B 不一样 et placer correctement 不
\item À la fin de cette leçon, je sais ajouter un adjectif après 一样 pour préciser le point de comparaison (一样高, 一样大)
\item À la fin de cette leçon, je sais distinguer 跟 et 和 dans cette structure et connaître leur interchangeabilité
\item À la fin de cette leçon, je sais présenter deux aspects d'une situation avec 一方面……,另一方面……
\item À la fin de cette leçon, je sais exprimer une hypothèse nuancée avec 如果……,但是/可是……
\item À la fin de cette leçon, je sais éviter les erreurs typiques des francophones (ordre des mots, place de 不, oubli de 一样)
\item À la fin de cette leçon, je sais transformer des phrases simples en phrases comparatives et argumentatives
\item À la fin de cette leçon, je sais employer ces trois structures dans un court texte cohérent
\end{itemize}
\end{objectifs}

\begin{prerequis}
\begin{itemize}
\item Les pronoms personnels et la structure de base Sujet-Verbe-Complément (Seconde)
\item Les adjectifs qualificatifs courants : 大, 小, 高, 好, 多, 贵, 漂亮 (Seconde)
\item La négation avec 不 et 没 (Seconde)
\item La comparaison de supériorité avec 比 (leçon antérieure de 1ère)
\item Les connecteurs simples 因为……所以 et 虽然……但是 (début de 1ère)
\end{itemize}
\end{prerequis}

\newpage

\coursec{Observation et découverte des exemples}

\courssub{Situation d'accroche : au marché Mokolo}

Imaginez la scène. Au marché Mokolo de Yaoundé, un samedi matin, deux élèves de Première comparent leurs téléphones portables. L'un s'exclame : « Mon téléphone est comme le tien, mais le mien est plus cher ! » Plus tard, assis devant un beignet-haricot, ils parlent de leurs études. L'un dit : « D'un côté, je veux apprendre le chinois ; de l'autre, je dois réviser l'anglais. » Son ami lui répond : « Si tu es intelligent mais que tu ne travailles pas, tu ne réussiras pas. »

En trois phrases de la vie quotidienne, ces deux élèves viennent d'utiliser trois outils grammaticaux essentiels de toute langue : la \cle{comparaison d'égalité} (« comme le tien », « aussi... que »), l'\cle{argumentation en deux points} (« d'un côté... de l'autre ») et l'\cle{hypothèse combinée à une opposition} (« si tu es intelligent... mais... »).

\textbf{Problématique :} comment exprimer en chinois la comparaison d'égalité \zh{跟/和……一样}, l'argumentation en deux points \zh{一方面……另一方面} et l'hypothèse avec opposition \zh{如果……但是/可是} ? C'est l'objet de cette leçon. Avant d'énoncer les règles, nous allons d'abord \emph{observer} des exemples authentiques, comme le font les vrais sinologues.

\courssub{Lecture guidée : six exemples à observer}

Lisons attentivement les phrases suivantes. Pour chacune, le chinois est donné en caractères, puis en pinyin, puis en traduction française. Ne cherchez pas encore la règle : contentez-vous de lire, d'écouter le rythme de la phrase et de repérer ce qui se répète.

\begin{textebox}[Trois phrases de comparaison]
我的学校跟你的学校一样。\\
\emph{Wǒ de xuéxiào gēn nǐ de xuéxiào yíyàng.}\\
« Mon école est pareille à la tienne. »\\[0.4em]
他跟我一样高。\\
\emph{Tā gēn wǒ yíyàng gāo.}\\
« Il est aussi grand que moi. »\\[0.4em]
今天和昨天不一样热。\\
\emph{Jīntiān hé zuótiān bù yíyàng rè.}\\
« Aujourd'hui, il ne fait pas aussi chaud qu'hier. »
\end{textebox}

\begin{exemplebox}[Trois exemples supplémentaires]
\textbf{Exemple 4 — argumentation :}\\
\zh{一方面我想学汉语，另一方面我要准备考试。}\\
\emph{Yì fāngmiàn wǒ xiǎng xué Hànyǔ, lìng yì fāngmiàn wǒ yào zhǔnbèi kǎoshì.}\\
« D'un côté je veux apprendre le chinois, de l'autre je dois préparer l'examen. »\\[0.5em]
\textbf{Exemple 5 — hypothèse et opposition :}\\
\zh{如果你很聪明，但是不努力，就不能成功。}\\
\emph{Rúguǒ nǐ hěn cōngming, dànshì bù nǔlì, jiù bù néng chénggōng.}\\
« Si tu es intelligent mais que tu ne travailles pas, tu ne peux pas réussir. »\\[0.5em]
\textbf{Exemple 6 — comparaison avec 和 :}\\
\zh{她和我一样漂亮。}\\
\emph{Tā hé wǒ yíyàng piàoliang.}\\
« Elle est aussi belle que moi. »
\end{exemplebox}

\courssub{Repérage : soulignons les mots-outils}

\begin{methode}[Observer avant de mémoriser : la démarche en quatre étapes]
Face à une nouvelle structure chinoise, le bon élève ne récite pas d'abord la règle : il \emph{observe}. Voici la démarche à suivre pour cette leçon et pour toutes les leçons de grammaire de l'année.
\begin{enumerate}
\item \textbf{Lire} les exemples à voix haute, en respectant les tons, jusqu'à ce que la phrase sonne naturelle.
\item \textbf{Repérer les mots-outils} : souligner les petits mots qui reviennent dans chaque exemple (\zh{跟}, \zh{和}, \zh{一样}, \zh{一方面}, \zh{另一方面}, \zh{如果}, \zh{但是}, \zh{可是}).
\item \textbf{Formuler une hypothèse} : deviner la place de chaque élément et le sens de la structure, en comparant avec le français.
\item \textbf{Vérifier avec la règle} : confronter son hypothèse à l'énoncé grammatical du professeur, puis la corriger si nécessaire.
\end{enumerate}
\end{methode}

Appliquons l'étape 2. Dans les six exemples, soulignez dans votre cahier les mots-outils suivants :

\begin{center}
\begin{tabularx}{\linewidth}{|X|X|X|X|}
\hline
\rowcolor{popblueL} Caractères & Pinyin & Nature & Traduction \\
\hline
跟 & gēn & préposition & avec, comme (comparaison) \\
\hline
和 & hé & préposition / conjonction & avec, et \\
\hline
一样 & yíyàng & adjectif & pareil, identique \\
\hline
不一样 & bù yíyàng & adjectif nié & pas pareil, différent \\
\hline
一方面 & yì fāngmiàn & connecteur & d'un côté, d'une part \\
\hline
另一方面 & lìng yì fāngmiàn & connecteur & de l'autre côté, d'autre part \\
\hline
如果 & rúguǒ & conjonction & si (hypothèse) \\
\hline
但是 & dànshì & conjonction & mais \\
\hline
可是 & kěshì & conjonction & mais (plus oral) \\
\hline
就 & jiù & adverbe & alors, donc (conséquence) \\
\hline
\end{tabularx}
\end{center}

\begin{attention}[跟 et 和 : deux mots, un même rôle dans la comparaison]
Dans la structure de comparaison, \zh{跟} (gēn) et \zh{和} (hé) sont \textbf{interchangeables} : \zh{我跟你一样} et \zh{我和你一样} signifient tous les deux « je suis comme toi ». En revanche, attention : \zh{和} peut aussi être la conjonction « et » (\zh{我和你} = « toi et moi »). C'est le mot \zh{一样} placé plus loin dans la phrase qui indique qu'il s'agit d'une \emph{comparaison} et non d'une simple énumération. Cherchez toujours le couple « \zh{跟/和} ... \zh{一样} ».
\end{attention}

\courssub{Hypothèses des élèves : où se place chaque élément ?}

Avant la formalisation, essayons de deviner la structure à partir des exemples observés. Voici les hypothèses qu'une classe de Première formule généralement — et que nous vérifierons dans les sections suivantes.

\begin{itemize}
\item \textbf{Hypothèse 1 (comparaison) :} dans \zh{他跟我一样高}, l'ordre semble être : A + \zh{跟/和} + B + \zh{一样} (+ adjectif). Le français dit « aussi grand \emph{que moi} », le chinois semble dire « avec moi pareil grand » : la personne comparée (B) arrive \emph{avant} l'adjectif.
\item \textbf{Hypothèse 2 (négation) :} dans \zh{今天和昨天不一样热}, la négation \zh{不} se place juste devant \zh{一样}, et l'adjectif \zh{热} reste après.
\item \textbf{Hypothèse 3 (argumentation) :} \zh{一方面} ouvre le premier point, \zh{另一方面} ouvre le second ; chacun est suivi d'une phrase complète, séparée par une virgule.
\item \textbf{Hypothèse 4 (hypothèse + opposition) :} \zh{如果} introduit la condition, \zh{但是/可是} introduit l'opposition, et \zh{就} annonce la conséquence finale.
\end{itemize}

\begin{attention}[Le piège de la traduction mot à mot]
L'erreur la plus fréquente des francophones est de traduire « aussi grand que moi » mot à mot et de produire la phrase fautive *\zh{一样高跟我}. C'est impossible en chinois ! L'ordre chinois est rigide : \textbf{A + \zh{跟/和} + B + \zh{一样} + adjectif}. Le terme de comparaison (B, « moi ») se place \emph{avant} l'adjectif, jamais après. Retenez le schéma : le chinois dit littéralement « lui avec moi pareil grand ».
\end{attention}

\courssub{Carte mentale : les trois structures de la leçon}

\begin{popfigure}
\begin{center}\tcbox[colback=popblue!50,colframe=popblue!50,coltext=white,arc=9pt,boxsep=4pt,left=6pt,right=6pt]{\bfseries\small Leçon 20 Trois structures}\end{center}
\vspace{-2pt}
\begin{tcbraster}[raster columns=2,raster equal height=rows,raster column skip=6pt,raster row skip=6pt,enhanced,blankest]
\begin{tcolorbox}[enhanced,colback=white,colframe=popgreen!45,boxrule=0.9pt,arc=3pt,title={COMPARAISON 跟/和……一样},fonttitle=\bfseries\scriptsize,coltitle=white,colbacktitle=popgreen!45,left=4pt,right=4pt,top=2pt,bottom=2pt,boxsep=1pt,toptitle=1.5pt,bottomtitle=1.5pt,halign=left]
\begin{itemize}[nosep,leftmargin=*,itemsep=1pt,font=\scriptsize]
\item 他跟我一样高。 aussi grand que moi
\item 不一样热 pas aussi chaud
\end{itemize}
\end{tcolorbox}
\begin{tcolorbox}[enhanced,colback=white,colframe=poporange!50,boxrule=0.9pt,arc=3pt,title={ARGUMENTATION 一方面……另一方面},fonttitle=\bfseries\scriptsize,coltitle=white,colbacktitle=poporange!50,left=4pt,right=4pt,top=2pt,bottom=2pt,boxsep=1pt,toptitle=1.5pt,bottomtitle=1.5pt,halign=left]
\begin{itemize}[nosep,leftmargin=*,itemsep=1pt,font=\scriptsize]
\item d'un côté… de l'autre…
\end{itemize}
\end{tcolorbox}
\begin{tcolorbox}[enhanced,colback=white,colframe=poppurple!45,boxrule=0.9pt,arc=3pt,title={HYPOTHÈSE + OPPOSITION 如果……但是/可是},fonttitle=\bfseries\scriptsize,coltitle=white,colbacktitle=poppurple!45,left=4pt,right=4pt,top=2pt,bottom=2pt,boxsep=1pt,toptitle=1.5pt,bottomtitle=1.5pt,halign=left]
\begin{itemize}[nosep,leftmargin=*,itemsep=1pt,font=\scriptsize]
\item si… mais… 就 = alors
\end{itemize}
\end{tcolorbox}
\end{tcbraster}

\legende{Carte mentale des trois structures grammaticales de la leçon 20.}
\end{popfigure}

\courssub{Premier bilan d'observation}

\begin{center}
\begin{tabularx}{\linewidth}{|X|X|X|}
\hline
\rowcolor{popblueL} Structure & Exemple observé (caractères + pinyin) & Traduction \\
\hline
A 跟/和 B 一样 & 我的学校跟你的学校一样。Wǒ de xuéxiào gēn nǐ de xuéxiào yíyàng. & Mon école est pareille à la tienne. \\
\hline
A 跟/和 B 一样 + adjectif & 他跟我一样高。Tā gēn wǒ yíyàng gāo. & Il est aussi grand que moi. \\
\hline
A 跟/和 B 不一样 + adjectif & 今天和昨天不一样热。Jīntiān hé zuótiān bù yíyàng rè. & Il ne fait pas aussi chaud qu'hier. \\
\hline
一方面…，另一方面… & 一方面我想学汉语，另一方面我要准备考试。Yì fāngmiàn wǒ xiǎng xué Hànyǔ, lìng yì fāngmiàn wǒ yào zhǔnbèi kǎoshì. & D'un côté je veux apprendre le chinois, de l'autre je dois préparer l'examen. \\
\hline
如果…，但是…，就… & 如果你很聪明，但是不努力，就不能成功。Rúguǒ nǐ hěn cōngming, dànshì bù nǔlì, jiù bù néng chénggōng. & Si tu es intelligent mais que tu ne travailles pas, tu ne peux pas réussir. \\
\hline
\end{tabularx}
\end{center}

\begin{exoresolu}[Repérage des mots-outils]
Énoncé : dans la phrase \zh{我的汉语跟你的汉语不一样。}, souligne les mots-outils de la comparaison, indique si la phrase est affirmative ou négative, puis traduis-la.
\tcblower
Solution détaillée :
\begin{enumerate}
\item Les mots-outils sont \zh{跟} (gēn, « avec/comme ») et \zh{不一样} (bù yíyàng, « pas pareil »). La présence de \zh{不} devant \zh{一样} montre qu'il s'agit d'une comparaison \textbf{négative}.
\item La structure est : A (\zh{我的汉语}) + \zh{跟} + B (\zh{你的汉语}) + \zh{不一样}. Ici, il n'y a pas d'adjectif après \zh{不一样} : la phrase affirme simplement que les deux choses sont différentes.
\item Traduction : « Mon chinois n'est pas pareil au tien » (c'est-à-dire : « mon chinois est différent du tien »).
\end{enumerate}
Vérification : le terme de comparaison (\zh{你的汉语}) est bien placé \emph{avant} \zh{不一样} — conforme à notre hypothèse 1.
\end{exoresolu}

\begin{exercice}[À vous d'observer]
Dans chaque phrase, souligne les mots-outils et dis à quelle structure (comparaison, argumentation, hypothèse + opposition) elle appartient :
\begin{enumerate}
\item \zh{杜阿拉跟雅温得不一样大。}\emph{（Dù'ālā gēn Yǎwēndé bù yíyàng dà.）}
\item \zh{一方面我喜欢足球，另一方面我也喜欢篮球。}
\item \zh{如果明天下雨，我们就不去市场。}
\item \zh{她和我一样高。}
\end{enumerate}
\end{exercice}

\begin{corrige}[Exercice « À vous d'observer »]
\begin{enumerate}
\item Mots-outils : \zh{跟} et \zh{不一样}. Structure de \textbf{comparaison négative} : « Douala n'est pas aussi grande que Yaoundé. »
\item Mots-outils : \zh{一方面} et \zh{另一方面}. Structure d'\textbf{argumentation} : « D'un côté j'aime le football, de l'autre j'aime aussi le basketball. »
\item Mots-outils : \zh{如果} et \zh{就}. Structure d'\textbf{hypothèse} : « S'il pleut demain, nous n'irons pas au marché. »
\item Mots-outils : \zh{和} et \zh{一样}. Structure de \textbf{comparaison affirmative} : « Elle est aussi grande que moi. »
\end{enumerate}
\end{corrige}

\begin{aretenir}[Ce que l'observation nous a déjà appris]
\begin{itemize}
\item La comparaison d'égalité suit l'ordre \textbf{A + \zh{跟/和} + B + \zh{一样} (+ adjectif)} ; la négation se fait avec \zh{不} devant \zh{一样}.
\item L'argumentation en deux points utilise le couple \zh{一方面……另一方面……}.
\item L'hypothèse s'introduit avec \zh{如果} et la conséquence avec \zh{就} ; l'opposition s'exprime avec \zh{但是} ou \zh{可是}.
\item On ne traduit jamais mot à mot depuis le français : la place du terme de comparaison est différente dans les deux langues.
\end{itemize}
\end{aretenir}

Nos hypothèses sont posées. Dans la section suivante, nous formaliserons rigoureusement la première structure : \zh{A 跟/和 B (不)一样}.

\coursec{Structure 1 : A 跟/和 B (不)一样}

Dans la section précédente, nous avons observé des phrases comme \zh{喀麦隆跟中国不一样大。} et \zh{我的汉语水平跟你的差不多。} Nous avons remarqué que le mot \zh{一样} (\emph{yíyàng}, « pareil, identique ») revenait systématiquement, précédé de deux éléments comparés reliés par \zh{跟} ou \zh{和}. Nous allons maintenant étudier cette structure de manière rigoureuse : sa forme affirmative, sa forme négative, la place exacte de la négation \zh{不}, l'emploi avec ou sans adjectif, et la différence fondamentale avec la structure comparative en \zh{比} déjà rencontrée.

\courssub{Observation : que remarque-t-on ?}

\begin{exemplebox}[Quatre phrases à observer]
\begin{enumerate}
\item \zh{这本书跟那本书一样贵。} \emph{Zhè běn shū gēn nà běn shū yíyàng guì.} « Ce livre est aussi cher que celui-là. »
\item \zh{雅温得和杜阿拉不一样大。} \emph{Yǎwēndé hé Dù'ālā bù yíyàng dà.} « Yaoundé et Douala n'ont pas la même taille. »
\item \zh{我的房间跟你的不一样大。} \emph{Wǒ de fángjiān gēn nǐ de bù yíyàng dà.} « Ma chambre n'est pas aussi grande que la tienne. »
\item \zh{我们俩都一样高。} \emph{Wǒmen liǎ dōu yíyàng gāo.} « Nous sommes tous les deux aussi grands. »
\end{enumerate}
\end{exemplebox}

Observons attentivement. Dans les quatre phrases, on retrouve toujours le même squelette : deux éléments A et B, reliés par \zh{跟} ou \zh{和}, puis le mot \zh{一样}. Quand la phrase est négative, \zh{不} se place juste \textbf{avant} \zh{一样}, jamais ailleurs. Quand on précise le critère de la comparaison (cher, grand, haut), l'adjectif vient \textbf{après} \zh{一样}. C'est exactement ce que nous allons formaliser.

\courssub{La forme affirmative : A + 跟/和 + B + 一样 (+ adjectif)}

\begin{propriete}[Schéma de la structure affirmative]
\textbf{Comparaison globale (sans adjectif) :}

A + 跟/和 + B + 一样

\textbf{Comparaison sur un critère précis (avec adjectif) :}

A + 跟/和 + B + 一样 + adjectif

\zh{跟} (\emph{gēn}) et \zh{和} (\emph{hé}) sont parfaitement interchangeables dans cette structure ; \zh{跟} est plus fréquent à l'oral, \zh{和} un peu plus neutre. Le sens est : « A est pareil à B », « A est aussi (adjectif) que B ».
\end{propriete}

\begin{attention}[跟 et 和 sont ici des prépositions, pas des conjonctions]
Dans cette structure, \zh{跟/和} introduit le \textbf{terme de comparaison} : le groupe \zh{跟 B} / \zh{和 B} est un complément placé \textbf{avant} le prédicat \zh{一样}, comme dans \zh{我跟你说} « je te parle ». Conséquence importante : on ne peut pas intervertir librement A et B sans changer le sens (c'est A qu'on compare à B, et non l'inverse), et le sujet de la phrase reste A. Ne confondez pas avec \zh{和} conjonction « et » : \zh{我和他都很高。} « Lui et moi, nous sommes grands tous les deux » (deux sujets) n'est pas la même phrase que \zh{我跟他一样高。} « Je suis aussi grand que lui » (comparaison).
\end{attention}

\begin{popfigure}
\begin{tikzpicture}[node distance=0.35cm,
  bloc/.style={draw=popblue, fill=popblueL, rounded corners=3pt, minimum height=0.9cm, align=center, font=\small},
  fleche/.style={-{Stealth[length=2.5mm]}, thick, popdark}]
\node[bloc] (a) {A};
\node[bloc, right=of a] (gen) {跟 / 和};
\node[bloc, right=of gen] (b) {B};
\node[bloc, right=of b, draw=poporange, fill=poporangeL] (bu) {(不)};
\node[bloc, right=of bu, draw=poppurple, fill=poppurpleL] (yy) {一样};
\node[bloc, right=of yy, draw=popgreen, fill=popgreenL] (adj) {(+ adjectif)};
\draw[fleche] (a) -- (gen);
\draw[fleche] (gen) -- (b);
\draw[fleche] (b) -- (bu);
\draw[fleche] (bu) -- (yy);
\draw[fleche] (yy) -- (adj);
\end{tikzpicture}
\legende{Schéma linéaire de la structure : A + 跟/和 + B + (不) + 一样 (+ adjectif). Les éléments entre parenthèses sont facultatifs.}
\end{popfigure}

\begin{definition}[一样 \emph{yíyàng} : pareil, identique]
\zh{一样} est un adjectif qui signifie « pareil, identique, le même ». Dans la structure de comparaison, il joue le rôle du français « aussi ... que » ou « le/la même ... que ». Il est \textbf{obligatoire} : c'est lui qui porte l'idée d'égalité. Attention à la prononciation : \zh{一} seul se dit \emph{yī}, mais devant un quatrième ton il passe au deuxième ton : \emph{yíyàng}.
\end{definition}

\courssub{Emploi sans adjectif : la comparaison globale}

Quand aucun adjectif ne suit \zh{一样}, la comparaison porte sur l'ensemble : les deux éléments sont « pareils », « identiques », sans préciser sur quel point.

\begin{exemplebox}[Comparaison globale]
\begin{itemize}
\item \zh{这个词跟那个词一样。} \emph{Zhège cí gēn nàge cí yíyàng.} « Ce mot est identique à celui-là. »
\item \zh{我的想法和你的想法一样。} \emph{Wǒ de xiǎngfǎ hé nǐ de xiǎngfǎ yíyàng.} « Mon avis et ton avis sont les mêmes. »
\item \zh{今天的作业跟昨天的一样。} \emph{Jīntiān de zuòyè gēn zuótiān de yíyàng.} « Les devoirs d'aujourd'hui sont les mêmes que ceux d'hier. »
\end{itemize}
\end{exemplebox}

Remarquez dans le troisième exemple l'emploi de \zh{的} : \zh{昨天的} signifie « ceux d'hier » ; le nom (\zh{作业}) n'est pas répété, exactement comme le pronom possessif français « les miens, ceux de... ».

\courssub{Emploi avec adjectif : préciser le critère}

Le plus souvent, on ajoute un adjectif après \zh{一样} pour dire sur quel point les deux éléments sont égaux : \zh{一样高} « aussi grand (taille) », \zh{一样大} « aussi grand (dimension, âge) », \zh{一样贵} « aussi cher », \zh{一样多} « aussi nombreux », \zh{一样好} « aussi bon », \zh{一样远} « aussi loin ».

\begin{exemplebox}[Comparaison sur un critère précis]
\begin{itemize}
\item \zh{我弟弟跟我一样高。} \emph{Wǒ dìdi gēn wǒ yíyàng gāo.} « Mon petit frère est aussi grand que moi. »
\item \zh{杜阿拉和雅温得不一样热。} \emph{Dù'ālā hé Yǎwēndé bù yíyàng rè.} « Douala et Yaoundé ne sont pas aussi chaudes l'une que l'autre. »
\item \zh{这个市场的东西跟那个市场的一样多。} \emph{Zhège shìchǎng de dōngxi gēn nàge shìchǎng de yíyàng duō.} « Il y a autant de choses dans ce marché que dans celui-là. »
\item \zh{汉语和法语一样有意思。} \emph{Hànyǔ hé Fǎyǔ yíyàng yǒu yìsi.} « Le chinois est aussi intéressant que le français. »
\end{itemize}
\end{exemplebox}

\begin{attention}[Pas de 很 devant l'adjectif dans cette structure]
Après \zh{一样}, l'adjectif seul suffit : on dit \zh{我跟他一样高。}, jamais \zh{*我跟他一样很高。} En effet, \zh{一样} joue déjà le rôle de marqueur du prédicat adjectival ; ajouter \zh{很} serait un pléonasme fautif. De même, on ne met pas \zh{很} après \zh{比} : \zh{我比他高。}, jamais \zh{*我比他很高。}
\end{attention}

\begin{attention}[一样 + adjectif, pas 一样 + verbe psychologique]
La structure \zh{跟……一样} se combine presque toujours avec un \textbf{adjectif}. Dire \zh{*我跟他一样喜欢足球} pour « j'aime le football autant que lui » est très lourd ; on préfère \zh{我跟他一样，也喜欢足球。} \emph{Wǒ gēn tā yíyàng, yě xǐhuan zúqiú.} « Comme lui, j'aime aussi le football. » Retenez : après \zh{一样}, placez un adjectif, pas un verbe.
\end{attention}

\courssub{La forme négative : la place de 不}

\begin{propriete}[Règle de la place de 不]
Dans la forme négative, \zh{不} se place \textbf{immédiatement avant} \zh{一样} :

A + 跟/和 + B + \textbf{不} + 一样 (+ adjectif)

La négation porte sur \zh{一样} (« pas pareil »), pas sur \zh{跟}. On dit \zh{A 跟 B 不一样}, jamais \zh{*A 不跟 B 一样}. Le sens est : « A n'est pas pareil à B », « A n'est pas aussi (adjectif) que B ».
\end{propriete}

\begin{exemplebox}[Formes négatives]
\begin{itemize}
\item \zh{我的房间跟你的不一样大。} \emph{Wǒ de fángjiān gēn nǐ de bù yíyàng dà.} « Ma chambre n'est pas aussi grande que la tienne. »
\item \zh{今天跟昨天不一样冷。} \emph{Jīntiān gēn zuótiān bù yíyàng lěng.} « Aujourd'hui, il ne fait pas aussi froid qu'hier. »
\item \zh{这本书跟那本书不一样。} \emph{Zhè běn shū gēn nà běn shū bù yíyàng.} « Ce livre n'est pas pareil à celui-là. »
\end{itemize}
\end{exemplebox}

\begin{attention}[Piège : ne jamais dire *A 不跟 B 一样]
Sous l'influence du français (« A n'est pas comme B »), les francophones placent souvent la négation avant \zh{跟} : \zh{*我不跟他一样高。} C'est \textbf{incorrect}. En chinois, la négation nie le fait d'être « pareil », donc \zh{不} colle à \zh{一样} : \zh{我跟他不一样高。} « Je ne suis pas aussi grand que lui. »
\end{attention}

\courssub{Insistance avec 都}

\begin{propriete}[都 pour insister sur l'égalité]
Quand les deux éléments comparés forment un groupe (souvent un sujet pluriel), on peut ajouter \zh{都} (\emph{dōu}, « tous ») avant \zh{一样} pour insister : « tous les deux pareils ».

Sujet pluriel + 都 + 一样 (+ adjectif)
\end{propriete}

\begin{exemplebox}[都 + 一样]
\begin{itemize}
\item \zh{我们俩都一样高。} \emph{Wǒmen liǎ dōu yíyàng gāo.} « Nous sommes tous les deux aussi grands. »
\item \zh{这两个城市都一样热闹。} \emph{Zhè liǎng ge chéngshì dōu yíyàng rènao.} « Ces deux villes sont toutes les deux aussi animées. »
\end{itemize}
\end{exemplebox}

\courssub{一样 contre 比 : égalité ou supériorité ?}

Il ne faut pas confondre la structure en \zh{一样}, qui exprime l'\textbf{égalité}, avec la structure en \zh{比}, qui exprime la \textbf{supériorité} (ou l'infériorité). Le français utilise des tournures proches (« aussi grand que » / « plus grand que »), ce qui brouille les élèves : en chinois, les deux structures sont totalement distinctes et ne se mélangent jamais.

\begin{popfigure}
\begin{tikzpicture}[node distance=0.5cm and 2.2cm,
  tete/.style={rounded corners=4pt, minimum height=0.9cm, align=center, font=\small\bfseries, text=white},
  corps/.style={rounded corners=4pt, minimum height=2.6cm, align=center, font=\small, text width=5.2cm}]
\node[tete, fill=poporange] (t1) {比 : supériorité};
\node[tete, fill=popblue, right=of t1] (t2) {一样 : égalité};
\node[corps, fill=poporangeL, below=0.15cm of t1, align=center] (c1){A 比 B + adjectif\\[0.2cm]我比他高。\\ \emph{Wǒ bǐ tā gāo.}\\ « Je suis plus grand que lui. »};
\node[corps, fill=popblueL, below=0.15cm of t2, align=center] (c2){A 跟/和 B 一样 + adjectif\\[0.2cm]我跟他一样高。\\ \emph{Wǒ gēn tā yíyàng gāo.}\\ « Je suis aussi grand que lui. »};
\end{tikzpicture}
\legende{比 marque la supériorité (A dépasse B) ; 一样 marque l'égalité (A = B).}
\end{popfigure}

\begin{center}
\begin{tabularx}{\linewidth}{|X|X|X|}
\hline
\rowcolor{popblueL} Structure & Exemple (caractères + pinyin) & Traduction \\
\hline
A 比 B + Adj. & \zh{杜阿拉比雅温得热。} \emph{Dù'ālā bǐ Yǎwēndé rè.} & Douala est plus chaude que Yaoundé. \\
\hline
A 跟/和 B 一样 + Adj. & \zh{杜阿拉跟雅温得一样热。} \emph{Dù'ālā gēn Yǎwēndé yíyàng rè.} & Douala est aussi chaude que Yaoundé. \\
\hline
A 跟/和 B 不一样 + Adj. & \zh{杜阿拉跟雅温得不一样热。} \emph{Dù'ālā gēn Yǎwēndé bù yíyàng rè.} & Douala n'est pas aussi chaude que Yaoundé. \\
\hline
A 跟/和 B 一样 (sans Adj.) & \zh{这两本书一样。} \emph{Zhè liǎng běn shū yíyàng.} & Ces deux livres sont pareils. \\
\hline
Sujet pluriel + 都 + 一样 + Adj. & \zh{他们俩都一样大。} \emph{Tāmen liǎ dōu yíyàng dà.} & Ils ont tous les deux le même âge. \\
\hline
\end{tabularx}
\end{center}

\begin{attention}[Deux erreurs classiques des francophones]
\begin{enumerate}
\item \textbf{Oublier 一样} : dire \zh{*我跟他高。} pour « je suis aussi grand que lui » est incorrect. Sans \zh{一样}, la phrase ne veut rien dire. \zh{一样} est le cœur de la structure : il est \textbf{obligatoire}.
\item \textbf{Confondre 一样 et 一点儿} : \zh{一样} signifie « pareil », \zh{一点儿} signifie « un peu ». Une forme comme \zh{*一样一点儿} est impossible. Pour dire « un peu plus grand », on utilise \zh{比} : \zh{我比他高一点儿。} \emph{Wǒ bǐ tā gāo yìdiǎnr.} « Je suis un peu plus grand que lui. »
\end{enumerate}
\end{attention}

\begin{savaistu}[La variante soutenue 与……一样]
À l'écrit formel (articles, discours officiels), on rencontre aussi la structure \zh{与……一样}, où \zh{与} (\emph{yǔ}) est l'équivalent soutenu de \zh{跟/和} : \zh{中国与喀麦隆一样重视教育。} \emph{Zhōngguó yǔ Kāmàilóng yíyàng zhòngshì jiàoyù.} « La Chine attache autant d'importance que le Cameroun à l'éducation. » Au niveau Première, \zh{跟} et \zh{和} suffisent amplement ; sachez simplement reconnaître \zh{与} si vous le croisez dans un texte.
\end{savaistu}

\courssub{Exercice résolu}

\begin{exoresolu}[Traduisez en chinois]
Énoncé. Traduisez les phrases suivantes en chinois, en utilisant la structure \zh{跟/和……(不)一样} : 1. « Ma grande sœur est aussi grande que moi. » 2. « Le riz n'est pas aussi cher que le poulet. » 3. « Son école est pareille à la mienne. »
\tcblower
Solution détaillée.

1. On identifie les deux éléments : A = \zh{我姐姐} (ma grande sœur), B = \zh{我}. Critère : grand (taille) = \zh{高}. Phrase affirmative : \zh{我姐姐跟我一样高。} \emph{Wǒ jiějie gēn wǒ yíyàng gāo.}

2. Phrase négative : \zh{不} se place avant \zh{一样}. A = \zh{米饭} (le riz), B = \zh{鸡肉} (le poulet), critère : \zh{贵}. \zh{米饭跟鸡肉不一样贵。} \emph{Mǐfàn gēn jīròu bù yíyàng guì.} Attention : ne pas écrire \zh{*米饭不跟鸡肉一样贵。}

3. Comparaison globale, sans adjectif : \zh{他的学校跟我的一样。} \emph{Tā de xuéxiào gēn wǒ de yíyàng.} On remarque que le second \zh{学校} n'est pas répété : \zh{我的} suffit, comme « la mienne » en français.
\end{exoresolu}

\begin{exercice}[À vous de jouer]
Transformez les phrases suivantes selon le modèle indiqué entre parenthèses, puis traduisez-les en français.

1. \zh{这本书} / \zh{那本书} / \zh{贵} (affirmatif avec \zh{和})
2. \zh{今天} / \zh{昨天} / \zh{热} (négatif)
3. \zh{我} / \zh{他} / \zh{高} (affirmatif avec \zh{都} et sujet pluriel)
4. \zh{我的房间} / \zh{你的房间} (comparaison globale, négatif)
5. Traduisez en chinois : « Le football et le basketball sont aussi intéressants l'un que l'autre. »
\end{exercice}

\begin{corrige}[Exercice « À vous de jouer »]
1. \zh{这本书和那本书一样贵。} \emph{Zhè běn shū hé nà běn shū yíyàng guì.} « Ce livre est aussi cher que celui-là. »

2. \zh{今天跟昨天不一样热。} \emph{Jīntiān gēn zuótiān bù yíyàng rè.} « Aujourd'hui, il ne fait pas aussi chaud qu'hier. » (\zh{不} bien placé avant \zh{一样}.)

3. \zh{我们俩都一样高。} \emph{Wǒmen liǎ dōu yíyàng gāo.} « Nous sommes tous les deux aussi grands. »

4. \zh{我的房间跟你的房间不一样。} \emph{Wǒ de fángjiān gēn nǐ de fángjiān bù yíyàng.} « Ma chambre n'est pas pareille à la tienne. » (On peut aussi dire plus simplement \zh{我的房间跟你的不一样。})

5. \zh{足球和篮球一样有意思。} \emph{Zúqiú hé lánqiú yíyàng yǒu yìsi.} Critère : \zh{有意思} « intéressant » ; pas de \zh{很} après \zh{一样}.
\end{corrige}

\begin{aretenir}[L'essentiel de la structure]
\begin{itemize}
\item Affirmatif : A + 跟/和 + B + 一样 (+ adjectif) — « A est aussi ... que B ».
\item Négatif : \zh{不} se place \textbf{avant} \zh{一样} : A 跟 B 不一样 (+ adjectif).
\item \zh{跟} et \zh{和} sont interchangeables ; \zh{跟} domine à l'oral. Ce sont ici des prépositions introduisant le terme de comparaison.
\item Sans adjectif, la comparaison est globale (« pareil ») ; avec adjectif, elle précise le critère (\zh{一样高}, \zh{一样贵}, \zh{一样多}), sans jamais ajouter \zh{很}.
\item \zh{都} peut insister sur l'égalité dans un groupe : \zh{我们都一样。}
\item \zh{一样} = égalité ; \zh{比} = supériorité. Ne jamais mélanger les deux structures.
\end{itemize}
\end{aretenir}

\coursec{Structure 2 : 一方面……，另一方面……}

Après avoir appris à comparer deux éléments avec \zh{跟/和……（不）一样}, nous allons maintenant découvrir une structure qui permet de \cle{présenter deux aspects d'une même situation} : \zh{一方面……，另一方面……}. Bonne nouvelle pour les francophones : cette structure correspond presque mot à mot au français « d'une part... d'autre part » ou « d'un côté... de l'autre ». C'est l'une des structures les plus faciles à mémoriser, à condition de respecter quelques règles précises.

\courssub{Observation : découvrir la structure dans un texte}

\begin{textebox}[Pourquoi apprendre le chinois ?]
一方面汉语很有意思，另一方面汉语很有用。一方面我想去中国看看，另一方面我的钱不够。一方面学习汉语很难，另一方面我的老师帮助我。\\
\emph{Yì fāngmiàn Hànyǔ hěn yǒu yìsi, lìng yì fāngmiàn Hànyǔ hěn yǒuyòng. Yì fāngmiàn wǒ xiǎng qù Zhōngguó kànkan, lìng yì fāngmiàn wǒ de qián bù gòu. Yì fāngmiàn xuéxí Hànyǔ hěn nán, lìng yì fāngmiàn wǒ de lǎoshī bāngzhù wǒ.}\\
« D'une part le chinois est intéressant, d'autre part il est utile. D'un côté je veux aller voir la Chine, de l'autre je n'ai pas assez d'argent. D'un côté apprendre le chinois est difficile, de l'autre mon professeur m'aide. »
\end{textebox}

\textbf{Questions d'observation :}
\begin{enumerate}
\item Combien de fois le mot \zh{一方面} apparaît-il ? Et \zh{另一方面} ?
\item Où se placent-ils dans chaque proposition : au début, au milieu ou à la fin ?
\item Quelle ponctuation sépare les deux propositions ?
\item Les deux aspects présentés sont-ils toujours du même type (tous positifs, tous négatifs) ?
\end{enumerate}

Vous l'avez remarqué : \zh{一方面} ouvre toujours la première proposition et \zh{另一方面} ouvre toujours la seconde. Les deux propositions sont séparées par une virgule chinoise \zh{，}. Et les deux aspects peuvent être tous deux positifs (intéressant + utile), ou l'un positif et l'autre négatif (vouloir partir + manquer d'argent).

\courssub{La règle : forme et construction}

\begin{propriete}[Règle de construction]
La structure \zh{一方面……，另一方面……} encadre \cle{deux propositions complètes} :
\begin{center}
\textbf{一方面 + proposition 1 ，另一方面 + proposition 2 。}
\end{center}
\begin{itemize}
\item \zh{一方面} (yì fāngmiàn, « un aspect ») se place \textbf{en tête} de la première proposition ;
\item \zh{另一方面} (lìng yì fāngmiàn, « l'autre aspect ») se place \textbf{en tête} de la seconde proposition ;
\item les deux propositions sont séparées par la virgule chinoise \zh{，} ;
\item on n'utilise \textbf{jamais} \zh{一方面} seul sans son pendant \zh{另一方面} : la balance doit avoir ses deux plateaux.
\end{itemize}
\end{propriete}

\begin{definition}[Vocabulaire de la structure]
\begin{center}
\begin{tabularx}{\linewidth}{|X|X|X|X|}
\hline
\rowcolor{popblueL} Caractères & Pinyin & Nature & Traduction \\
\hline
一方面 & yì fāngmiàn & locution & d'une part, d'un côté \\
\hline
另一方面 & lìng yì fāngmiàn & locution & d'autre part, de l'autre côté \\
\hline
方面 & fāngmiàn & nom & aspect, côté, domaine \\
\hline
另 & lìng & adjectif & autre, un autre \\
\hline
\end{tabularx}
\end{center}
\end{definition}

\begin{attention}[Piège n° 1 : jamais 二方面 !]
Le second élément est \zh{另一方面} avec \zh{另} (lìng, « autre »), et \textbf{jamais} \zh{二方面}. En français on dit bien « d'\emph{une} part... d'\emph{autre} part » et non « d'une part... de \emph{deux} part » : la logique chinoise est identique. \zh{另} signifie « l'autre, un autre ». Écrire \zh{二方面} est une faute classique des élèves francophones qui raisonnent « un... deux ».
\end{attention}

\courssub{Emplois : deux aspects d'une même situation}

La structure sert à présenter \cle{deux aspects, deux raisons ou deux facettes} d'une même situation. Ces deux aspects peuvent être :

\begin{enumerate}
\item \textbf{Complémentaires et positifs} (deux avantages) :\\
\zh{一方面汉语很有意思，另一方面汉语很有用。}\\
\emph{Yì fāngmiàn Hànyǔ hěn yǒu yìsi, lìng yì fāngmiàn Hànyǔ hěn yǒuyòng.}\\
« D'une part le chinois est intéressant, d'autre part il est utile. »

\item \textbf{Complémentaires et négatifs} (deux difficultés) :\\
\zh{一方面汉字很多，另一方面发音也不容易。}\\
\emph{Yì fāngmiàn Hànzì hěn duō, lìng yì fāngmiàn fāyīn yě bù róngyì.}\\
« D'un côté il y a beaucoup de caractères, de l'autre la prononciation n'est pas facile non plus. »

\item \textbf{Contrastés} (un avantage, un inconvénient) :\\
\zh{一方面我想去中国，另一方面我的钱不够。}\\
\emph{Yì fāngmiàn wǒ xiǎng qù Zhōngguó, lìng yì fāngmiàn wǒ de qián bù gòu.}\\
« D'un côté je veux aller en Chine, de l'autre je n'ai pas assez d'argent. »
\end{enumerate}

\begin{exemplebox}[Exemples contextualisés Cameroun]
\zh{一方面雅温得的天气很舒服，另一方面交通不太方便。}\\
\emph{Yì fāngmiàn Yǎwēndé de tiānqì hěn shūfu, lìng yì fāngmiàn jiāotōng bú tài fāngbiàn.}\\
« D'un côté le temps à Yaoundé est agréable, de l'autre les transports ne sont pas très pratiques. »\\[4pt]
\zh{一方面杜阿拉的工作很多，另一方面房租很贵。}\\
\emph{Yì fāngmiàn Dù'ālā de gōngzuò hěn duō, lìng yì fāngmiàn fángzū hěn guì.}\\
« D'une part il y a beaucoup de travail à Douala, d'autre part les loyers sont chers. »\\[4pt]
\zh{一方面我喜欢吃恩多雷，另一方面我也喜欢吃中国菜。}\\
\emph{Yì fāngmiàn wǒ xǐhuan chī ēnduōléi, lìng yì fāngmiàn wǒ yě xǐhuan chī Zhōngguó cài.}\\
« D'un côté j'aime manger le ndolé, de l'autre j'aime aussi la cuisine chinoise. »
\end{exemplebox}

\courssub{Le sujet : identique ou différent}

\begin{propriete}[Le sujet dans les deux propositions]
Deux cas sont possibles :
\begin{itemize}
\item \textbf{Sujet identique} : le sujet peut être placé \cle{avant} \zh{一方面} ; il commande alors les deux propositions et n'a pas besoin d'être répété.\\
\zh{他一方面学习汉语，另一方面学习英语。}\\
\emph{Tā yì fāngmiàn xuéxí Hànyǔ, lìng yì fāngmiàn xuéxí Yīngyǔ.}\\
« D'une part il apprend le chinois, d'autre part il apprend l'anglais. »
\item \textbf{Sujets différents} : chaque proposition a son propre sujet, placé \cle{après} \zh{一方面} / \zh{另一方面}.\\
\zh{一方面我想学汉语，另一方面我的父母想让我学英语。}\\
\emph{Yì fāngmiàn wǒ xiǎng xué Hànyǔ, lìng yì fāngmiàn wǒ de fùmǔ xiǎng ràng wǒ xué Yīngyǔ.}\\
« D'un côté moi je veux apprendre le chinois, de l'autre mes parents veulent que j'apprenne l'anglais. »
\end{itemize}
\end{propriete}

\begin{attention}[Place du sujet : deux positions possibles]
Quand le sujet est commun, on peut dire :\\
\zh{我一方面学习，一方面工作。} (sujet avant, il gouverne les deux)\\
ou \zh{一方面我学习，另一方面我工作。} (sujet répété dans chaque proposition).\\
En revanche, si les sujets sont différents, ils doivent obligatoirement apparaître \textbf{après} \zh{一方面} et \zh{另一方面}. Ne dites jamais \zh{*我一方面他想……}.
\end{attention}

\courssub{Schéma en balance}

\begin{popfigure}
\begin{center}
\begin{tikzpicture}[scale=1]
% pivot
\fill[popdark] (0,0) circle (0.12);
\draw[popdark, very thick] (-0.5,-0.9) -- (0,0) -- (0.5,-0.9);
\draw[popdark, very thick] (-0.5,-0.9) -- (0.5,-0.9);
% flèche sujet
\node[fill=poppurpleL, rounded corners, align=center, font=\small] at (0,1.5) {Sujet commun\\ \zh{我} (wǒ)};
\draw[-{Stealth}, thick, poppurple] (0,1.15) -- (0,0.2);
% barre
\draw[popdark, very thick] (-3.2,0) -- (3.2,0);
% plateau gauche
\draw[popblue, thick] (-3.2,0) -- (-3.7,-0.7);
\draw[popblue, thick] (-3.2,0) -- (-2.7,-0.7);
\draw[popblue, thick] (-3.9,-0.7) arc (180:360:0.7);
\node[fill=popblueL, rounded corners, align=center, font=\small] at (-3.2,-1.7) {\zh{一方面} + P1\\ yì fāngmiàn\\ \zh{我想去中国}};
% plateau droit
\draw[poporange, thick] (3.2,0) -- (2.7,-0.7);
\draw[poporange, thick] (3.2,0) -- (3.7,-0.7);
\draw[poporange, thick] (2.5,-0.7) arc (180:360:0.7);
\node[fill=poporangeL, rounded corners, align=center, font=\small] at (3.2,-1.7) {\zh{另一方面} + P2\\ lìng yì fāngmiàn\\ \zh{我的钱不够}};
\end{tikzpicture}
\end{center}
\legende{La structure en balance : le sujet commun est le pivot, les deux propositions sont les deux plateaux.}
\end{popfigure}

\courssub{Comparaison avec le français}

\begin{center}
\begin{tabularx}{\linewidth}{|X|X|X|}
\hline
\rowcolor{popblueL} Chinois & Français & Remarque \\
\hline
一方面……，另一方面…… & d'une part... d'autre part & correspondance quasi directe \\
\hline
一方面……，另一方面…… & d'un côté... de l'autre (côté) & registre plus courant, identique \\
\hline
\zh{另} (lìng) & « autre » & jamais « deux » : pas de \zh{二方面} \\
\hline
virgule \zh{，} entre les propositions & virgule française & ponctuation parallèle \\
\hline
\end{tabularx}
\end{center}

C'est un \cle{point de repère facile} pour les francophones : l'ordre des idées est le même qu'en français. La seule vigilance concerne la place du sujet et l'emploi obligatoire des deux éléments ensemble.

\begin{attention}[Erreurs fréquentes des francophones]
\begin{itemize}
\item Utiliser \zh{一方面} seul : \zh{*一方面我想去中国。} est incomplet — il manque \zh{另一方面}.
\item Écrire \zh{*二方面} au lieu de \zh{另一方面}.
\item Oublier la virgule chinoise \zh{，} entre les deux propositions.
\item Traduire mot à mot « d'une part... mais de l'autre » : en chinois, on n'ajoute \textbf{pas} \zh{但是} dans cette structure ; le contraste est déjà porté par \zh{一方面……另一方面……}.
\end{itemize}
\end{attention}

\courssub{Méthode : rédiger un paragraphe argumentatif}

\begin{methode}[Construire un paragraphe avec 一方面……另一方面……]
\begin{enumerate}
\item \textbf{Poser le thème} en une phrase simple : \zh{学习汉语有好处也有困难。} \emph{Xuéxí Hànyǔ yǒu hǎochù yě yǒu kùnnan.} « Apprendre le chinois a des avantages et des difficultés. »
\item \textbf{Développer l'aspect 1} avec \zh{一方面} : \zh{一方面汉语很有意思，……}
\item \textbf{Développer l'aspect 2} avec \zh{另一方面} : \zh{另一方面汉字很难写。}
\item \textbf{Conclure} par une phrase de synthèse : \zh{但是我很喜欢汉语。} \emph{Dànshì wǒ hěn xǐhuan Hànyǔ.} « Mais j'aime beaucoup le chinois. »
\end{enumerate}
\end{methode}

\begin{exoresolu}[Transformation de phrases]
Énoncé : reliez les deux phrases suivantes avec \zh{一方面……另一方面……}.\\
\zh{喀麦隆的足球很有名。喀麦隆的音乐也很好。}
\tcblower
Solution : les deux phrases ont le même sujet (\zh{喀麦隆}, le Cameroun) et présentent deux aspects positifs complémentaires. On place donc le sujet commun en tête, puis les deux propositions :\\
\zh{喀麦隆一方面足球很有名，另一方面音乐也很好。}\\
\emph{Kāmàilóng yì fāngmiàn zúqiú hěn yǒumíng, lìng yì fāngmiàn yīnyuè yě hěn hǎo.}\\
« D'une part le football camerounais est très célèbre, d'autre part sa musique est très bonne aussi. »\\
On pouvait aussi répéter le sujet : \zh{一方面喀麦隆的足球很有名，另一方面喀麦隆的音乐也很好。}
\end{exoresolu}

\begin{exercice}[À vous]
Traduisez en chinois avec \zh{一方面……另一方面……} :
\begin{enumerate}
\item D'un côté j'aime Douala, de l'autre il y fait trop chaud.
\item D'une part mon frère apprend le chinois, d'autre part ma sœur apprend l'anglais.
\item D'un côté je veux regarder le match de football, de l'autre je dois étudier.
\end{enumerate}
\end{exercice}

\begin{corrige}[Exercice]
\begin{enumerate}
\item \zh{一方面我喜欢杜阿拉，另一方面那里太热了。} \emph{Yì fāngmiàn wǒ xǐhuan Dù'ālā, lìng yì fāngmiàn nàlǐ tài rè le.}
\item \zh{一方面我的哥哥学习汉语，另一方面我的姐姐学习英语。} \emph{Yì fāngmiàn wǒ de gēge xuéxí Hànyǔ, lìng yì fāngmiàn wǒ de jiějie xuéxí Yīngyǔ.} (sujets différents : chacun après sa locution)
\item \zh{一方面我想看足球比赛，另一方面我应该学习。} \emph{Yì fāngmiàn wǒ xiǎng kàn zúqiú bǐsài, lìng yì fāngmiàn wǒ yīnggāi xuéxí.}
\end{enumerate}
\end{corrige}

\begin{aretenir}[L'essentiel]
\begin{itemize}
\item Schéma : \textbf{一方面 + P1 ，另一方面 + P2 。}
\item Toujours les deux éléments ensemble ; jamais \zh{二方面}.
\item Sujet commun : avant \zh{一方面} ou répété ; sujets différents : après chaque locution.
\item Correspondance directe avec « d'une part... d'autre part » — un allié pour les francophones.
\end{itemize}
\end{aretenir}

\coursec{Structure 3 : 如果……,但是/可是……}

Dans la section précédente, nous avons appris à présenter les deux faces d'une même situation avec \zh{一方面……,另一方面……}. Nous allons maintenant découvrir une structure qui exprime non plus la coexistence de deux aspects, mais le \cle{contraste} entre une \cle{hypothèse} et une réalité qui la contrecarre : \zh{如果……,但是/可是……}. C'est l'équivalent chinois du français « même si..., cependant... » ou « si..., mais... ».

\courssub{Observation : découvrir la structure en contexte}

Lisons d'abord ce petit dialogue entre deux élèves d'un lycée de Yaoundé, la veille d'un contrôle de chinois.

\begin{textebox}[Dialogue : la veille du contrôle]
A : 明天要考试了,你准备好了吗？\\
\emph{Míngtiān yào kǎoshì le, nǐ zhǔnbèi hǎo le ma?}\\
« Demain il y a l'examen, tu es prêt ? »

B : 如果明天不下雨,可是我还是有一点儿紧张。\\
\emph{Rúguǒ míngtiān bú xiàyǔ, kěshì wǒ háishì yǒu yìdiǎnr jǐnzhāng.}\\
« Même s'il ne pleut pas demain, je suis quand même un peu nerveux. »

A : 如果你有时间,但是你不想复习,那怎么办呢？\\
\emph{Rúguǒ nǐ yǒu shíjiān, dànshì nǐ bù xiǎng fùxí, nà zěnme bàn ne?}\\
« Si tu as le temps mais que tu ne veux pas réviser, qu'est-ce qu'on fait alors ? »

B : 我很想复习！如果考试很难,可是我不会放弃。\\
\emph{Wǒ hěn xiǎng fùxí! Rúguǒ kǎoshì hěn nán, kěshì wǒ bú huì fàngqì.}\\
« J'ai vraiment envie de réviser ! Même si l'examen est difficile, je n'abandonnerai pas. »
\end{textebox}

Observons les phrases qui contiennent \zh{如果} : dans chacune, la première partie introduit une \textbf{hypothèse} avec \zh{如果} (« si »), et la seconde partie, introduite par \zh{但是} ou \zh{可是} (« mais, cependant »), affirme quelque chose qui \textbf{va contre la conséquence attendue}. Normalement, s'il ne pleut pas, on est content ; ici, l'élève est nerveux quand même. Normalement, si l'examen est difficile, on a envie d'abandonner ; ici, il refuse d'abandonner.

\courssub{La règle : forme et emploi}

\begin{propriete}[Schéma de la structure]
\textbf{如果 + hypothèse,但是/可是 + proposition principale contrastée}

\begin{itemize}
\item \zh{如果} (rúguǒ, « si ») introduit l'hypothèse. Il peut se placer \textbf{en tête de phrase} ou \textbf{après le sujet} : \zh{如果明天下雨} = \zh{明天如果下雨} (« s'il pleut demain »).
\item \zh{但是} (dànshì) et \zh{可是} (kěshì) introduisent la proposition principale et marquent l'\textbf{opposition} : la conséquence logique attendue est contrecarrée ou nuancée.
\item \zh{但是} et \zh{可是} sont \textbf{interchangeables} ; \zh{可是} est plus oral, \zh{但是} est plus neutre et plus fréquent à l'écrit.
\item On ajoute souvent \zh{还} (hái, « quand même, encore ») ou \zh{还是} (háishì, « quand même ») dans la proposition principale pour renforcer l'idée de persistance.
\end{itemize}
\end{propriete}

\begin{definition}[Emploi : hypothèse contrecarrée]
On utilise \zh{如果……,但是/可是……} pour dire « même si..., cependant... » : on imagine une situation possible (\zh{如果}), puis on affirme que la réalité ou la décision reste contraire à ce qu'on attendrait (\zh{但是/可是}). La nuance est celle de la \textbf{concession} ou de la \textbf{persistance malgré l'obstacle}.
\end{definition}

\begin{exemplebox}[Exemples traduits]
\begin{itemize}
\item \zh{如果明天下雨,可是我们还要上课。}\\
\emph{Rúguǒ míngtiān xiàyǔ, kěshì wǒmen hái yào shàngkè.}\\
« Même s'il pleut demain, nous aurons quand même cours. »
\item \zh{如果你有时间,但是你不想学习,那也没有办法。}\\
\emph{Rúguǒ nǐ yǒu shíjiān, dànshì nǐ bù xiǎng xuéxí, nà yě méiyǒu bànfǎ.}\\
« Si tu as le temps mais que tu ne veux pas étudier, on ne peut rien y faire. »
\item \zh{如果他很忙,可是他每天都学汉语。}\\
\emph{Rúguǒ tā hěn máng, kěshì tā měitiān dōu xué Hànyǔ.}\\
« Même s'il est occupé, il apprend le chinois tous les jours. »
\item \zh{我如果有钱,但是也不会买很贵的手机。}\\
\emph{Wǒ rúguǒ yǒu qián, dànshì yě bú huì mǎi hěn guì de shǒujī.}\\
« Même si j'avais de l'argent, je n'achèterais pas un téléphone très cher. » (ici \zh{如果} est placé après le sujet \zh{我})
\item \zh{如果杜阿拉离雅温得很远,可是很多人还是每天去那里工作。}\\
\emph{Rúguǒ Dù'ālā lí Yǎwēndé hěn yuǎn, kěshì hěn duō rén háishì měitiān qù nàli gōngzuò.}\\
« Même si Douala est loin de Yaoundé, beaucoup de gens y vont quand même travailler chaque jour. »
\end{itemize}
\end{exemplebox}

\courssub{如果……但是…… ou 如果……就…… ? Ne pas confondre !}

Vous connaissez déjà la structure simple \zh{如果……就……} (« si..., alors... »), où \zh{就} marque la \textbf{conséquence logique}. Dans notre nouvelle structure, \zh{但是/可是} marque au contraire l'\textbf{opposition}. Les deux connecteurs de la proposition principale s'excluent mutuellement.

\begin{propriete}[Règle d'incompatibilité : 就 OU 但是/可是, jamais les deux]
Dans la proposition principale, on choisit :
\begin{itemize}
\item soit \zh{就} → conséquence logique : \zh{如果明天下雨,我们就不去踢球。} (« S'il pleut demain, nous n'irons pas jouer au foot. »)
\item soit \zh{但是/可是} → opposition : \zh{如果明天下雨,可是我们还要去踢球。} (« Même s'il pleut demain, nous irons quand même jouer au foot. »)
\end{itemize}
La combinaison *\zh{如果……就但是……} est \textbf{incorrecte} : on ne peut pas marquer à la fois la conséquence et l'opposition dans la même proposition.
\end{propriete}

\begin{center}
\begin{tabularx}{\linewidth}{|X|X|X|}
\hline
\rowcolor{popblueL} Structure & Exemple (caractères + pinyin) & Traduction \\
\hline
如果……就…… (conséquence) & \zh{如果下雨,我就不出门。} \emph{Rúguǒ xiàyǔ, wǒ jiù bù chūmén.} & S'il pleut, je ne sors pas. \\
\hline
如果……但是/可是…… (opposition) & \zh{如果下雨,可是我还要出门。} \emph{Rúguǒ xiàyǔ, kěshì wǒ hái yào chūmén.} & Même s'il pleut, je sors quand même. \\
\hline
虽然……但是…… (fait réel) & \zh{虽然下雨了,但是我还出门了。} \emph{Suīrán xiàyǔ le, dànshì wǒ hái chūmén le.} & Bien qu'il ait plu, je suis quand même sorti. \\
\hline
\end{tabularx}
\end{center}

\begin{aretenir}[Pour l'examen : 如果……但是 ou 虽然……但是 ?]
\zh{虽然……但是……} exprime l'opposition sur un \textbf{fait réel, constaté} (« bien que..., mais... ») : \zh{虽然他很忙,但是他每天学汉语。} (« Bien qu'il soit occupé, il apprend le chinois chaque jour. » — c'est un fait). \zh{如果……但是……} porte sur une \textbf{hypothèse} (« même si..., ... quand même ») : \zh{如果他很忙,可是他还会学汉语。} (« Même s'il était occupé, il apprendrait quand même le chinois. » — on imagine le cas). À l'examen, repérez si la situation est réelle (\zh{虽然}) ou supposée (\zh{如果}).
\end{aretenir}

\begin{popfigure}
\begin{tikzpicture}[node distance=0.6cm and 1.4cm,
boite/.style={rectangle, rounded corners=4pt, draw=popblue, fill=popblueL, align=center, inner sep=6pt, font=\small},
boite2/.style={rectangle, rounded corners=4pt, draw=poporange, fill=poporangeL, align=center, inner sep=6pt, font=\small},
boite3/.style={rectangle, rounded corners=4pt, draw=popgreen, fill=popgreenL, align=center, inner sep=6pt, font=\small}]
\node[boite, align=center] (hyp){Hypothèse\\\zh{如果明天下雨}\\\emph{« s'il pleut demain »}};
\node[boite2, right=of hyp, yshift=1.2cm, align=center] (att){Conséquence attendue\\\zh{我们不上课}\\\emph{« pas de cours »}};
\node[boite3, right=of hyp, yshift=-1.2cm, align=center] (rea){Réalité contrastée\\\zh{可是我们还要上课}\\\emph{« cours quand même »}};
\draw[-{Stealth}, very thick, poporange] (hyp) -- (att);
\draw[popdark, very thick] ($(hyp)!0.5!(att)$) ++(-0.12,-0.18) -- ++(0.24,0.36);
\draw[popdark, very thick] ($(hyp)!0.5!(att)$) ++(-0.12,0.18) -- ++(0.24,-0.36);
\draw[-{Stealth}, very thick, popgreen] (hyp) -- (rea);
\end{tikzpicture}
\legende{Schéma de la structure : l'hypothèse introduite par 如果 ne produit pas la conséquence attendue (flèche barrée) ; la proposition principale en 但是/可是 affirme la réalité contrastée (flèche pleine).}
\end{popfigure}

\courssub{Comparaison avec le français et pièges des francophones}

Bonne nouvelle : le français possède aussi la combinaison « si..., mais... » (« Si tu as le temps, mais tu ne veux pas étudier... »), ce qui facilite le transfert. Attention cependant : l'\textbf{ordre chinois est rigide} — \zh{如果} ouvre toujours la subordonnée, \zh{但是/可是} ouvre toujours la principale — alors que le français permet plus de souplesse.

\begin{attention}[Erreurs fréquentes des francophones]
\begin{itemize}
\item \textbf{Inverser l'ordre} par calque du français « mais si... » : *\zh{但是如果明天下雨,我们还要上课。} est \textbf{incorrect}. L'ordre chinois est obligatoirement \zh{如果……,但是/可是……}.
\item \textbf{Cumuler} \zh{就} et \zh{但是} : *\zh{如果下雨,我就但是不去。} est impossible. On choisit l'un ou l'autre.
\item \textbf{Oublier le sujet} de la proposition principale quand il change : on dit \zh{如果你有时间,但是你不想学习} et non *\zh{如果你有时间,但是不想学习} si le sujet de la deuxième proposition doit être clair (ici, reprendre \zh{你}).
\item \textbf{Oublier la virgule chinoise} \zh{,} entre les deux propositions.
\end{itemize}
\end{attention}

\begin{savaistu}[可是 tout seul, à l'oral]
À l'oral, \zh{可是} peut s'employer \textbf{seul, en début de phrase}, pour exprimer une protestation ou une objection vive, sans structure complète : \zh{可是我不会!} \emph{Kěshì wǒ bú huì!} (« Mais je ne sais pas, moi ! ») ; \zh{可是我没有钱!} \emph{Kěshì wǒ méiyǒu qián!} (« Mais je n'ai pas d'argent ! »). C'est un marqueur très fréquent dans les conversations quotidiennes en Chine, un peu comme notre « mais... ! » indigné.
\end{savaistu}

\courssub{Exercice résolu et entraînement}

\begin{exoresolu}[Transformer avec 如果……但是/可是……]
Énoncé : transformez la phrase suivante en hypothèse contrastée. Phrase de départ : \zh{虽然天气很热,但是他还跑步。} (« Bien qu'il fasse très chaud, il court quand même. ») — Exprimez la même idée comme une supposition avec \zh{如果}.
\tcblower
Solution détaillée :
\begin{enumerate}
\item On remplace le fait réel (\zh{虽然}) par une hypothèse (\zh{如果}) : \zh{如果天气很热}.
\item On conserve le connecteur d'opposition dans la principale : \zh{但是} ou \zh{可是}.
\item On garde \zh{还} (« quand même ») qui renforce la persistance.
\end{enumerate}
Résultat : \zh{如果天气很热,可是他还跑步。} \emph{Rúguǒ tiānqì hěn rè, kěshì tā hái pǎobù.} « Même s'il fait très chaud, il court quand même. »
Vérification : pas de \zh{就} avec \zh{可是} ; \zh{如果} bien placé en tête de la subordonnée ; virgule chinoise entre les deux propositions.
\end{exoresolu}

\begin{exercice}[À vous de jouer]
\begin{enumerate}
\item Traduisez en chinois : « Même si le ndolé est un peu cher, j'en mange quand même chaque semaine. » (le ndolé se dit \zh{恩多莱} \emph{Ēnduōlái} en chinois)
\item Complétez avec \zh{就}, \zh{但是} ou \zh{可是} : \zh{如果明天有考试,我\_\_\_\_\_\_ 不去看电影。} puis \zh{如果明天有考试,\_\_\_\_\_\_ 我还想去踢球。}
\item Corrigez la phrase fautive : *\zh{但是如果他有时间,可是他不学习。}
\end{enumerate}
\end{exercice}

\begin{corrige}[Exercice « À vous de jouer »]
\begin{enumerate}
\item \zh{如果恩多莱有点儿贵,可是我还是每个星期都吃。} \emph{Rúguǒ Ēnduōlái yǒudiǎnr guì, kěshì wǒ háishì měi ge xīngqī dōu chī.} (On accepte \zh{但是} à la place de \zh{可是}.)
\item Première phrase : \zh{就} (conséquence logique : \zh{如果明天有考试,我就不去看电影。} « S'il y a un examen demain, je n'irai pas au cinéma. »). Seconde phrase : \zh{可是} ou \zh{但是} (opposition : \zh{如果明天有考试,可是我还想去踢球。} « Même s'il y a un examen demain, je veux quand même aller jouer au foot. »).
\item L'ordre est inversé et les connecteurs sont cumulés : il faut \zh{如果他有时间,但是他不学习。} ou \zh{如果他有时间,可是他不学习。} — \zh{如果} en tête de subordonnée, un seul connecteur d'opposition dans la principale.
\end{enumerate}
\end{corrige}

\begin{aretenir}[L'essentiel de la structure 3]
\begin{itemize}
\item Schéma : \textbf{如果 + hypothèse,但是/可是 + proposition contrastée} = « même si..., ... quand même ».
\item \zh{但是} (neutre, écrit) et \zh{可是} (oral) sont interchangeables.
\item \zh{如果} se place en tête de phrase ou après le sujet.
\item On choisit \zh{就} (conséquence) OU \zh{但是/可是} (opposition), jamais les deux.
\item \zh{还} / \zh{还是} renforcent l'idée de « quand même » dans la proposition principale.
\item Ne pas confondre avec \zh{虽然……但是……} (fait réel) : \zh{如果……但是……} reste une hypothèse.
\end{itemize}
\end{aretenir}

\coursec{Comparaison avec le français et pièges des francophones}

Après avoir analysé la structure conditionnelle \zh{如果……, 但是/可是……} (Section 4), nous devons maintenant prendre du recul. Les trois structures de cette leçon (\zh{跟/和……(不)一样}, \zh{一方面……, 另一方面……}, \zh{如果……, 但是/可是……}) partagent un point commun : elles organisent la pensée par \textbf{comparaison}, \textbf{nuance} ou \textbf{concession}. Or, le français et le chinois n'organisent pas ces relations logiques de la même manière. Cette section est votre \textbf{filet de sécurité} : elle recense les pièges où tombent 90\% des apprenants francophones et vous donne la méthode pour les éviter.

\courssub{Le piège de l'ordre des mots : le terme de comparaison « migre »}

En français, la structure comparative d'égalité place le terme de comparaison \textbf{après} l'adjectif et le mot-outil « que » :
\begin{center}
\textit{Sujet + verbe + aussi + Adjectif + \textbf{que} + Terme de comparaison}
\end{center}
\zh{Je suis aussi grand \textbf{que toi}.}

En chinois, le terme de comparaison (introduit par \zh{跟} ou \zh{和}) se place \textbf{avant} le prédicat \zh{一样} (+ adjectif). C'est une inversion totale de perspective : le chinois pose d'abord le \textbf{référent} (la règle de l'antériorité du thème/rhème), puis évalue le sujet par rapport à lui.

\begin{propriete}[Règle d'or : Position du terme de comparaison]
\textbf{Structure :} Sujet A + \zh{跟/和} + Sujet B + \zh{一样} (+ Adjectif) \\
\textbf{Traduction littérale :} « A \textbf{avec} B \textbf{est pareil} (grand). » \\
\textbf{Français :} « A est aussi adjectif \textbf{que} B. »
\end{propriete}

\begin{exemplebox}[Inversion Français $\leftrightarrow$ Chinois]
\begin{itemize}
    \item \zh{我跟你一样高。} \textit{(Wǒ gēn nǐ yīyàng gāo.)} — « Je suis aussi grand que toi. » \textit{[Litt. : Moi avec toi même grand.]}
    \item \zh{喀麦隆的天气跟中国南方的天气一样热。} \textit{(Kāmàilóng de tiānqì gēn Zhōngguó nánfāng de tiānqì yīyàng rè.)} — « Le temps au Cameroun est aussi chaud que dans le Sud de la Chine. »
    \item \zh{这道菜跟我妈妈做的一样好吃。} \textit{(Zhè dào cài gēn wǒ māma zuò de yīyàng hǎochī.)} — « Ce plat est aussi bon que celui que fait ma mère. »
\end{itemize}
\end{exemplebox}

\begin{attention}[Piège n°1 : L'ordre « Adjectif + 跟/和 » est INTERDIT]
\emph{Ne calquez jamais l'ordre français.}
\begin{center}
\begin{tabular}{ll}
\zh{* 我一样高跟你。} & \textbf{FAUX} (Ordre français calqué) \\
\zh{* 这本书一样贵跟那本。} & \textbf{FAUX} \\
\zh{\textbf{我跟你一样高。}} & \textbf{CORRECT} \\
\zh{\textbf{这本书跟那本一样贵。}} & \textbf{CORRECT}
\end{tabular}
\end{center}
\emph{Astuce mnémotechnique :} En chinois, on « pose le décor » (\zh{跟谁 ?}) avant de jouer la scène (\zh{一样}).
\end{attention}

\courssub{Le piège de la négation : \cle{不} ne va qu'avec \cle{一样}}

En français, la négation porte sur l'adverbe de degré (« \textbf{pas} aussi... que ») ou sur le verbe (« n'est \textbf{pas} aussi... que »). En chinois, la négation est \textbf{strictement attachée à \zh{一样}}. Elle ne précède jamais \zh{跟} ou \zh{和}.

\begin{propriete}[Place de la négation dans la comparaison d'égalité/différence]
\begin{itemize}
    \item \textbf{Égalité :} A \zh{跟/和} B \zh{一样} (+ Adj).
    \item \textbf{Différence (Inégalité) :} A \zh{跟/和} B \zh{\textbf{不}一样} (+ Adj).
\end{itemize}
Traduction : « A n'est \textbf{pas} aussi adjectif que B » ou « A est différent de B (sur le plan de l'adjectif) ».
\end{propriete}

\begin{exemplebox}[Négation correcte]
\begin{itemize}
    \item \zh{我跟他\textbf{不}一样高。} \textit{(Wǒ gēn tā \textbf{bù} yīyàng gāo.)} — « Je ne suis pas aussi grand que lui. »
    \item \zh{雅温得的冬天跟北京的冬天\textbf{不}一样冷。} \textit{(Yāwēndé de dōngtiān gēn Běijīng de dōngtiān \textbf{bù} yīyàng lěng.)} — « L'hiver à Yaoundé n'est pas aussi froid qu'à Pékin. »
    \item \zh{这件衣服跟那件\textbf{不}一样贵。} \textit{(Zhè jiàn yīfu gēn nà jiàn \textbf{bù} yīyàng guì.)} — « Ce vêtement n'est pas aussi cher que celui-là. »
\end{itemize}
\end{exemplebox}

\begin{attention}[Piège n°2 : \zh{*我不跟他一样高} — Erreur de place de \textbf{不}]
\begin{center}
\begin{tabular}{ll}
\zh{* 我\textbf{不}跟他一样高。} & \textbf{FAUX} (Négation portée sur le verbe/préposition) \\
\zh{* 跟他我\textbf{不}一样高。} & \textbf{FAUX} \\
\zh{\textbf{我跟他\textbf{不}一样高。}} & \textbf{CORRECT} (\zh{不} colle à \zh{一样})
\end{tabular}
\end{center}
\emph{Explication :} \zh{不跟} signifierait « ne pas suivre / ne pas être avec », ce qui change le sens. La structure \zh{跟/和...一样} est un bloc indivisible pour la comparaison ; \zh{不} doit s'insérer \textbf{à l'intérieur} du bloc, juste devant \zh{一样}.
\end{attention}

\courssub{Le piège du « double marqueur » et de l'oubli de \cle{一样}}

Deux erreurs symétriques sont fréquentes :
\begin{enumerate}
    \item \textbf{L'oubli de \zh{一样}} : On dit \zh{*我跟你高} (calque de « Je suis grand comme toi » ou confusion avec \zh{比}).
    \item \textbf{Le double marqueur \zh{*一样比}} : Mélange des deux structures de comparaison (égalité + supériorité).
\end{enumerate}

\begin{propriete}[Exclusion mutuelle : \zh{比} vs \zh{一样}]
\begin{center}
\begin{tabular}{|l|c|c|}
\hline
\textbf{Structure} & \textbf{Sens} & \textbf{Marqueurs} \\
\hline
\zh{A 比 B + Adj} & Supériorité (A > B) & \zh{比} seul, \textbf{pas de \zh{一样}} \\
\hline
\zh{A 跟/和 B 一样 (+ Adj)} & Égalité (A = B) & \zh{跟/和} + \zh{一样}, \textbf{pas de \zh{比}} \\
\hline
\zh{A 跟/和 B 不一样 (+ Adj)} & Inégalité (A $\neq$ B) & \zh{跟/和} + \zh{不一样}, \textbf{pas de \zh{比}} \\
\hline
\end{tabular}
\end{center}
\end{propriete}

\begin{attention}[Piège n°3 : \zh{*他比我一样高} — Mélange interdit]
\begin{center}
\begin{tabular}{ll}
\zh{* 他\textbf{比}我\textbf{一样}高。} & \textbf{FAUX} (Contradiction logique : \zh{比} implique différence, \zh{一样} implique égalité) \\
\zh{\textbf{他比我高。}} & \textbf{CORRECT} (Il est plus grand que moi.) \\
\zh{\textbf{他跟我一样高。}} & \textbf{CORRECT} (Il est aussi grand que moi.)
\end{tabular}
\end{center}
\emph{Règle de décision :} Avant de parler, demandez-vous : « Est-ce que je veux dire \textbf{« plus que »} (\zh{比}) ou \textbf{« autant que / pas autant que »} (\zh{一样}) ? » \textbf{Jamais les deux à la fois.}
\end{attention}

\courssub{Pièges lexicaux et structurels sur \zh{一方面……, 另一方面……}}

Cette structure correspond au français « D'un côté..., de l'autre côté... » ou « En même temps..., mais... ». Elle exprime une \textbf{analyse dialectique} (deux aspects d'une même réalité, souvent contradictoires ou complémentaires).

\begin{attention}[Piège n°4 : \zh{*二方面} au lieu de \zh{另一方面}]
\begin{center}
\begin{tabular}{ll}
\zh{* \textbf{一}方面我很忙，\textbf{二}方面我很高兴。} & \textbf{FAUX} (\zh{二方面} n'existe pas dans cette structure) \\
\zh{\textbf{一方面}我很忙，\textbf{另一方面}我很高兴。} & \textbf{CORRECT} (D'un côté je suis occupé, de l'autre je suis content.)
\end{tabular}
\end{center}
\emph{Explication :} \zh{一} (premier) appelle logiquement \zh{另} (l'autre / l'autre côté), pas \zh{二} (deuxième). \zh{另一方面} = « l'autre face / l'autre aspect ».
\end{attention}

\begin{exemplebox}[Utilisation correcte : Nuance et concession]
\begin{itemize}
    \item \zh{学汉字\textbf{一方面}难，\textbf{另一方面}很有意思。} \textit{(Xué hànzì \textbf{yīfāngmiàn} nán, \textbf{lìngfāngmiàn} hěn yǒuyìsi.)} — « Apprendre les caractères est d'un côté difficile, de l'autre intéressant. »
    \item \zh{去中国旅游\textbf{一方面}要花很多钱，\textbf{另一方面}能提高中文水平。} \textit{(Qù Zhōngguó lǚyóu \textbf{yīfāngmiàn} yào huā hěn duō qián, \textbf{lìngfāngmiàn} néng tígāo Zhōngwén shuǐpíng.)} — « Voyager en Chine coûte d'un côté beaucoup d'argent, de l'autre ça permet d'améliorer son chinois. »
\end{itemize}
\end{exemplebox}

\courssub{Pièges sur la structure conditionnelle-concessive \zh{如果……, 但是/可是……}}

En français, « Si..., mais... » est souvent maladroit ; on préfère « Si..., \textbf{cependant}... » ou « Bien que..., ... ». En chinois, \zh{如果} (condition) + \zh{但是/可是} (concession/réalité contrastée) est une structure standard pour exprimer : « \textbf{Même si} l'hypothèse se réalise, le fait principal ne change pas » ou « \textbf{Au cas où}..., \textbf{néanmoins}... ».

\begin{attention}[Piège n°5 : Cumul \zh{就} + \zh{但是} — Redondance logique]
Dans une phrase \zh{如果..., ...}, le mot \zh{就} (alors) marque la conséquence logique dans la proposition principale. \zh{但是/可是} marque l'opposition. Ils s'excluent mutuellement dans la même proposition.
\begin{center}
\begin{tabular}{ll}
\zh{* 如果下雨，我\textbf{就}\textbf{但是}不去。} & \textbf{FAUX} (\zh{就} = conséquence attendue ; \zh{但是} = rupture) \\
\zh{如果下雨，我\textbf{就}不去了。} & \textbf{CORRECT} (Condition $\rightarrow$ Conséquence logique) \\
\zh{如果下雨，我\textbf{但是/可是}还要去。} & \textbf{CORRECT} (Condition $\rightarrow$ Concession : malgré la pluie, j'y vais)
\end{tabular}
\end{center}
\emph{Astuce :} Si vous voulez dire « Même si..., quand même... », n'utilisez \textbf{pas} \zh{就}. Utilisez \zh{还是} (toujours/néanmoins) ou \zh{也} (aussi) avant le verbe principal : \zh{如果下雨，我\textbf{也/还是}要去。}
\end{attention}

\begin{exemplebox}[Distinction \zh{就} vs \zh{但是/可是} dans le contexte \zh{如果}]
\begin{itemize}
    \item \textbf{Conséquence logique (\zh{just}) :} \zh{如果明天有考试，我今晚\textbf{就}要复习。} \textit{(S'il y a un examen demain, je \textbf{devrai} réviser ce soir.)}
    \item \textbf{Concession / Volonté ferme (\zh{但是/可是} + \zh{还是/也}) :} \zh{如果明天下大雨，我\textbf{可是}\textbf{还是}要去踢足球。} \textit{(Même s'il pleut fort demain, je \textbf{tiens quand même} à aller jouer au foot.)}
\end{itemize}
\end{exemplebox}

\courssub{Synthèse visuelle : Tableau de transfert Français $\rightarrow$ Chinois}

Le tableau ci-dessous résume les correspondances et les erreurs types (marquées d'une croix rouge) pour les trois structures de la leçon.

\begin{popfigure}
\legende{Tableau de transfert contrastif : correspondances et pièges (croix rouges = erreurs fréquentes)}
\begin{tikzpicture}[
    node distance=0.8cm and 1.5cm,
    box/.style={rectangle, draw=popblue, thick, rounded corners=3pt, minimum height=1.1cm, text width=3.8cm, align=center, font=\small, fill=popblueL},
    arrow/.style={-Stealth, thick, popdark},
    labelbox/.style={font=\small\bfseries, text width=3.5cm, align=center},
    errbox/.style={draw=red, thick, rounded corners, minimum height=1.1cm, text width=3.8cm, align=center, font=\small, fill=poppinkL}
]
% Colonnes titres
\node[labelbox] (fr) at (0, 3.5) {\textbf{Français (Source)}};
\node[labelbox] (zh) at (6, 3.5) {\textbf{Chinois (Cible)}};
\node[labelbox] (pi) at (12, 3.5) {\textbf{Pièges à éviter}};

% Ligne 1 : Aussi... que
\node[box, align=center] (f1) at (0, 2.2){« Aussi Adj \textbf{que} B »\\(Égalité)};
\node[box, align=center] (z1) at (6, 2.2){\zh{A 跟/和 B 一样 (+ Adj)}\\(Ordre : A -- B -- 一样)};
\node[errbox, align=center] (p1) at (12, 2.2){\textcolor{red}{\Large $\times$} \zh{* A 一样 Adj 跟 B}\\ \zh{* A 不跟 B 一样 Adj}};

% Ligne 2 : Pas aussi... que
\node[box, align=center] (f2) at (0, 1.0){« \textbf{Pas} aussi Adj \textbf{que} B »\\(Inégalité / Infériorité)};
\node[box, align=center] (z2) at (6, 1.0){\zh{A 跟/和 B \textbf{不}一样 (+ Adj)}\\(\zh{不} \textbf{devant} \zh{一样})};
\node[errbox, align=center] (p2) at (12, 1.0){\textcolor{red}{\Large $\times$} \zh{* A 不跟 B 一样 Adj}\\ \zh{* A 跟 B 不 Adj (oubli \zh{一样})}};

% Ligne 3 : D'un côté / De l'autre
\node[box, align=center] (f3) at (0, -0.2){« D'un côté..., de l'autre... »\\(Dialectique)};
\node[box, align=center] (z3) at (6, -0.2){\zh{一方面..., 另一方面...}\\(\zh{一}... \zh{另}...)};
\node[errbox, align=center] (p3) at (12, -0.2){\textcolor{red}{\Large $\times$} \zh{* 一方面..., 二方面...}\\ \zh{* 一方面..., 另外...}};

% Ligne 4 : Si... mais / cependant
\node[box, align=center] (f4) at (0, -1.4){« Si..., \textbf{cependant}/pourtant... »\\(Concession conditionnelle)};
\node[box, align=center] (z4) at (6, -1.4){\zh{如果..., 但是/可是...}\\(Pas de \zh{就} !)};
\node[errbox, align=center] (p4) at (12, -1.4){\textcolor{red}{\Large $\times$} \zh{* 如果..., 就但是...}\\ \zh{* 如果..., 就可是...}};

% Flèches
\draw[arrow] (f1) -- (z1);
\draw[arrow] (f2) -- (z2);
\draw[arrow] (f3) -- (z3);
\draw[arrow] (f4) -- (z4);
\draw[arrow, poporange] (z1) -- (p1);
\draw[arrow, poporange] (z2) -- (p2);
\draw[arrow, poporange] (z3) -- (p3);
\draw[arrow, poporange] (z4) -- (p4);
\end{tikzpicture}
\end{popfigure}

\courssub{Méthode : Auto-correction en trois questions (Check-list mentale)}

Avant de valider une phrase contenant l'une de ces trois structures, appliquez systématiquement ce protocole. Il ne prend que 3 secondes et élimine 95\% des erreurs.

\begin{methode}[Protocole d'auto-correction « 3 Questions »]
\textbf{Étape 1 : Égalité ou Supériorité ?}
\begin{itemize}
    \item Si le sens est « A > B » (plus grand, plus cher, mieux) $\rightarrow$ \textbf{Structure \zh{比}} (\zh{A 比 B Adj}). \textbf{Interdit :} \zh{一样}, \zh{跟/和}.
    \item Si le sens est « A = B » (aussi... que) ou « A $\neq$ B » (pas aussi... que) $\rightarrow$ \textbf{Structure \zh{跟/和...一样}}.
\end{itemize}

\textbf{Étape 2 : Où est \zh{不} ? (Uniquement pour \zh{跟/和...一样})}
\begin{itemize}
    \item Vérifiez que \zh{不} est \textbf{immédiatement devant \zh{一样}}.
    \item \zh{跟/和 B \textbf{不}一样 Adj} = Correct.
    \item \zh{\textbf{不}跟/和 B 一样 Adj} = \textbf{FAUX}.
\end{itemize}

\textbf{Étape 3 : \zh{一样} est-il présent ? (Pour \zh{跟/和})}
\begin{itemize}
    \item \zh{跟/和} appelle \textbf{obligatoirement} \zh{一样} (ou \zh{不一样}).
    \item \zh{* A 跟 B Adj} (sans \zh{一样}) = \textbf{FAUX} (sauf verbes d'action comme \zh{跟 B 学习}, mais pas pour adjectifs d'état).
    \item \zh{* A 比 B 一样 Adj} = \textbf{FAUX} (Mélange interdit).
\end{itemize}

\textbf{Étape Bonus (Structures 2 \& 3) :}
\begin{itemize}
    \item \zh{一方面} $\rightarrow$ vérifiez la présence de \zh{另一方面} (pas \zh{二方面}).
    \item \zh{如果} + \zh{但是/可是} $\rightarrow$ vérifiez l'\textbf{absence de \zh{just}} dans la seconde proposition.
\end{itemize}
\end{methode}

\courssub{Atelier de correction collective : 5 phrases types}

Appliquez la méthode ci-dessus pour corriger ces 5 phrases « pièges » rédigées par des élèves de Première A. La correction détaillée suit dans l'exercice résolu.

\begin{exoresolu}[Atelier de correction : Repérez et corrigez les 5 erreurs]
\textbf{Consigne :} Chaque phrase contient \textbf{une seule erreur majeure} liée aux structures de la leçon. Identifiez l'erreur, citez la règle violée (Piège n°1 à 5) et réécrivez la phrase correcte.

\begin{enumerate}
    \item \zh{* 我弟弟跟我一样高。} (Contexte : Le frère cadet est \textbf{plus grand} que le locuteur.)
    \item \zh{* 这双鞋不跟那双一样贵。}
    \item \zh{* 学习中文一方面很难，二方面很有用。}
    \item \zh{* 如果你不喜欢吃辣，你就但是可以点这个菜。}
    \item \zh{* 北京的冬天比喀麦隆的冬天一样冷。}
\end{enumerate}

\tcblower

\textbf{Solution détaillée et analyse des pièges :}

\begin{enumerate}
    \item \textbf{Phrase :} \zh{* 我弟弟跟我一样高。} \\
    \textbf{Erreur :} \textbf{Piège n°3 (Choix structure : \zh{比} vs \zh{一样})}. Le contexte dit « plus grand » (supériorité), mais la structure utilisée est l'égalité (\zh{一样}). \\
    \textbf{Correction :} \zh{我弟弟\textbf{比}我高。} \textit{(Wǒ dìdi \textbf{bǐ} wǒ gāo.)} — Mon frère cadet est plus grand que moi. \\
    \textit{Note :} On ne dit pas \zh{比我一样高}.

    \item \textbf{Phrase :} \zh{* 这双鞋不跟那双一样贵。} \\
    \textbf{Erreur :} \textbf{Piège n°2 (Place de \zh{不})}. La négation \zh{不} est placée devant la préposition \zh{跟} au lieu de devant \zh{一样}. \\
    \textbf{Correction :} \zh{这双鞋跟那双\textbf{不}一样贵。} \textit{(Zhè shuāng xié gēn nà shuāng \textbf{bù} yīyàng guì.)} — Ces chaussures ne sont pas aussi chères que celles-là.

    \item \textbf{Phrase :} \zh{* 学习中文一方面很难，二方面很有用。} \\
    \textbf{Erreur :} \textbf{Piège n°4 (\zh{二方面} au lieu de \zh{另一方面})}. L'antonyme de \zh{一面/一方面} dans cette structure figée est \zh{另一面/另一方面}. \\
    \textbf{Correction :} \zh{学习中文\textbf{一方面}很难，\textbf{另一方面}很有用。} \textit{(Xuéxí Zhōngwén \textbf{yīfāngmiàn} hěn nán, \textbf{lìngfāngmiàn} hěn yǒuyòng.)} — Apprendre le chinois est d'un côté difficile, de l'autre utile.

    \item \textbf{Phrase :} \zh{* 如果你不喜欢吃辣，你就但是可以点这个菜。} \\
    \textbf{Erreur :} \textbf{Piège n°5 (Cumul \zh{就} + \zh{但是})}. \zh{就} marque la conséquence logique attendue ; \zh{但是} marque l'opposition. Ils s'annulent. Le sens ici est concessif (« Même si tu n'aimes pas le piment, tu \textbf{peux quand même} commander »). \\
    \textbf{Correction :} \zh{如果你不喜欢吃辣，你\textbf{还是/也}\textbf{可以}点这个菜。} \textit{(Rúguǒ nǐ bù xǐhuan chī là, nǐ \textbf{háishì/yě kěyǐ} diǎn zhège cài.)} — Si tu n'aimes pas le piment, tu peux quand même commander ce plat. \\
    \textit{Alternative (si opposition réelle) :} \zh{如果你不喜欢吃辣，\textbf{但是}这个菜不辣，你可以点。} (Si tu n'aimes pas le piment, \textbf{mais} ce plat n'est pas pimenté, tu peux le commander.)

    \item \textbf{Phrase :} \zh{* 北京的冬天比喀麦隆的冬天一样冷。} \\
    \textbf{Erreur :} \textbf{Piège n°3 (Double marqueur \zh{比} + \zh{一样})}. Mélange interdit de la structure de supériorité (\zh{比}) et d'égalité (\zh{一样}). Il faut choisir selon la réalité climatique (Pékin est \textbf{plus} froid que Yaoundé). \\
    \textbf{Correction :} \zh{北京的冬天\textbf{比}喀麦隆的冬天冷。} \textit{(Běijīng de dōngtiān \textbf{bǐ} Kāmàilóng de dōngtiān lěng.)} — L'hiver à Pékin est plus froid qu'au Cameroun. \\
    \textit{Si égalité (hypothétique) :} \zh{北京的冬天\textbf{跟}喀麦隆的冬天\textbf{一样}冷。}
\end{enumerate}
\end{exoresolu}

\courssub{Exercices d'application individuelle}

\begin{exercice}[Entraînement ciblé : Traduction et Transformation]
\textbf{A. Traduisez en chinois (caractères + pinyin) en évitant les pièges.}
\begin{enumerate}
    \item Le ndolé n'est pas aussi épicé que le mapo tofu (麻婆豆腐 \textit{mápó dòufu}).
    \item D'un côté, le chinois est difficile ; de l'autre, c'est très utile pour le travail.
    \item Si tu as le temps, viens quand même me voir à Douala. (Indice : concession \zh{还是/也}).
    \item Mon père est aussi grand que mon oncle (舅舅 \textit{jiùjiu} ou 叔叔 \textit{shūshu}).
    \item Ce téléphone n'est pas aussi cher que l'iPhone.
\end{enumerate}

\textbf{B. Transformez les phrases suivantes en utilisant la structure indiquée.}
\begin{enumerate}
    \item \zh{我哥哥比我高。} $\rightarrow$ Structure \zh{跟/和...不一样} (Exprimez : « Je ne suis pas aussi grand que mon frère »).
    \item \zh{这本书很贵，那本书也很贵。} $\rightarrow$ Structure \zh{跟/和...一样}.
    \item \zh{如果明天下雨，我们不去爬山。} $\rightarrow$ Structure \zh{如果..., 但是/可是...} (Exprimez : « Même s'il pleut, nous irons quand même »).
\end{enumerate}
\end{exercice}

\begin{corrige}[Exercice 9]
\textbf{A. Traduction}
\begin{enumerate}
    \item \zh{恩多莱 跟麻婆豆腐不一样辣。} \textit{(Ēnduōlái gēn Mápó dòufu bù yīyàng là.)} \textit{[Piège n°2 respecté : \zh{不} devant \zh{一样}]}.
    \item \zh{学中文\textbf{一方面}难，\textbf{另一方面}对工作很有用。} \textit{(Xué Zhōngwén \textbf{yīfāngmiàn} nán, \textbf{lìngfāngmiàn} duì gōngzuò hěn yǒuyòng.)} \textit{[Piège n°4 respecté : \zh{另} pas \zh{二}]}.
    \item \zh{如果你有时间，你\textbf{还是/也}来杜阿拉看我吧。} \textit{(Rúguǒ nǐ yǒu shíjiān, nǐ \textbf{háishì/yě} lái Dùālā kàn wǒ ba.)} \textit{[Piège n°5 respecté : pas de \zh{就}, usage de \zh{还是/也} pour concession]}.
    \item \zh{我爸爸跟我叔叔/舅舅\textbf{一样高}。} \textit{(Wǒ bàba gēn wǒ shūshu/jiùjiu \textbf{yīyàng gāo}.)} \textit{[Piège n°1 respecté : ordre A 跟 B 一样 Adj]}.
    \item \zh{这个手机跟iPhone\textbf{不}一样贵。} \textit{(Zhège shǒujī gēn iPhone \textbf{bù} yīyàng guì.)} \textit{[Piège n°2 respecté]}.
\end{enumerate}

\textbf{B. Transformation}
\begin{enumerate}
    \item \zh{我跟我哥哥\textbf{不}一样高。} \textit{(Wǒ gēn wǒ gēge \textbf{bù} yīyàng gāo.)} \textit{Note : La phrase source « 我哥哥比我高 » (Mon frère > Moi) implique « 我跟我哥哥不一样高 » (Moi < Mon frère). Le sujet change pour conserver la vérité du fait.}
    \item \zh{这本书跟那本书\textbf{一样贵}。} \textit{(Zhè běn shū gēn nà běn shū \textbf{yīyàng guì}.)}
    \item \zh{如果明天下雨，我们\textbf{但是/可是}\textbf{还是}要去爬山。} \textit{(Rúguǒ míngtiān xiàyǔ, wǒmen \textbf{dànshì/kěshì} \textbf{háishì} yào qù páshān.)} \textit{Structure concessionnelle complète.}
\end{enumerate}
\end{corrige}

\courssub{Bilan des points de vigilance (Check-list pour l'examen)}

\begin{aretenir}[Récapitulatif des 5 pièges mortels pour l'évaluation]
\begin{enumerate}
    \item \textbf{Ordre des mots :} \zh{Sujet A + 跟/和 + Sujet B + 一样 + Adj}. \textbf{Jamais} Adj avant \zh{跟/和}.
    \item \textbf{Négation :} \zh{不} \textbf{collé} à \zh{一样} (\zh{不一样}). \textbf{Jamais} \zh{不跟/不和}.
    \item \textbf{Présence de \zh{一样} :} Obligatoire après \zh{跟/和} pour les adjectifs. \textbf{Jamais} \zh{*跟 B Adj}.
    \item \textbf{Choix \zh{比} vs \zh{一样} :} \zh{比} = Supériorité (différence) ; \zh{一样} = Égalité. \textbf{Exclusifs}. Pas de \zh{*比...一样}.
    \item \textbf{Correspondances figées :} \zh{一方面} $\leftrightarrow$ \zh{另方面} (pas \zh{二方面}) ; \zh{如果} $\leftrightarrow$ \zh{但是/可是} (pas \zh{就}).
\end{enumerate}
\end{aretenir}

\begin{savaistu}[Le saviez-vous ? — La logique du « Même » (一样) dans la culture chinoise]
Le caractère \zh{样} (yàng) signifie à l'origine « modèle, forme, apparence, échantillon » (radical \zh{木} bois + phonétique \zh{羊} mouton : le bois taillé selon le modèle du mouton ?). \zh{一样} (yī yàng) signifie littéralement « un seul modèle / une seule forme ».
En chinois, dire \zh{A 跟 B 一样}, c'est dire : « A et B sont taillés dans le \textbf{même modèle} ». C'est une vision \textbf{essentialiste} de la comparaison : on ne mesure pas une différence de degré (comme \zh{比}), on constate une \textbf{identité de nature}.
C'est pourquoi \zh{不一样} (pas le même modèle) est si fort : cela implique une différence fondamentale, pas juste une question de « plus » ou « moins ». Cela explique aussi pourquoi on ne peut pas mélanger \zh{比} (mesure de degré) et \zh{一样} (identité de modèle) : c'est une contradiction logique dans la langue elle-même.
\end{savaistu}

\coursec{Exercices d'application et de transformation}

Après avoir comparé le chinois et le français et repéré les pièges qui guettent les francophones (ordre des mots, place de \zh{不}, oubli de l'adjectif après \zh{一样}), il est temps de passer à la pratique. Cette section suit une progression en quatre étapes : d'abord des exercices d'\cle{application guidée} (les mots-outils sont donnés), puis des exercices de \cle{transformation} (d'une structure à l'autre), ensuite de la \cle{traduction} (version et thème), et enfin une \cle{production libre} sur un sujet camerounais. C'est exactement le parcours d'une épreuve de chinois au lycée.

\begin{popfigure}
\begin{tikzpicture}[node distance=0.6cm]
\node[draw=popblue, fill=popblueL, rounded corners, minimum width=3.2cm, minimum height=1cm, align=center] (e1) {\textbf{Étape 1}\\ Application guidée};
\node[draw=poporange, fill=poporangeL, rounded corners, minimum width=3.2cm, minimum height=1cm, align=center, right=of e1] (e2) {\textbf{Étape 2}\\ Transformation};
\node[draw=popgreen, fill=popgreenL, rounded corners, minimum width=3.2cm, minimum height=1cm, align=center, right=of e2] (e3) {\textbf{Étape 3}\\ Traduction};
\node[draw=poppurple, fill=poppurpleL, rounded corners, minimum width=3.2cm, minimum height=1cm, align=center, right=of e3] (e4) {\textbf{Étape 4}\\ Production libre};
\draw[-{Stealth}, very thick, popdark] (e1) -- (e2);
\draw[-{Stealth}, very thick, popdark] (e2) -- (e3);
\draw[-{Stealth}, very thick, popdark] (e3) -- (e4);
\end{tikzpicture}
\legende{Frise de progressivité des exercices : du guidé vers le libre.}
\end{popfigure}

\begin{methode}[Transformer \zh{比} en \zh{一样} : la règle d'or]
Avant toute transformation, pose-toi une question de \emph{sens}, pas de grammaire :
\begin{enumerate}
\item Lis la phrase avec \zh{比} : elle exprime une \textbf{différence} (\zh{他比我高} « il est plus grand que moi »).
\item Demande-toi : l'égalité est-elle logique ou voulue ? \zh{A 比 B 高} ne peut PAS devenir \zh{A 跟 B 一样高} sans changer le sens !
\item La transformation n'est possible que dans deux cas : soit on inverse le rapport avec la négation (\zh{他比我高} → \zh{我不跟他一样高} « je ne suis pas aussi grand que lui »), soit le contexte affirme réellement l'égalité.
\item Sens inverse : \zh{A 跟 B 不一样大} peut devenir \zh{B 比 A 大} ou \zh{A 比 B 小} — choisis selon le sens.
\end{enumerate}
\end{methode}

\begin{attention}[Repères de traduction français → chinois]
Dans les traductions, repère les signaux français avant d'écrire :
\begin{itemize}
\item « aussi... que » → \zh{……跟/和……一样} (+ adjectif) ;
\item « pas aussi... que » → \zh{……跟/和……不一样} (+ adjectif) ;
\item « plus... que » → \zh{比} (ne pas mélanger les deux structures !) ;
\item « d'une part... d'autre part » → \zh{一方面……,另一方面……} ;
\item « si..., mais... » → \zh{如果……,但是/可是……}.
\end{itemize}
\end{attention}

\courssub{Exercice 1 : application guidée — compléter avec \zh{跟}/\zh{和}, \zh{一样}, \zh{不一样}}

\begin{exercice}[Complète les phrases]
Complète avec \zh{跟} ou \zh{和}, \zh{一样} ou \zh{不一样} (parfois les deux sont possibles).
\begin{enumerate}
\item 我的学校\_\_\_\_你的学校\_\_\_\_大。(Mes deux écoles sont aussi grandes l'une que l'autre.)
\item 今天的天气\_\_\_\_昨天\_\_\_\_热。(Aujourd'hui il ne fait pas aussi chaud qu'hier.)
\item 喀麦隆的足球\_\_\_\_中国的足球\_\_\_\_有名吗？(Le football camerounais est-il aussi célèbre que le football chinois ?)
\item 这本书\_\_\_\_那本书\_\_\_\_有意思。(Ce livre est aussi intéressant que celui-là.)
\item 我妹妹\_\_\_\_我\_\_\_\_高，她比我矮。(Ma petite sœur n'est pas aussi grande que moi.)
\item 杜阿拉的天气\_\_\_\_雅温得的天气\_\_\_\_。(Le climat de Douala n'est pas le même que celui de Yaoundé.)
\end{enumerate}
\end{exercice}

\begin{corrige}[Exercice 1]
\begin{enumerate}
\item \zh{我的学校跟/和你的学校一样大。}\\
\emph{Wǒ de xuéxiào gēn / hé nǐ de xuéxiào yíyàng dà.}\\
« Mon école est aussi grande que la tienne. » — Égalité affirmative : \zh{一样}.
\item \zh{今天的天气跟/和昨天不一样热。}\\
\emph{Jīntiān de tiānqì gēn / hé zuótiān bù yíyàng rè.}\\
« Aujourd'hui il ne fait pas aussi chaud qu'hier. » — « pas aussi... que » : \zh{不} avant \zh{一样}.
\item \zh{喀麦隆的足球跟/和中国的足球一样有名吗？}\\
\emph{Kāmàilóng de zúqiú gēn / hé Zhōngguó de zúqiú yíyàng yǒumíng ma?}\\
« Le football camerounais est-il aussi célèbre que le football chinois ? » — Question : \zh{吗} final, structure inchangée.
\item \zh{这本书跟/和那本书一样有意思。}\\
\emph{Zhè běn shū gēn / hé nà běn shū yíyàng yǒu yìsi.}\\
« Ce livre est aussi intéressant que celui-là. » — Adjectif dissyllabique possible après \zh{一样}.
\item \zh{我妹妹跟/和我不一样高。}\\
\emph{Wǒ mèimei gēn / hé wǒ bù yíyàng gāo.}\\
« Ma petite sœur n'est pas aussi grande que moi. » — La suite « elle est plus petite que moi » confirme la négation.
\item \zh{杜阿拉的天气跟/和雅温得的天气不一样。}\\
\emph{Dùā'lā de tiānqì gēn / hé Yǎwēndé de tiānqì bù yíyàng.}\\
« Le climat de Douala n'est pas le même que celui de Yaoundé. » — Sans adjectif : « n'est pas le même que ».
\end{enumerate}
\end{corrige}

\courssub{Exercice 2 : thème — traduire en chinois}

\begin{exoresolu}[Traduis en chinois]
Énoncé : traduis les deux phrases suivantes. a) « Ma ville est aussi grande que ta ville. » b) « Aujourd'hui il ne fait pas aussi froid qu'hier. »
\tcblower
Solution détaillée :
\begin{enumerate}
\item Phrase a) : signal « aussi... que » → \zh{跟/和……一样} + adjectif \zh{大}. Sujet : \zh{我的城市} ; terme de comparaison : \zh{你的城市}.\\
\zh{我的城市跟你的城市一样大。}\\
\emph{Wǒ de chéngshì gēn nǐ de chéngshì yíyàng dà.}
\item Phrase b) : signal « pas aussi... que » → \zh{不} + \zh{一样} + adjectif \zh{冷}. En chinois, « aujourd'hui » suffit, pas besoin de « il fait ».\\
\zh{今天跟昨天不一样冷。}\\
\emph{Jīntiān gēn zuótiān bù yíyàng lěng.}\\
(ou \zh{今天的天气跟昨天不一样冷。} \emph{Jīntiān de tiānqì gēn zuótiān bù yíyàng lěng.})
\end{enumerate}
Piège évité : ne pas écrire \zh{*今天比昨天不冷} — mélange interdit des structures \zh{比} et \zh{不一样}.
\end{exoresolu}

\courssub{Exercice 3 : transformation — de \zh{比} à \zh{一样} et inversement}

\begin{exercice}[Transforme quand c'est possible]
Transforme chaque phrase : de \zh{比} vers \zh{(不)一样}, ou l'inverse. Si la transformation change le sens, explique pourquoi.
\begin{enumerate}
\item \zh{我哥哥比我高。}(Mon grand frère est plus grand que moi.)
\item \zh{汉语跟法语不一样难。}(Le chinois n'est pas aussi difficile que le français.)
\item \zh{这个房间跟那个房间一样干净。}(Cette chambre est aussi propre que celle-là.)
\item \zh{杜阿拉比雅温得热。}(Douala est plus chaude que Yaoundé.)
\end{enumerate}
\end{exercice}

\begin{corrige}[Exercice 3]
\begin{enumerate}
\item \zh{我不跟哥哥一样高。}\\
\emph{Wǒ bù gēn gēge yíyàng gāo.}\\
« Je ne suis pas aussi grand que mon grand frère. » — On inverse le rapport avec la négation : le sens est conservé. (On ne peut PAS écrire \zh{*我哥哥跟我一样高} : ce serait faux !)
\item \zh{法语比汉语难。}\\
\emph{Fǎyǔ bǐ Hànyǔ nán.}\\
« Le français est plus difficile que le chinois. » — « Pas aussi difficile que » = « l'autre est plus difficile » : transformation correcte.
\item \zh{这个房间不比那个房间脏。}\\
\emph{Zhège fángjiān bù bǐ nàge fángjiān zāng.}\\
« Cette chambre n'est pas plus sale que celle-là. » — Transformation possible seulement en changeant d'adjectif (propre ↔ sale) ; sinon on garde \zh{一样}, plus naturel.
\item \zh{雅温得跟杜阿拉不一样热。}\\
\emph{Yǎwēndé gēn Dùā'lā bù yíyàng rè.}\\
« Yaoundé n'est pas aussi chaude que Douala. » — Inversion + négation : sens conservé.
\end{enumerate}
\end{corrige}

\begin{exemplebox}[Corrigé type à retenir]
\zh{雅温得跟杜阿拉不一样大。}\\
\emph{Yǎwēndé gēn Dùā'lā bù yíyàng dà.}\\
« Yaoundé n'est pas aussi grande que Douala. »\\
Remarque l'ordre : sujet + \zh{跟} + terme de comparaison + \zh{不} + \zh{一样} + adjectif. C'est le schéma à reproduire dans toutes tes productions.
\end{exemplebox}

\courssub{Exercice 4 : application — relier avec \zh{一方面……,另一方面……}}

\begin{exercice}[Relie les deux propositions]
Relie chaque paire d'idées avec \zh{一方面……,另一方面……}.
\begin{enumerate}
\item École : notre école est grande. / Notre école a de bons professeurs.
\item Famille : mon père travaille à Douala. / Ma mère travaille à Yaoundé.
\item Sport : jouer au football est bon pour la santé. / Jouer au football me rend heureux.
\item Études : j'aime le chinois. / Le chinois est utile pour l'avenir.
\end{enumerate}
\end{exercice}

\begin{corrige}[Exercice 4]
\begin{enumerate}
\item \zh{我们学校一方面很大，另一方面有很好的老师。}\\
\emph{Wǒmen xuéxiào yìfāngmiàn hěn dà, lìng yìfāngmiàn yǒu hěn hǎo de lǎoshī.}\\
« Notre école est d'une part grande, d'autre part elle a de très bons professeurs. »
\item \zh{一方面我爸爸在杜阿拉工作，另一方面我妈妈在雅温得工作。}\\
\emph{Yìfāngmiàn wǒ bàba zài Dùā'lā gōngzuò, lìng yìfāngmiàn wǒ māma zài Yǎwēndé gōngzuò.}\\
« D'une part mon père travaille à Douala, d'autre part ma mère travaille à Yaoundé. » — Quand les sujets sont différents, \zh{一方面} se place avant chaque sujet.
\item \zh{踢足球一方面对身体好，另一方面让我很高兴。}\\
\emph{Tī zúqiú yìfāngmiàn duì shēntǐ hǎo, lìng yìfāngmiàn ràng wǒ hěn gāoxìng.}\\
« Jouer au football est d'une part bon pour le corps, d'autre part cela me rend heureux. »
\item \zh{我一方面喜欢汉语，另一方面汉语对我的未来很有用。}\\
\emph{Wǒ yìfāngmiàn xǐhuan Hànyǔ, lìng yìfāngmiàn Hànyǔ duì wǒ de wèilái hěn yǒuyòng.}\\
« J'aime d'une part le chinois, d'autre part le chinois est très utile pour mon avenir. »
\end{enumerate}
\end{corrige}

\courssub{Exercice 5 : transformation — de \zh{如果……就……} à \zh{如果……但是/可是……}}

\begin{exercice}[Change la conséquence en opposition]
Transforme chaque phrase : remplace la conséquence logique (\zh{就}) par une opposition (\zh{但是/可是}).
\begin{enumerate}
\item \zh{如果明天下雨，我们就不去运动场。}(S'il pleut demain, nous n'irons pas au terrain de sport.)
\item \zh{如果我有时间，我就学习汉语。}(Si j'ai le temps, j'étudie le chinois.)
\item \zh{如果他去中国，他就会说汉语了。}(S'il va en Chine, il saura parler chinois.)
\end{enumerate}
\end{exercice}

\begin{corrige}[Exercice 5]
\begin{enumerate}
\item \zh{如果明天下雨，但是我们还是去运动场。}\\
\emph{Rúguǒ míngtiān xiàyǔ, dànshì wǒmen háishì qù yùndòngchǎng.}\\
« S'il pleut demain, nous irons quand même au terrain. » — \zh{还是} renforce l'opposition.
\item \zh{如果我有时间，可是我不能学习汉语，因为我要帮妈妈。}\\
\emph{Rúguǒ wǒ yǒu shíjiān, kěshì wǒ bù néng xuéxí Hànyǔ, yīnwèi wǒ yào bāng māma.}\\
« Si j'ai le temps, mais je ne peux pas étudier le chinois, car je dois aider maman. »
\item \zh{如果他去中国，但是他不想学习汉语。}\\
\emph{Rúguǒ tā qù Zhōngguó, dànshì tā bù xiǎng xuéxí Hànyǔ.}\\
« S'il va en Chine, mais il ne veut pas apprendre le chinois. »
\end{enumerate}
Règle appliquée : \zh{如果} introduit l'hypothèse ; \zh{就} annonce une conséquence attendue, \zh{但是/可是} annonce un obstacle ou un contraste. On supprime toujours \zh{就} quand on met \zh{但是/可是}.
\end{corrige}

\courssub{Exercice 6 : version — traduire un mini-texte}

\begin{textebox}[Les deux villes de mon pays]
雅温得跟杜阿拉不一样。雅温得的天气跟杜阿拉不一样热，可是杜阿拉比雅温得大。\\
一方面我喜欢雅温得，因为那里有我的学校；另一方面我也喜欢杜阿拉，因为那里有海。\\
如果我有时间，但是我没有钱，我就不能去杜阿拉旅行。\\
这两个城市都跟中国的城市不一样，它们有自己的特点。\\
\emph{Yǎwēndé gēn Dùā'lā bù yíyàng. Yǎwēndé de tiānqì gēn Dùā'lā bù yíyàng rè, kěshì Dùā'lā bǐ Yǎwēndé dà. Yìfāngmiàn wǒ xǐhuan Yǎwēndé, yīnwèi nàli yǒu wǒ de xuéxiào; lìng yìfāngmiàn wǒ yě xǐhuan Dùā'lā, yīnwèi nàli yǒu hǎi. Rúguǒ wǒ yǒu shíjiān, dànshì wǒ méiyǒu qián, wǒ jiù bù néng qù Dùā'lā lǚxíng. Zhè liǎng gè chéngshì dōu gēn Zhōngguó de chéngshì bù yíyàng, tāmen yǒu zìjǐ de tèdiǎn.}
\end{textebox}

\begin{exercice}[Traduis le texte en français]
Traduis le texte ci-dessus. Souligne dans ta traduction les trois structures étudiées.
\end{exercice}

\begin{corrige}[Exercice 6]
« Yaoundé et Douala ne sont pas pareilles. Le climat de Yaoundé n'est pas aussi chaud que celui de Douala, mais Douala est plus grande que Yaoundé. D'une part, j'aime Yaoundé parce que mon école s'y trouve ; d'autre part, j'aime aussi Douala parce qu'il y a la mer. Si j'ai le temps, mais que je n'ai pas d'argent, je ne peux pas aller voyager à Douala. Ces deux villes ne ressemblent pas aux villes chinoises : elles ont leurs propres caractéristiques. »\\
Structures repérées : \zh{跟……不一样} (×3), \zh{比} (×1), \zh{一方面……另一方面} (×1), \zh{如果……但是……} (×1).
\end{corrige}

\courssub{Exercice 7 : production libre — comparer Yaoundé et Douala}

\begin{exercice}[Rédige cinq phrases]
Rédige cinq phrases comparant Yaoundé et Douala. Utilise au moins une fois chaque structure : \zh{跟/和……一样}(+ adjectif), \zh{跟/和……不一样}(+ adjectif), \zh{一方面……另一方面……}, \zh{如果……但是/可是……}.
\end{exercice}

\begin{corrige}[Exercice 7 — production possible]
\begin{enumerate}
\item \zh{杜阿拉跟雅温得不一样大，杜阿拉比较大。}\\
\emph{Dùā'lā gēn Yǎwēndé bù yíyàng dà, Dùā'lā bǐjiào dà.}\\
« Douala n'est pas de la même taille que Yaoundé, Douala est plus grande. »
\item \zh{这两个城市跟中国的城市一样有意思。}\\
\emph{Zhè liǎng gè chéngshì gēn Zhōngguó de chéngshì yíyàng yǒu yìsi.}\\
« Ces deux villes sont aussi intéressantes que les villes chinoises. »
\item \zh{雅温得的天气跟杜阿拉不一样热。}\\
\emph{Yǎwēndé de tiānqì gēn Dùā'lā bù yíyàng rè.}\\
« Il ne fait pas aussi chaud à Yaoundé qu'à Douala. »
\item \zh{一方面杜阿拉有海，另一方面雅温得有我的家。}\\
\emph{Yìfāngmiàn Dùā'lā yǒu hǎi, lìng yìfāngmiàn Yǎwēndé yǒu wǒ de jiā.}\\
« D'une part Douala a la mer, d'autre part Yaoundé a ma maison. »
\item \zh{如果杜阿拉很热，但是那里的人很快乐。}\\
\emph{Rúguǒ Dùā'lā hěn rè, dànshì nàli de rén hěn kuàilè.}\\
« S'il fait chaud à Douala, les gens y sont pourtant joyeux. »
\end{enumerate}
\end{corrige}

\courssub{Correction collective : justifier chaque règle}

\begin{propriete}[Barème type examen]
À l'examen, chaque phrase de comparaison est notée sur les mots-outils : \textbf{1 point par mot-outil correctement placé}.
\begin{itemize}
\item \zh{跟} ou \zh{和} entre les deux termes : 1 point ;
\item \zh{不} placé \emph{avant} \zh{一样} (jamais après, jamais avant \zh{跟}) : 1 point ;
\item \zh{一样} présent : 1 point ;
\item adjectif après \zh{一样} quand le sens l'exige : 1 point ;
\item pour \zh{一方面……另一方面} : les deux moitiés présentes : 2 points ;
\item pour \zh{如果……但是/可是} : absence de \zh{就} : 1 point.
\end{itemize}
\end{propriete}

\begin{methode}[Correction collective en classe]
\begin{enumerate}
\item Un élève écrit sa phrase au tableau ; la classe identifie d'abord la \textbf{structure} utilisée.
\item On vérifie ensuite chaque mot-outil un par un, selon le barème ci-dessus.
\item On justifie à voix haute : « \zh{不} est avant \zh{一样} parce que... », « on n'a pas mis \zh{就} parce que \zh{但是} exprime l'opposition... ».
\item On propose enfin une variante correcte (par exemple la même idée avec \zh{比} au lieu de \zh{不一样}).
\end{enumerate}
\end{methode}

\begin{attention}[Dernières vérifications avant de rendre ta copie]
\begin{itemize}
\item Jamais \zh{*不一样跟} ni \zh{*跟……比} : les structures \zh{跟……一样} et \zh{比} ne se mélangent pas.
\item \zh{不} se colle à \zh{一样} : \zh{不一样}, pas \zh{*不跟……一样}.
\item Après \zh{如果……但是/可是}, supprime bien \zh{就}.
\item \zh{一方面……另一方面} présente deux aspects du \emph{même} sujet ; ce n'est pas une opposition vraie (pour opposer, utilise \zh{但是}).
\end{itemize}
\end{attention}

\begin{aretenir}[L'essentiel de la séance d'exercices]
\begin{itemize}
\item « aussi... que » → \zh{A 跟/和 B 一样 + Adj} ; « pas aussi... que » → \zh{A 跟/和 B 不一样 + Adj} ; « plus... que » → \zh{A 比 B + Adj}.
\item Transformer \zh{比} en \zh{一样} exige d'inverser les termes et d'ajouter \zh{不} pour garder le sens.
\item \zh{一方面……,另一方面……} = « d'une part... d'autre part » (deux aspects complémentaires).
\item \zh{如果……,但是/可是……} = hypothèse + opposition ; on supprime \zh{就}.
\end{itemize}
\end{aretenir}

\newpage

\coursec{Activité d'intégration}

\coursec{Leçon 20 — Grammaire : 比较字句 \zh{跟/和……(不)一样} 与 \zh{一方面……另一方面……} 与 \zh{如果……但是/可是……}}

\courssub{Objectifs de la leçon}

À l'issue de cette leçon, l'élève sera capable de :
\begin{itemize}
\item \textbf{Comprendre} et \textbf{produire} les trois structures cibles : comparaison d'égalité/différence (\zh{A 跟/和 B (不)一样 [+ 形容词]}), énumération de deux aspects contrastés (\zh{一方面……,另一方面……}), hypothèse conditionnelle avec concession/restriction (\zh{如果……,但是/可是……}) ;
\item \textbf{Situer} l'ordre des mots chinois par rapport au français (place de l'adjectif après \zh{一样}, place de \zh{但是/可是} en début de proposition principale) ;
\item \textbf{Mobiliser} le vocabulaire du thème « Citoyenneté et ouverture au monde » (coopération Cameroun-Chine, apprentissage des langues, débouchés professionnels) ;
\item \textbf{Rédiger} un paragraphe d'opinion argumenté (5--6 phrases) en caractères simplifiés corrects, avec ponctuation chinoise ;
\item \textbf{Participer} à un débat oral structuré en respectant la prononciation des tons (notamment \zh{yíyàng}, \zh{rúguǒ}, \zh{kěshì}).
\end{itemize}

\courssub{Situation déclenchante : La Semaine des langues}

Dans le cadre de la « Semaine des langues » du lycée, le club de chinois organise un débat ouvert aux classes de Première. Le thème : \emph{« Apprendre le chinois au Cameroun : atouts et défis »}. Le message ci-dessous est publié sur le groupe de discussion du club pour lancer la réflexion.

\begin{textebox}[Message du club de chinois — Groupe de discussion]
同学们好！下星期我们俱乐部有一个讨论会：「在喀麦隆学汉语：好处和困难」。\\
Tóngxuémen hǎo! Xià ge xīngqī wǒmen jùlèbù yǒu yí ge tǎolùnhuì: « Zài Kāmàilóng xué Hànyǔ: hǎochu hé kùnnan ».\\
« Chers camarades, bonjour ! La semaine prochaine, notre club organise une discussion : "Apprendre le chinois au Cameroun : avantages et difficultés". »\\[0.3em]
一方面，汉语跟英语不一样，汉字很难；另一方面，中国跟喀麦隆的合作越来越多，会说汉语很有用。\\
Yīfāngmiàn, Hànyǔ gēn Yīngyǔ bù yíyàng, Hànzì hěn nán; lìngyī fāngmiàn, Zhōngguó gēn Kāmàilóng de hézuò yuèláiyuè duō, huì shuō Hànyǔ hěn yǒuyòng.\\
« D'un côté, le chinois diffère de l'anglais, les caractères sont difficiles ; de l'autre, la coopération entre la Chine et le Cameroun augmente, parler chinois est très utile. »\\[0.3em]
如果你想说一说你的看法，请准备三个句子。欢迎大家参加！\\
Rúguǒ nǐ xiǎng shuō yi shuō nǐ de kànfǎ, qǐng zhǔnbèi sān ge jùzi. Huānyíng dàjiā cānjiā!\\
« Si tu veux exprimer ton avis, prépare trois phrases. Bienvenue à tous ! »
\end{textebox}

\courssub{Étape 1 — Observation et analyse guidée (10 min, individuel)}

Lisez le message et répondez par écrit :
\begin{enumerate}
\item Quel est le thème exact du débat ? Où et quand a-t-il lieu ?
\item Identifiez la structure \zh{一方面……,另一方面……}. Copiez la phrase chinoise, son pinyin et sa traduction.
\item Identifiez la comparaison avec \zh{跟……不一样}. Quels sont les deux termes comparés (A et B) ? Y a-t-il un adjectif après \zh{不一样} ?
\item Identifiez la structure \zh{如果……,……}. Quel est le verbe principal de la proposition \zh{如果} ? Quel mot-outil introduit la proposition principale (ici implicite, mais attendue) ?
\end{enumerate}

\begin{corrige}[Étape 1 — Réponses attendues]
\begin{enumerate}
\item Thème : « Apprendre le chinois au Cameroun : avantages et difficultés » (在喀麦隆学汉语：好处和困难). Lieu : au club de chinois (俱乐部). Moment : la semaine prochaine (下星期).
\item \zh{一方面，汉语跟英语不一样，汉字很难；另一方面，中国跟喀麦隆的合作越来越多，会说汉语很有用。} (Yīfāngmiàn, Hànyǔ gēn Yīngyǔ bù yíyàng, Hànzì hěn nán; lìngyī fāngmiàn, Zhōngguó gēn Kāmàilóng de hézuò yuèláiyuè duō, huì shuō Hànyǔ hěn yǒuyòng.) — « D'un côté, le chinois n'est pas comme l'anglais, les caractères sont difficiles ; de l'autre, la coopération Chine-Cameroun augmente, parler chinois est utile. »
\item Comparaison : A = \zh{汉语} (chinois), B = \zh{英语} (anglais). Structure : \zh{汉语跟英语不一样}. Pas d'adjectif après \zh{不一样} ici : la différence est globale (système d'écriture, grammaire, phonologie).
\item \zh{如果你想说一说你的看法} (Si tu veux dire ton avis). Verbe : \zh{想} (vouloir) + \zh{说一说} (dire un peu). La proposition principale commence par \zh{请} (imperatif/politesse), pas par \zh{但是/可是}. C'est une condition simple, pas une concession.
\end{enumerate}
\end{corrige}

\courssub{Étape 2 — Étude systématique des trois structures}

\begin{propriete}[Structure 1 : Comparaison d'égalité / différence — \zh{A 跟/和 B 一样 / 不一样 (+ 形容词)}]
\textbf{Schéma :} \zh{Sujet A + 跟/和 + Sujet B + (不)一样 (+ Adjectif)} \\[0.3em]
\textbf{Emplois :}
\begin{itemize}
\item \zh{一样} (yíyàng) : « être identique / pareil / semblable ». \textbf{L'adjectif (le cas échéant) se place APRÈS \zh{一样}}.
\item \zh{不一样} (bù yíyàng) : « être différent / ne pas être pareil ». Avec adjectif : « ne pas être aussi [Adj] que » (différence de degré).
\item \zh{跟} (gēn) et \zh{和} (hé) sont interchangeables (prépositions « avec / et »). \zh{跟} plus oral, \zh{和} plus écrit.
\end{itemize}
\textbf{Exemples :}
\begin{center}
\begin{tabularx}{\linewidth}{|X|X|X|}
\hline
\rowcolor{popblueL} Chinois (Caractères + Pinyin) & Structure & Traduction \\
\hline
\zh{雅温得跟北京一样热。} (Yǎwēndé gēn Běijīng yíyàng rè.) & A 跟 B 一样 Adj & Yaoundé est aussi chaud que Pékin. \\
\hline
\zh{汉语跟法语不一样难。} (Hànyǔ gēn Fǎyǔ bù yíyàng nán.) & A 跟 B 不一样 Adj & Le chinois n'est pas aussi difficile que le français. \\
\hline
\zh{我的书跟你的不一样。} (Wǒ de shū gēn nǐ de bù yíyàng.) & A 跟 B 不一样 (sans Adj) & Mon livre est différent du tien. \\
\hline
\end{tabularx}
\end{center}
\textbf{Attention :} \zh{不一样} nie l'égalité, pas l'adjectif. \zh{不一样难} = « pas aussi difficile que » (différence de degré), différent de \zh{不难} (pas difficile).
\end{propriete}

\begin{propriete}[Structure 2 : Deux aspects contrastés — \zh{一方面……, 另一方面……}]
\textbf{Schéma :} \zh{一方面，[Proposition 1]；另一方面，[Proposition 2] 。} \\[0.3em]
\textbf{Emplois :} Présenter deux facettes d'une même réalité (souvent : avantage / inconvénient ; cause / effet ; aspect positif / aspect négatif). Les deux volets sont symétriques.
\begin{itemize}
\item \zh{一方面} (yīfāngmiàn) = « D'un côté / En premier lieu ».
\item \zh{另一方面} (lìngyī fāngmiàn) = « De l'autre côté / En second lieu ».
\item Ponctuation : point-virgule chinois \zh{；} entre les deux volets, virgule après chaque marqueur.
\end{itemize}
\textbf{Exemples :}
\begin{center}
\begin{tabularx}{\linewidth}{|X|X|}
\hline
\rowcolor{popblueL} Chinois & Traduction \\
\hline
\zh{一方面汉字难写，另一方面语法不难。} (Yīfāngmiàn Hànzì nán xiě, lìngyī fāngmiàn yǔfǎ bù nán.) & D'un côté les caractères sont durs à écrire, de l'autre la grammaire n'est pas difficile. \\
\hline
\zh{学汉语一方面要花时间，另一方面很有用。} (Xué Hànyǔ yīfāngmiàn yào huā shíjiān, lìngyī fāngmiàn hěn yǒuyòng.) & Apprendre le chinois d'un côté prend du temps, de l'autre c'est très utile. \\
\hline
\end{tabularx}
\end{center}
\textbf{Variante :} \zh{一方面……，但/可是另一方面……} pour renforcer l'opposition.
\end{propriete}

\begin{propriete}[Structure 3 : Hypothèse conditionnelle avec restriction — \zh{如果……, 但是/可是……}]
\textbf{Schéma :} \zh{如果 + [Condition/Hypothèse] ， 但是/可是 + [Réalité restrictive / Conséquence nuancée] 。} \\[0.3em]
\textbf{Emplois :} \zh{如果} (rúguǒ) introduit une condition ou une supposition. \zh{但是} (dànshì) ou \zh{可是} (kěshì) introduit la proposition principale qui exprime un obstacle, une concession ou un résultat contraire à l'attente.
\begin{itemize}
\item \zh{如果} se place \textbf{au début} de la proposition subordonnée (sujet après \zh{如果}).
\item \zh{但是/可是} se place \textbf{au début} de la proposition principale (souvent après le sujet de la principale).
\item Souvent suivi de \zh{就} (jiù) dans la principale pour marquer la conséquence logique : \zh{如果……，可是……，就……}.
\end{itemize}
\textbf{Exemples :}
\begin{center}
\begin{tabularx}{\linewidth}{|X|X|}
\hline
\rowcolor{popblueL} Chinois & Traduction \\
\hline
\zh{如果下雨，可是我没带伞，就湿了。} (Rúguǒ xiàyǔ, kěshì wǒ méi dài sǎn, jiù shī le.) & S'il pleut, mais je n'ai pas pris de parapluie, [donc] je suis mouillé. \\
\hline
\zh{如果我想去中国，但是钱不够，我就得打工。} (Rúguǒ wǒ xiǎng qù Zhōngguó, dànshì qián bù gòu, wǒ jiù děi dǎgōng.) & Si je veux aller en Chine, mais l'argent ne suffit pas, je dois travailler. \\
\hline
\end{tabularx}
\end{center}
\textbf{Distinction :} \zh{虽然……但是/可是……} = Concession (fait réel) vs \zh{如果……但是/可是……} = Condition + obstacle/réalité contrastée.
\end{propriete}

\courssub{Étape 3 — Comparatif français / chinois \& Pièges francophones}

\begin{attention}[Erreurs fréquentes des apprenants francophones]
\begin{enumerate}
\item \textbf{Ordre de l'adjectif dans la comparaison :}
\begin{itemize}
\item \textbf{Faux :} \zh{汉语跟法语难不一样。} / \zh{雅温得跟北京大不一样。}
\item \textbf{Vrai :} \zh{汉语跟法语不一样难。} / \zh{雅温得跟北京不一样大。}
\item \textit{Explication :} En français « pas aussi difficile \textbf{que} », le « que » appelle l'adjectif avant. En chinois, \zh{不一样} est le noyau, l'adjectif le spécifie \textbf{après}.
\end{itemize}
\item \textbf{Oubli de \zh{跟/和} :}
\begin{itemize}
\item \textbf{Faux :} \zh{汉语一样英语。} / \zh{我一样你。}
\item \textbf{Vrai :} \zh{汉语跟英语一样。} / \zh{我跟你一样。}
\item \textit{Explication :} \zh{一样} n'est pas une préposition. Il faut \zh{跟/和} pour marquer le terme de comparaison.
\end{itemize}
\item \textbf{Placement de \zh{但是/可是} avec \zh{如果} :}
\begin{itemize}
\item \textbf{Faux :} \zh{如果我有钱但是我不去。} (Mais dans la même proposition).
\item \textbf{Vrai :} \zh{如果我有钱，但是我不想去。} ( Virgule obligatoire, \zh{但是} en tête de principale).
\item \textit{Explication :} \zh{如果} ouvre une parenthèse hypothétique ; \zh{但是} la referme pour opposer la réalité.
\end{itemize}
\item \textbf{Tons \& Sandhi :}
\begin{itemize}
\item \zh{一} (yī) $\to$ \zh{yí} (2e ton) devant un 4e ton (\zh{样} yàng) : \zh{yíyàng}.
\item \zh{不} (bù) $\to$ \zh{bú} (2e ton) devant un 4e ton : \zh{bú yíyàng}.
\item \zh{如果} (rúguǒ) : 2e + 3e ton. \zh{可是} (kěshì) : 3e + 4e ton (souvent kěshi à l'oral rapide).
\end{itemize}
\end{enumerate}
\end{attention}

\begin{aretenir}[Récapitulatif des trois structures — Fiche de révision]
\begin{center}
\begin{tabularx}{\linewidth}{|X|X|X|}
\hline
\rowcolor{popblueL} Structure & Schéma clé & Fonction communicative \\
\hline
\zh{A 跟/和 B (不)一样 (+ Adj)} & Adj \textbf{après} \zh{一样} & Comparer, évaluer une différence de degré ou de nature \\
\hline
\zh{一方面……，另一方面……} & \zh{；} sépare les volets & Argumenter, nuancer, présenter pour/contre \\
\hline
\zh{如果……，但是/可是……} & \zh{如果} (début sub.) / \zh{但是} (début princ.) & Formuler une hypothèse prudente, exprimer un obstacle \\
\hline
\end{tabularx}
\end{center}
\end{aretenir}

\courssub{Étape 4 — Lexique du thème : « Citoyenneté et ouverture au monde »}

\begin{center}
\begin{tabularx}{\linewidth}{|X|X|X|X|}
\hline
\rowcolor{popblueL} Caractères & Pinyin & Nature & Traduction \\
\hline
好处 & hǎochu & nom & avantage, point fort, bénéfice \\
\hline
困难 & kùnnan & nom / adj. & difficulté ; difficile \\
\hline
看法 & kànfǎ & nom & opinion, point de vue, avis \\
\hline
合作 & hézuò & nom / verbe & coopération ; coopérer \\
\hline
有用 & yǒuyòng & adj. & utile \\
\hline
有意思 & yǒu yìsi & adj. & intéressant, amusant \\
\hline
越来越 & yuèláiyuè & adv. & de plus en plus \\
\hline
机会 & jīhuì & nom & occasion, opportunité \\
\hline
工作 & gōngzuò & nom / verbe & travail, emploi ; travailler \\
\hline
参加 & cānjiā & verbe & participer à, prendre part à \\
\hline
努力 & nǔlì & verbe / adj. & faire des efforts ; assidu, courageux \\
\hline
语法 & yǔfǎ & nom & grammaire \\
\hline
汉字 & Hànzì & nom & caractère chinois \\
\hline
法语 & Fǎyǔ & nom & langue française \\
\hline
英语 & Yīngyǔ & nom & langue anglaise \\
\hline
\end{tabularx}
\end{center}

\courssub{Étape 5 — Entraînement guidé (Exercices d'application)}

\begin{exercice}[Exercice 1 : Transformation — Comparaison]
Transformez les phrases en utilisant \zh{跟/和……(不)一样(+ 形容词)}.
\begin{enumerate}
\item Le chinois est aussi utile que l'anglais. (utile = \zh{有用})
\item Yaoundé n'est pas aussi grande que Pékin. (grand = \zh{大})
\item Mon avis est différent du tien. (avis = \zh{看法})
\item La grammaire chinoise n'est pas aussi difficile que les caractères. (difficile = \zh{难})
\end{enumerate}
\end{exercice}


\begin{corrige}[Exercice 1]
\begin{enumerate}
\item \zh{汉语跟英语一样有用。} (Hànyǔ gēn Yīngyǔ yíyàng yǒuyòng.)
\item \zh{雅温得跟北京不一样大。} (Yǎwēndé gēn Běijīng bù yíyàng dà.)
\item \zh{我的看法跟你的不一样。} (Wǒ de kànfǎ gēn nǐ de bù yíyàng.)
\item \zh{汉语语法跟汉字不一样难。} (Hànyǔ yǔfǎ gēn Hànzì bù yíyàng nán.)
\end{enumerate}
\end{corrige}

\begin{exercice}[Exercice 2 : Complétion — 一方面……另一方面……]
Complétez les phrases pour exprimer un pour / contre.
\begin{enumerate}
\item Apprendre le chinois : \zh{一方面\_\_\_\_\_\_\_\_\_\_；另一方面\_\_\_\_\_\_\_\_\_\_。}
\item Travailler avec des Chinois : \zh{一方面\_\_\_\_\_\_\_\_\_\_；另一方面\_\_\_\_\_\_\_\_\_\_。}
\item Étudier en Chine : \zh{一方面\_\_\_\_\_\_\_\_\_\_；另一方面\_\_\_\_\_\_\_\_\_\_。}
\end{enumerate}
\end{exercice}


\begin{corrige}[Exercice 2 — Propositions de réponses]
\begin{enumerate}
\item \zh{一方面汉字很难写；另一方面语法比较简单。} / \zh{一方面要花很多时间；另一方面很有用。}
\item \zh{一方面有很多机会；另一方面文化差异大。} (文化差异 \zh{wénhuà chāyì} = différences culturelles).
\item \zh{一方面环境好；另一方面生活费贵。} (环境 \zh{huánjìng} = environnement ; 生活费 \zh{shēnghuófèi} = frais de vie).
\end{enumerate}
\end{corrige}

\begin{exercice}[Exercice 3 : Traduction — 如果……但是/可是……]
Traduisez en chinois (caractères + pinyin).
\begin{enumerate}
\item Si je vais en Chine, mais que je ne parle pas chinois, ce sera difficile.
\item Si tu veux réussir, mais que tu ne travailles pas, c'est impossible.
\item S'il pleut demain, mais que je n'ai pas de parapluie, je serai mouillé.
\end{enumerate}
\end{exercice}


\begin{corrige}[Exercice 3]
\begin{enumerate}
\item \zh{如果我去中国，但是我不会说汉语，就很难。} (Rúguǒ wǒ qù Zhōngguó, dànshì wǒ bù huì shuō Hànyǔ, jiù hěn nán.)
\item \zh{如果你想成功，但是你不努力，就不可能。} (Rúguǒ nǐ xiǎng chénggōng, dànshì nǐ bù nǔlì, jiù bù kěnéng.)
\item \zh{如果明天下雨，可是我没带伞，我就湿了。} (Rúguǒ míngtiān xiàyǔ, kěshì wǒ méi dài sǎn, wǒ jiù shī le.)
\end{enumerate}
\end{corrige}

\courssub{Étape 6 — Production intégrée : Le mini-débat \& Texte d'opinion}

\begin{methode}[Méthode : Rédiger un paragraphe d'opinion argumenté (5--6 phrases)]
Suivez ces 5 étapes pour passer de l'oral à l'écrit :
\begin{enumerate}
\item \textbf{Accroche (1 phrase) :} Donnez votre opinion générale avec \zh{我觉得 / 我认为 / 据说} + \zh{机会/挑战/好事}.
\item \textbf{Comparaison (1 phrase) :} Utilisez \zh{A 跟/和 B 不一样 (+ Adj)} pour situer le chinois par rapport à une langue connue (français, anglais, langue maternelle).
\item \textbf{Argument balancé (1 phrase) :} Utilisez \zh{一方面……，另一方面……} pour opposer coût (temps, effort) et gain (travail, culture, voyage).
\item \textbf{Hypothèse personnelle (1 phrase) :} Utilisez \zh{如果……，但是/可是……，就……} pour projeter votre avenir (voyage, étude, travail) avec un obstacle réaliste.
\item \textbf{Conclusion (1 phrase) :} Reprenez l'idée principale avec un connecteur (\zh{所以 / 因此 / 总之}) et éventuellement une concession \zh{虽然……可是……}.
\end{enumerate}
\textbf{Contrôle final :} Vérifiez l'ordre des mots (Adj après \zh{一样}), la ponctuation (\zh{，；。}), les classificateurs (\zh{一个机会}), les tons.
\end{methode}

\begin{experience}[Activité de classe : Mini-débat « Atouts et Défis »]
\textbf{Organisation :} Groupes de 4 élèves. Durée totale : 45 min.
\begin{enumerate}
\item \textbf{Préparation (15 min) :} Chaque élève écrit ses 3 phrases obligatoires (Comparaison, Pour/Contre, Si/Mais) sur une fiche. Le groupe s'entraide pour corriger les caractères et les tons.
\item \textbf{Présentation (20 min) :} Chaque groupe passe au tableau. Les 4 élèves enchaînent leurs 3 phrases (12 phrases/groupe). La classe note sur une \textbf{Grille d'écoute} :
\begin{center}
\begin{tabularx}{\linewidth}{|X|c|c|c|}
\hline
\rowcolor{popblueL} Structure entendue & Phrase notée (caractères/pinyin) & Correct ? (O/N) & Commentaire \\
\hline
\zh{A 跟/和 B 不一样 (+ Adj)} & & & \\
\hline
\zh{一方面……另一方面……} & & & \\
\hline
\zh{如果……但是/可是……} & & & \\
\hline
\end{tabularx}
\end{center}
\item \textbf{Échange (10 min) :} Questions de la classe (\zh{为什么？/ 举例？}). Vote pour le groupe le plus convaincant.
\end{enumerate}
\end{experience}

\begin{exoresolu}[Production modèle — Texte d'opinion pour affichage au club]
\textbf{Consigne :} Rédigez un paragraphe de 5 à 6 phrases sur le thème « Apprendre le chinois au Cameroun : atouts et défis » en intégrant les trois structures.
\tcblower
\textbf{Texte modèle :}\\[0.3em]
我觉得在喀麦隆学汉语是一个难得的机会。\\
Wǒ juéde zài Kāmàilóng xué Hànyǔ shì yí ge nándé de jīhuì.\\
« Je pense qu'apprendre le chinois au Cameroun est une opportunité rare. »\\[0.3em]
汉语跟法语不一样，汉字很难，但是语法不太难。\\
Hànyǔ gēn Fǎyǔ bù yíyàng, Hànzì hěn nán, dànshì yǔfǎ bù tài nán.\\
« Le chinois diffère du français : les caractères sont difficiles, mais la grammaire n'est pas trop difficile. »\\[0.3em]
一方面，学汉语要花很多时间和精力；另一方面，会说汉语找工作很容易。\\
Yīfāngmiàn, xué Hànyǔ yào huā hěn duō shíjiān hé jīnglì; lìngyī fāngmiàn, huì shuō Hànyǔ zhǎo gōngzuò hěn róngyì.\\
« D'un côté, apprendre le chinois demande beaucoup de temps et d'énergie ; de l'autre, parler chinois facilite la recherche d'emploi. »\\[0.3em]
如果我将来去中国工作，可是我的口语不好，我就必须现在多练习。\\
Rúguǒ wǒ jiānglái qù Zhōngguó gōngzuò, kěshì wǒ de kǒuyǔ bù hǎo, wǒ jiù bìxū xiànzài duō liànxí.\\
« Si plus tard je vais travailler en Chine, mais que mon oral n'est pas bon, je dois m'entraîner davantage maintenant. »\\[0.3em]
所以，虽然汉语跟英语不一样容易，可是它对未来非常有用。\\
Suǒyǐ, suīrán Hànyǔ gēn Yīngyǔ bù yíyàng róngyì, kěshì tā duì wèilái fēicháng yǒuyòng.\\
« Donc, bien que le chinois ne soit pas aussi facile que l'anglais, il est très utile pour l'avenir. »\\[0.5em]
\textbf{Analyse linguistique :}
\begin{itemize}
\item \textbf{Phrase 1} : \zh{难得} (nándé = rare, précieux) enrichit le vocabulaire. \zh{一个} (classificateur).
\item \textbf{Phrase 2} : \zh{不一样} sans adjectif (différence globale) + coordination \zh{但是} interne. Respect de l'ordre SVO.
\item \textbf{Phrase 3} : Structure \zh{一方面……另一方面……} canonique. \zh{时间和精力} (temps et énergie). \zh{找工作} (trouver du travail).
\item \textbf{Phrase 4} : \zh{如果……可是……就……} structure complète. \zh{将来} (futur). \zh{口语} (langage oral). \zh{必须} (devoir fort).
\item \textbf{Phrase 5} : Conclusion synthétique. Utilisation bonus de \zh{虽然……可是……} (concession réelle) + réemploi de \zh{不一样容易} (comparaison avec adjectif). Cohésion par \zh{所以}.
\end{itemize}
\end{exoresolu}

\courssub{Évaluation \& Grille de correction (5 points)}

\begin{center}
\begin{tabularx}{\linewidth}{|X|c|X|X|}
\hline
\rowcolor{popblueL} Critère & Points & Indicateurs de réussite \\
\hline
Exactitude des 3 structures (ordre, mots-outils, ponctuation) & 3 & \zh{不一样+Adj} correct ; \zh{一方面/另一方面} symétriques ; \zh{如果/但是} placements corrects (début de prop.) \\
\hline
Richesse lexicale & variée (vocabulaire leçon + personnel) & 1 & \zh{机会/困难/合作/有用/努力/口语/将来}... pas de répétition excessive \\
\hline
Forme (Caractères lisibles, traits corrects / Prononciation tons justes) & 1 & Écrit : ordre des traits respecté, pas de caractères inventés. Oral : \zh{yíyàng, rúguǒ, kěshì, bù yíyàng} distincts \\
\hline
\end{tabularx}
\end{center}

\courssub{Élément socioculturel : Cameroun — Chine}

\begin{savaistu}[La coopération linguistique Cameroun-Chine]
\begin{itemize}
\item \textbf{Instituts Confucius :} Deux au Cameroun (Université de Yaoundé II depuis 2007, Université de Douala depuis 2018). Cours grand public, HSK, formation professeurs.
\item \textbf{Bourses :} Chaque année, le gouvernement chinois offre des bourses d'études complètes (licence, master, doctorat, séjours linguistiques) aux étudiants camerounais via l'Ambassade de Chine à Yaoundé.
\item \textbf{Économie :} La Chine est le 1er partenaire commercial du Cameroun (projets infrastructures : barrage de Mékinac, stade d'Olembé, routes, port de Kribi). La maîtrise du chinois est un atout majeur sur le marché de l'emploi (BTP, commerce, interprariat, hôtellerie).
\item \textbf{Échanges culturels :} Fête du Printemps (Nouvel An chinois) célébrée à Yaoundé et Douala ; Semaine de la culture camerounaise en Chine. Jumelages villes : Yaoundé -- Pékin (accord 2015), Douala -- Shanghai.
\end{itemize}
\textbf{Prolongement :} Pour votre paragraphe final, ajoutez une comparaison locale : \zh{雅温得跟北京不一样大，可是发展都很快。} (Yaoundé n'est pas aussi grande que Pékin, mais toutes deux se développent vite).
\end{savaistu}

\begin{popfigure}
\begin{center}\tcbox[colback=popink,colframe=popink,coltext=white,arc=9pt,boxsep=4pt,left=6pt,right=6pt]{\bfseries\small Leçon 20 : Argumenter}\end{center}
\vspace{-2pt}
\begin{tcbraster}[raster columns=2,raster equal height=rows,raster column skip=6pt,raster row skip=6pt,enhanced,blankest]
\begin{tcolorbox}[enhanced,colback=white,colframe=popblue,boxrule=0.9pt,arc=3pt,title={Comparer \zh{跟/和...不一样}},fonttitle=\bfseries\scriptsize,coltitle=white,colbacktitle=popblue,left=4pt,right=4pt,top=2pt,bottom=2pt,boxsep=1pt,toptitle=1.5pt,bottomtitle=1.5pt,halign=left]
\begin{itemize}[nosep,leftmargin=*,itemsep=1pt,font=\scriptsize]
\item Égalité \zh{一样}
\item Différence \zh{不一样}
\item + Adj : degré
\end{itemize}
\end{tcolorbox}
\begin{tcolorbox}[enhanced,colback=white,colframe=poporange,boxrule=0.9pt,arc=3pt,title={Nuancer \zh{一方面...另一方面}},fonttitle=\bfseries\scriptsize,coltitle=white,colbacktitle=poporange,left=4pt,right=4pt,top=2pt,bottom=2pt,boxsep=1pt,toptitle=1.5pt,bottomtitle=1.5pt,halign=left]
\begin{itemize}[nosep,leftmargin=*,itemsep=1pt,font=\scriptsize]
\item Avantage
\item Inconvénient
\item Symétrie
\end{itemize}
\end{tcolorbox}
\begin{tcolorbox}[enhanced,colback=white,colframe=popgreen,boxrule=0.9pt,arc=3pt,title={Hypothétiser \zh{如果...但是/可是}},fonttitle=\bfseries\scriptsize,coltitle=white,colbacktitle=popgreen,left=4pt,right=4pt,top=2pt,bottom=2pt,boxsep=1pt,toptitle=1.5pt,bottomtitle=1.5pt,halign=left]
\begin{itemize}[nosep,leftmargin=*,itemsep=1pt,font=\scriptsize]
\item Condition \zh{如果}
\item Restriction \zh{但是/可是}
\item Conséquence \zh{就}
\end{itemize}
\end{tcolorbox}
\end{tcbraster}

\legende{Carte mentale des trois structures grammaticales de la leçon 20.}
\end{popfigure}

\courssub{Bilan \& Devoirs}

\begin{aretenir}[Bilan de la leçon 20]
\begin{itemize}
\item \textbf{Comparer :} \zh{A 跟/和 B (不)一样 (+ Adj)}. L'adjectif suit \zh{一样}. \zh{跟/和} obligatoires.
\item \textbf{Nuancer :} \zh{一方面……，另一方面……}. Structure binaire équilibrée, idéale pour « Pour / Contre ».
\item \textbf{Conditionner :} \zh{如果……，但是/可是……}. \zh{如果} en tête de subordonnée ; \zh{但是/可是} en tête de principale. Souvent + \zh{就}.
\item \textbf{Contexte :} Ces structures sont indispensables pour l'expression de l'opinion (épreuve orale/écrite du Bac, HSK 3-4, vie professionnelle).
\end{itemize}
\end{aretenir}

\begin{exercice}[Devoirs — Production écrite notée]
Rédigez un texte de 80--100 caractères chinois (environ 8--10 phrases) sur l'un des deux sujets au choix. Utilisez \textbf{au moins une fois chacune} des trois structures de la leçon. Soignez l'écriture (ordre des traits) et la ponctuation.
\begin{enumerate}
\item \textbf{Sujet A :} « Étudier en Chine vs Étudier au Cameroun » (比较：在中国留学 跟 在喀麦隆留学).
\item \textbf{Sujet B :} « Mon projet professionnel et le chinois » (如果……但是…… ; 一方面……另一方面……).
\end{enumerate}
\textbf{Critères :} Respect des schémas (3 pts) + Vocabulaire varié (2 pts) + Caractères \& Ponctuation (2 pts) + Cohérence globale (3 pts) = /10.
\end{exercice}

\newpage

\coursec{Méthodes et exercices résolus}

\coursec{Leçon 20 — Grammaire : Structures de comparaison, de nuance et de concession}

\begin{textebox}[Dialogue d'observation : Comparer, nuancer, décider]
\textbf{Contexte :} Armand (lycéen à Yaoundé) et Li Wei (étudiant chinois à Pékin) discutent en visio de leurs vies scolaires et de leurs projets.

\zh{阿曼：李伟，在喀麦隆学中文跟在中国不一样。} \\
\zh{李伟：哦？一方面老师很好，不严厉；另一方面机会少，难练口语吧？} \\
\zh{阿曼：没错。如果你想来雅温得，但是天气很热，要多喝水。} \\
\zh{李伟：如果下雨，但是我还是想去看足球赛。} \\
\zh{阿曼：好，我们用微信联系。一方面方便，另一方面别看太久，伤视力。}

\textbf{Pinyin :} \\
Āmàn: Lǐ Wěi, zài Kāmàilóng xué Zhōngwén gēn zài Zhōngguó bù yīyàng. \\
Lǐ Wěi: Ó? Yīfāngmiàn lǎoshī hěn hǎo, bù yánlì; lìngyī fāngmiàn jīhuì shǎo, nán liàn kǒuyǔ ba? \\
Āmàn: Méicuò. Rúguǒ nǐ xiǎng lái Yǎwēndé, dànshì tiānqì hěn rè, yào duō hē shuǐ. \\
Lǐ Wěi: Rúguǒ xiàyǔ, dànshì wǒ háishì xiǎng qù kàn zúqiú sài. \\
Āmàn: Hǎo, wǒmen yòng Wēixìn liánxì. Yīfāngmiàn fāngbiàn, lìngyī fāngmiàn bié kàn tài jiǔ, shāng shìlì.

\textbf{Traduction :} \\
Armand : Li Wei, étudier le chinois au Cameroun est différent de [le faire] en Chine. \\
Li Wei : Ah ? D'un côté les profs sont bien, pas stricts ; de l'autre les occasions sont rares, difficile de pratiquer l'oral, non ? \\
Armand : Exact. Si tu veux venir à Yaoundé, mais il fait très chaud, il faut boire beaucoup d'eau. \\
Li Wei : S'il pleut, mais j'irai quand même voir le match de foot. \\
Armand : D'accord, on se contacte sur WeChat. D'un côté c'est pratique, de l'autre ne regarde pas trop longtemps, ça abîme la vue.
\end{textebox}

\courssub{Structure 1 : La comparaison avec \cle{跟/和……(不)一样} [+ Adjectif]}

\begin{propriete}[Schéma et emplois de \cle{跟/和……(不)一样}]
\begin{center}
\begin{tabularx}{\linewidth}{|X|X|X|}
\hline
\rowcolor{popblueL} \textbf{Structure} & \textbf{Exemple (Caractères + Pinyin)} & \textbf{Traduction} \\
\hline
A \cle{跟/和} B \cle{一样} & \zh{我的手机 跟 你的手机 一样。} (Wǒ de shǒujī gēn nǐ de shǒujī yīyàng.) & Mon téléphone est pareil au tien. \\
\hline
A \cle{跟/和} B \cle{不一样} & \zh{喀麦隆的气候 和 北京的气候 不一样。} (Kāmàilóng de qìhòu hé Běijīng de qìhòu bù yīyàng.) & Le climat du Cameroun est différent de celui de Pékin. \\
\hline
A \cle{跟/和} B \cle{不一样} + \textbf{Adj} & \zh{中文作业 跟 法文作业 不一样 难。} (Zhōngwén zuòyè gēn Fǎwén zuòyè bù yīyàng nán.) & Les devoirs de chinois sont différemment difficiles (diffèrent par la difficulté) de ceux de français. \\
\hline
A \cle{跟/和} B \cle{一样} + \textbf{Adj} & \zh{他 跟 哥哥 一样 高。} (Tā gēn gēge yīyàng gāo.) & Il est aussi grand que son grand frère. \\
\hline
\end{tabularx}
\end{center}

\textbf{Règles clés :}
\begin{itemize}
    \item \cle{跟} (gēn) : oral, courant au Nord ; \cle{和} (hé) : écrit, formel, courant au Sud. Interchangeables ici.
    \item Ordre immuable : \textbf{Sujet A + 跟/和 + Objet B + (不)一样 (+ Adj)}.
    \item \cle{不一样 + Adj} : l'adjectif suit \textbf{sans} \cle{的}. Ex: \cle{不一样贵} (différemment cher), \cle{不一样忙} (différemment occupé).
    \item \cle{一样 + Adj} (égalité) : nécessite la présence de \cle{跟/和 B} devant. On ne dit pas \emph{*一样大}.
\end{itemize}
\end{propriete}

\begin{methode}[Construire une phrase de comparaison avec \cle{跟/和……(不)一样}]
\textbf{Objectif :} Dire que deux choses sont semblables ou différentes, avec ou sans précision adjectivale.

\textbf{Étapes :}
\begin{enumerate}
    \item \textbf{Identifier les deux termes} (A = Sujet, B = Objet de comparaison introduit par \cle{跟} ou \cle{和}).
    \item \textbf{Choisir la polarité} : \cle{一样} (identique/semblable) ou \cle{不一样} (différent).
    \item \textbf{(Optionnel) Ajouter un adjectif} après \cle{(不)一样} pour qualifier la différence ou l'égalité (\cle{不一样难}, \cle{一样高}).
    \item \textbf{Vérifier l'ordre :} \cle{A + 跟/和 + B + (不)一样 (+ Adj)}. \textbf{Jamais} de \cle{的} entre \cle{(不)一样} et l'adjectif.
\end{enumerate}
\end{methode}

\begin{exoresolu}[Comparer la vie scolaire : Yaoundé vs Pékin]
\textbf{Énoncé :} Traduis en chinois (caractères, pinyin) en utilisant \cle{跟/和……(不)一样}.
\begin{enumerate}
    \item Le climat de Yaoundé est différent de celui de Pékin.
    \item La cantine de mon lycée est pareille à celle de ton école.
    \item Les devoirs de chinois sont différemment difficiles que ceux de français. (Indice : \cle{不一样难}).
\end{enumerate}

\tcblower
\textbf{Solution détaillée :}

\textbf{Phrase 1 :}
\begin{itemize}
    \item \textbf{Analyse :} A = \cle{雅温得的气候} (Climat de Yaoundé), B = \cle{北京的气候}. Différence.
    \item \textbf{Vocabulaire :} 雅温得 \cle{Yǎwēndé} (Yaoundé), 北京 \cle{Běijīng} (Pékin), 气候 \cle{qìhòu} (climat).
    \item \textbf{Construction :} \cle{雅温得的气候 跟 北京的气候 不一样。}
    \item \textbf{Pinyin :} Yǎwēndé de qìhòu gēn Běijīng de qìhòu bù yīyàng.
\end{itemize}

\textbf{Phrase 2 :}
\begin{itemize}
    \item \textbf{Analyse :} A = \cle{我学校的食堂}, B = \cle{你学校的食堂}. Similarité. \cle{和} préféré à l'écrit.
    \item \textbf{Vocabulaire :} 学校 \cle{xuéxiào} (école), 食堂 \cle{shítáng} (cantine/resto U).
    \item \textbf{Construction :} \cle{我学校的食堂 和 你学校的食堂 一样。}
    \item \textbf{Pinyin :} Wǒ xuéxiào de shítáng hé nǐ xuéxiào de shítáng yīyàng.
\end{itemize}

\textbf{Phrase 3 :}
\begin{itemize}
    \item \textbf{Analyse :} A = \cle{中文作业}, B = \cle{法文作业}. Différence + adjectif \cle{难} (difficile). Structure \cle{不一样+Adj}.
    \item \textbf{Vocabulaire :} 中文 \cle{Zhōngwén} (chinois), 法文 \cle{Fǎwén} (français), 作业 \cle{zuòyè} (devoirs), 难 \cle{nán} (difficile).
    \item \textbf{Construction :} \cle{中文作业 跟 法文作业 不一样难。}
    \item \textbf{Pinyin :} Zhōngwén zuòyè gēn Fǎwén zuòyè bù yīyàng nán.
    \item \textbf{Note :} Sens : « Les devoirs de chinois et de français diffèrent par leur niveau de difficulté » (ne précise pas lequel est plus dur).
\end{itemize}
\end{exoresolu}

\courssub{Structure 2 : L'argumentation binaire avec \cle{一方面……，另一方面……}}

\begin{propriete}[Schéma et emplois de \cle{一方面……，另一方面……}]
\begin{center}
\begin{tabularx}{\linewidth}{|X|X|X|}
\hline
\rowcolor{poppurpleL} \textbf{Structure} & \textbf{Exemple (Caractères + Pinyin)} & \textbf{Traduction} \\
\hline
\cle{一方面}，[Sujet + Prédicat 1]；\cle{另一方面}，[Sujet + Prédicat 2] & \zh{一方面，微信很方便；另一方面，伤视力。} (Yīfāngmiàn, Wēixìn hěn fāngbiàn; lìngyī fāngmiàn, shāng shìlì.) & D'un côté WeChat est pratique ; de l'autre, ça abîme la vue. \\
\hline
\cle{一方面}，[Sujet + Prédicat 1]，\cle{而且}…；\cle{另一方面}，[Sujet + Prédicat 2] & \zh{一方面，老师好，而且不严；另一方面，机会少。} (Yīfāngmiàn, lǎoshī hǎo, érqiě bù yán; lìngyī fāngmiàn, jīhuì shǎo.) & D'un côté les profs sont bien et pas stricts ; de l'autre, peu d'occasions. \\
\hline
\end{tabularx}
\end{center}

\textbf{Règles clés :}
\begin{itemize}
    \item Couple indissociable : \cle{一方面} (yīfāngmiàn) / \cle{另一方面} (lìngyī fāngmiàn). Ne pas confondre avec \cle{一面……一面……} (actions simultanées : « tout en… »).
    \item Placement : En \textbf{thème} (début de proposition), suivies d'une virgule.
    \item Sujet : Souvent identique (omis dans la 2e partie) ou différent mais lié au même thème global.
    \item Ponctuation : Virgule après chaque marqueur ; point-virgule \textbf{；} pour séparer les deux grands segments si longs.
    \item Valeur : Contraste (avantages/inconvénients), concession, points de vue opposés.
\end{itemize}
\end{propriete}

\begin{methode}[Structurer un discours nuancé avec \cle{一方面……，另一方面……}]
\textbf{Étapes :}
\begin{enumerate}
    \item Identifier le \textbf{sujet unique} ou le thème global.
    \item Formuler l'aspect 1 (souvent positif / cause / pour) après \cle{一方面}，.
    \item Formuler l'aspect 2 (souvent négatif / conséquence / contre) après \cle{另一方面}，.
    \item Vérifier la cohérence : les deux propositions éclairent le même sujet.
    \item Ne \textbf{pas} ajouter \cle{但是} / \cle{可是} après \cle{另一方面} (redondant, la structure porte déjà le contraste).
\end{enumerate}
\end{methode}

\begin{exoresolu}[Analyser les réseaux sociaux : WeChat au Cameroun]
\textbf{Énoncé :} Rédige un court paragraphe (3-4 phrases) sur WeChat pour lycéens camerounais. Utilise \cle{一方面……，另一方面……}. Vocabulaire : \cle{方便} (pratique), \cle{沟通} (communiquer), \cle{影响} (affecter), \cle{学习} (études), \cle{视力} (vue).

\tcblower
\textbf{Solution détaillée :}

\textbf{Texte complet :}
\zh{微信对喀麦隆高中生很重要。一方面，微信很方便，可以跟家人朋友沟通；另一方面，如果用太多，会影响学习和视力。}

\textbf{Pinyin :}
Wēixìn duì Kāmàilóng gāozhōngshēng hěn zhòngyào. Yīfāngmiàn, Wēixìn hěn fāngbiàn, kěyǐ gēn jiārén péngyou gōutōng; lìngyī fāngmiàn, rúguǒ yòng tài duō, huì yǐngxiǎng xuéxí hé shìlì.

\textbf{Analyse :}
\begin{itemize}
    \item Phrase 1 : Constat général (\cle{对……很重要}).
    \item Phrase 2 : Structure cible. Sujet \cle{微信} pour \cle{方便} ; sujet nul (générique) pour \cle{用太多} et \cle{影响}.
    \item Combinaison : \cle{如果} inséré dans la 2e partie pour nuancer (condition d'excès).
\end{itemize}
\end{exoresolu}

\courssub{Structure 3 : La concession réelle avec \cle{如果……，但是/可是……}}

\begin{propriete}[Schéma et emplois de \cle{如果……，但是/可是……}]
\begin{center}
\begin{tabularx}{\linewidth}{|X|X|X|}
\hline
\rowcolor{poporangeL} \textbf{Structure} & \textbf{Exemple (Caractères + Pinyin)} & \textbf{Traduction} \\
\hline
\cle{如果} + [S1 + P1]，\cle{但是} + [S2 + P2] & \zh{如果下雨，但是我们还是去看球赛。} (Rúguǒ xiàyǔ, dànshì wǒmen háishì qù kàn qiúsài.) & S'il pleut, nous irons quand même voir le match. \\
\hline
\cle{如果} + [S1 + P1]，\cle{可是} + [S2 + P2] & \zh{如果他很累，可是他还坚持复习。} (Rúguǒ tā hěn lèi, kěshì tā hái jiānchí fùxí.) & S'il est très fatigué, pourtant il persévère à réviser. \\
\hline
\cle{如果} + [S1 + P1]，\cle{但是} + [S2 + P2]，\cle{所以} + [Conséquence] & \zh{如果我有时间，但是我没钱，所以我不去。} (Rúguǒ wǒ yǒu shíjiān, dànshì wǒ méi qián, suǒyǐ wǒ bù qù.) & Si j'ai le temps, mais je n'ai pas d'argent, donc je n'y vais pas. \\
\hline
\end{tabularx}
\end{center}

\textbf{Règles clés :}
\begin{itemize}
    \item \cle{如果} (rúguǒ) : « si », introduit l'hypothèse/condition. Place : \textbf{avant le sujet} (ou avant adv. temps/lieu).
    \item \cle{但是} (dànshì) : « mais », formel, contraste fort. \cle{可是} (kěshì) : « mais / pourtant », oral, nuance affective/regret. Place : \textbf{avant le sujet de la clause principale}.
    \item \textbf{Interdiction :} Pas de \cle{了} (aspect accompli) dans la clause \cle{如果}. L'action n'est pas réalisée.
    \item Marqueurs de persistance dans la clause principale : \cle{还是} (háishì, « quand même »), \cle{也} (yě, « aussi »), \cle{还} (hái, « encore/encore plus »).
\end{itemize}
\end{propriete}

\begin{methode}[Exprimer un contraste réel avec \cle{如果……，但是/可是……}]
\textbf{Étapes :}
\begin{enumerate}
    \item Poser la condition/hypothèse : \cle{如果} + Sujet 1 + Verbe/Adj.
    \item Introduire le contraste/la réalité : \cle{但是} (écrit/formel) ou \cle{可是} (oral/émotion) + Sujet 2 + Verbe/Adj.
    \item Ajouter \cle{还是}/\cle{还}/\cle{也} avant le verbe principal pour renforcer le « malgré tout ».
    \item Vérifier : Pas de \cle{了} après le verbe de la clause \cle{如果}.
\end{enumerate}
\end{methode}

\begin{exoresolu}[Condition météo et décision : Match à Douala / Révision]
\textbf{Énoncé :} Complète en chinois (caractères + pinyin) avec \cle{如果……，但是/可是……}. Traduis.
\begin{enumerate}
    \item \_\_\_\_\_\_ 下雨 \_\_\_\_\_\_，\_\_\_\_\_\_ 我们 \_\_\_\_\_\_ 去 看 球赛。 (Si pluie $\rightarrow$ mais nous allons quand même).
    \item \_\_\_\_\_\_ 他 \_\_\_\_\_\_ 很 累 \_\_\_\_\_\_，\_\_\_\_\_\_ 他 \_\_\_\_\_\_ 坚持 复习 汉字。 (S'il fatigué $\rightarrow$ mais il persévère).
\end{enumerate}

\tcblower
\textbf{Solution détaillée :}

\textbf{Phrase 1 :} (Contraste fort, décision maintenue $\rightarrow$ \cle{但是} + \cle{还是})
\zh{如果下雨，但是我们还是去看球赛。} \\
Rúguǒ xiàyǔ, dànshì wǒmen \textbf{háishì} qù kàn qiúsài. \\
S'il pleut, mais nous irons \textbf{quand même} voir le match.

\textbf{Phrase 2 :} (État physique vs effort mental, nuanceorale $\rightarrow$ \cle{可是} + \cle{还})
\zh{如果他很累，可是他还坚持复习汉字。} \\
Rúguǒ tā hěn lèi, kěshì tā \textbf{hái} jiānchí fùxí hànzì. \\
S'il est très fatigué, \textbf{pourtant} il persiste encore à réviser les caractères.

\textbf{Piège évité :} \emph{*如果下雨了，我们去...} $\rightarrow$ \cle{了} incompatible avec \cle{如果} (hypothèse non réalisée).
\end{exoresolu}

\courssub{Synthèse et production : Combiner les trois structures}

\begin{methode}[Méthode « 3P » : Rédiger un paragraphe argumentatif cohérent (Niveau Bac LV2)]
\textbf{Objectif :} Produire un texte de 60-80 caractères intégrant : Comparaison (\cle{不一样}), Nuance (\cle{一方面/另一方面}), Condition contrastée (\cle{如果/但是}).

\textbf{Étapes :}
\begin{enumerate}
    \item \textbf{Planifier (5 min) :} Schématiser : Sujet $\rightarrow$ Comparaison factuelle $\rightarrow$ Analyse pour/contre $\rightarrow$ Conseil conditionnel.
    \item \textbf{Produire (10 min) :} Rédiger \textbf{directement en chinois} (pas de traduction mot à mot).
        \begin{itemize}
            \item Phrase 1 : \cle{A 跟 B 不一样} (Fait).
            \item Phrase 2 : \cle{一方面……；另一方面……} (Analyse).
            \item Phrase 3 : \cle{如果……，但是……} (Conseil/Prédiction).
        \end{itemize}
    \item \textbf{Polir (5 min) :} Vérifier : Ordre mots (Suj-\cle{跟}-Obj-\cle{不一样}), Ponctuation (， 。 ；), Tons pinyin, Caractères lisibles.
\end{enumerate}
\textbf{Mots-outils :} \cle{而且} (de plus), \cle{不过} (cependant - léger), \cle{所以} (donc), \cle{还是} (quand même).
\end{methode}

\begin{exoresolu}[Rédaction : « Étudier le chinois à Yaoundé vs Pékin »]
\textbf{Énoncé :} Écris un message à Li Ming (ami chinois). Contraintes :
\begin{itemize}
    \item Compare l'environnement d'apprentissage (Yaoundé vs Pékin) avec \cle{跟……不一样}.
    \item Avantage + Inconvénient d'étudier à Yaoundé avec \cle{一方面……，另一方面……}.
    \item Conseil conditionnel avec \cle{如果……，但是……}.
\end{itemize}

\tcblower
\textbf{Solution détaillée :}

\textbf{1. Planification :}
\begin{itemize}
    \item Comp : \cle{语言环境} (environnement linguistique) Yaoundé $\neq$ Pékin.
    \item Nuance : + Profs sympas (\cle{老师好，不严厉}) ; - Peu pratique oral (\cle{机会少，难练口语}).
    \item Condition : Si tu viens (\cle{来喀麦隆}), mais chaud (\cle{天气热}) $\rightarrow$ boire eau (\cle{多喝水}).
\end{itemize}

\textbf{2. Production (Texte chinois $\approx$ 70 caractères) :}
\zh{在喀麦隆学中文跟在中国不一样。一方面，老师很好，不严厉；另一方面，机会少，难练口语。如果你来喀麦隆，但是天气很热，要多喝水。}

\textbf{3. Annotation :}
\begin{itemize}
    \item \zh{在喀麦隆学中文跟在中国不一样。} (Zài Kāmàilóng xué Zhōngwén gēn zài Zhōngguó bù yīyàng.) Étudier le chinois au Cameroun diffère de la Chine.
    \item \zh{一方面，老师很好，不严厉；} (Yīfāngmiàn, lǎoshī hěn hǎo, bù yánlì;) D'un côté, profs bien, pas stricts ;
    \item \zh{另一方面，机会少，难练口语。} (lìngyī fāngmiàn, jīhuì shǎo, nán liàn kǒuyǔ.) de l'autre, occasions rares, dur de pratiquer l'oral.
    \item \zh{如果你来喀麦隆，但是天气很热，要多喝水。} (Rúguǒ nǐ lái Kāmàilóng, dànshì tiānqì hěn rè, yào duō hē shuǐ.) Si tu viens au Cameroun, mais il fait chaud, faut boire de l'eau.
\end{itemize}

\textbf{4. Vérification :} Ordre \cle{跟……不一样} OK. Couple \cle{一方面/另一方面} + \cle{；} OK. \cle{如果} tête clause 1, \cle{但是} tête clause 2 OK. \cle{要} (obligation) pour conseil OK.
\end{exoresolu}

\courssub{Atelier de transformation : Du français vers le chinois (Éviter les calques)}

\begin{methode}[Transformer : Comparatif français $\rightarrow$ Structure chinoise]
\textbf{Étapes systématiques :}
\begin{enumerate}
    \item Repérer le comparatif FR : \emph{« A est différent de B »}, \emph{« A ressemble à B »}, \emph{« A n'est pas comme B »}.
    \item Extraire A (Sujet) et B (Objet).
    \item Choisir \cle{跟} (oral) ou \cle{和} (écrit).
    \item Appliquer : \cle{A + 跟/和 + B + 不一样 / 一样}.
    \item Gérer l'adjectif :
        \begin{itemize}
            \item « Différemment + Adj » $\rightarrow$ \cle{不一样 + Adj} (sans \cle{的}).
            \item « Aussi + Adj que » $\rightarrow$ \cle{跟……一样 + Adj}.
        \end{itemize}
    \item Supprimer prépositions FR (\emph{de, à, que}) : elles n'existent pas dans cette structure.
\end{enumerate}
\end{methode}

\begin{exoresolu}[Transformation guidée : Pièges à éviter]
\textbf{Énoncé :} Transforme en chinois correct (Caractères + Pinyin). Identifie le calque.

\begin{center}
\begin{tabularx}{\linewidth}{|X|X|X|}
\hline
\rowcolor{popblueL} \textbf{Français (Source)} & \textbf{Chinois (Cible)} & \textbf{Analyse / Piège évité} \\
\hline
1. Mon téléphone est différent de ton téléphone. & \zh{我的手机 跟 你的手机 不一样。} (Wǒ de shǒujī gēn nǐ de shǒujī bù yīyàng.) & \textbf{Calque :} \emph{*我手机不同你手机} (\cle{不同} verbe statif, structure diffère) ou \emph{*不一样跟...} (ordre). \textbf{Règle :} \cle{跟/和} introduit B \textbf{avant} \cle{不一样}. \\
\hline
2. La nourriture camerounaise n'est pas comme la nourriture chinoise. & \zh{喀麦隆的菜 和 中国的菜 不一样。} (Kāmàilóng de cài hé Zhōngguó de cài bù yīyàng.) & \textbf{Nuance :} \cle{不像} (bù xiàng) = « ne ressemble pas à » (apparence). \cle{不一样} = « n'est pas pareil » (différence générale : goût, style). Programme impose \cle{不一样}. \\
\hline
3. Les prix à Douala sont différemment chers qu'à Yaoundé. & \zh{杜拉的价格 跟 雅温得的价格 不一样贵。} (Dùlā de jiàgé gēn Yǎwēndé de jiàgé bù yīyàng guì.) & \textbf{Structure clé :} \cle{不一样 + 贵}. \textbf{Pas de \cle{的}}. Sens : différence de degré de cherté (sans dire lequel est plus cher). \\
\hline
4. D'un côté je veux aller en Chine, de l'autre je n'ai pas d'argent. & \zh{一方面，我想去中国；另一方面，我没钱。} (Yīfāngmiàn, wǒ xiǎng qù Zhōngguó; lìngyī fāngmiàn, wǒ méi qián.) & \textbf{Piège :} Ne pas ajouter \cle{但是} ! \cle{一方面/另一方面} \textbf{suffisent} à marquer le contraste argumentatif. Sujet \cle{我} répété pour clarté. \\
\hline
5. Si j'ai le temps, mais je n'ai pas d'argent, je ne vais pas. & \zh{如果我有时间，但是我没钱，所以我不去。} (Rúguǒ wǒ yǒu shíjiān, dànshì wǒ méi qián, suǒyǐ wǒ bù qù.) & \textbf{Logique :} Condition (\cle{有时间}) + Contraste réel (\cle{没钱}) $\rightarrow$ Conséquence (\cle{所以不去}). \textbf{Crucial :} Pas de \cle{了} dans clause \cle{如果}. \cle{但是} en tête clause principale. \\
\hline
\end{tabularx}
\end{center}
\end{exoresolu}

\begin{savaistu}[Éducation et numérique : Cameroun vs Chine]
\textbf{Système éducatif :}
\begin{itemize}
    \item \textbf{Cameroun :} Système bilingue (français/anglais). Secondaire : 7 ans (4 ans 1er cycle + 3 ans 2e cycle). Bac à la fin de Terminale. Chinois = LV2/LV3 optionnel, en développement (Instituts Confucius à Yaoundé, Douala, Buea).
    \item \textbf{Chine :} 9 ans obligatoires (6 primaire + 3 collège). Secondaire : 3 ans (lycée \cle{高中}). \cle{高考} (Gāokǎo) examen national ultra-compétitif pour l'université. Anglais obligatoire dès le primaire ; japonais/coréen/français options rares.
\end{itemize}

\textbf{WeChat (\cle{微信 Wēixìn}) :}
\begin{itemize}
    \item \textbf{Chine :} « Super-app » indispensable (paiement \cle{微信支付}, admin, travail, social). 1,3+ milliard d'utilisateurs. Intégré à la vie quotidienne (marchés, taxis, hôpitaux).
    \item \textbf{Cameroun :} Utilisé par la diaspora chinoise, étudiants en chinois, importateurs (Douala, Yaoundé). Fonction paiement limitée (nécessite compte bancaire chinois). Principal usage : messagerie, groupes de classe, suivi cours en ligne.
\end{itemize}

\textbf{Point culturel :} « \cle{严厉} » (strict/sévère). En Chine, l'autorité professorale est traditionnellement forte (\cle{师道尊严}), devoirs nombreux, discipline ferme. Au Cameroun, respect dû au \cle{Maître/Professeur} mais relation souvent plus accessible, dialogue possible. \cle{不严厉} perçu comme avantage par élèves camerounais = bienveillance, non laxisme.
\end{savaistu}

\begin{attention}[Les 5 erreurs qui coûtent des points au Bac (Première A)]
\begin{enumerate}
    \item \textbf{Ordre des mots : \cle{跟/和} placé APRÈS \cle{不一样}.}
    \emph{Erreur :} \zh{*我 不一样 跟 你。} $\rightarrow$ \textbf{Correct :} \zh{我 跟 你 不一样。} (Sujet + \cle{跟/和} + Objet + \cle{不一样}). La préposition introduit le terme de comparaison \textbf{avant} le prédicat.
    
    \item \textbf{Insertion abusive de \cle{的} entre \cle{不一样} et l'adjectif.}
    \emph{Erreur :} \zh{*这道菜 跟 那道菜 不一样 的 辣。} $\rightarrow$ \textbf{Correct :} \zh{这道菜 跟 那道菜 不一样 辣。} \cle{不一样} modifie directement l'adjectif (fonction adverbiale), pas de \cle{的}.
    
    \item \textbf{Confusion \cle{一方面} / \cle{一面} et omission de \cle{另一方面}.}
    \emph{Erreur :} \zh{*一面 我 忙，一面 我 累。} (Signifie : « Je suis occupé \textbf{en même temps que} je suis fatigué » = actions simultanées). $\rightarrow$ \textbf{Correct (contraste) :} \zh{一方面 我 忙，另一方面 我 累。} (D'un côté je suis occupé, de l'autre je suis fatigué = arguments contraires). Couple indissociable.
    
    \item \textbf{Utilisation de \cle{了} dans la clause \cle{如果} (hypothèse).}
    \emph{Erreur :} \zh{*如果 下雨 了，我 不 去。} $\rightarrow$ \textbf{Correct :} \zh{如果 下雨，我 不 去。} \cle{如果} = condition non réalisée $\rightarrow$ pas d'aspect accompli \cle{了}. \cle{了} possible seulement clause principale si action réalisée *après* condition (ex: \cle{如果下雨，我就不去了}).
    
    \item \textbf{Calque « Si... mais... » $\rightarrow$ Double sujet mal géré / répétition verbe.}
    \emph{Erreur :} \zh{*如果 下雨，下雨 但是 我 去。} $\rightarrow$ \textbf{Correct :} \zh{如果 下雨，但是 我 还是 去。} Sujet clause principale (\cle{我}) suit \cle{但是}. Ne pas répéter sujet/verbe clause \cle{如果}. Utiliser \cle{还是} / \cle{也} / \cle{还} pour marquer la persistance.
\end{enumerate}
\end{attention}

\begin{center}
\begin{tabularx}{\linewidth}{|X|X|X|X|}
\hline
\rowcolor{popblueL} \textbf{Caractères} & \textbf{Pinyin} & \textbf{Nature} & \textbf{Traduction} \\
\hline
雅温得 & Yǎwēndé & n.p. & Yaoundé (capitale) \\
\hline
杜拉 & Dùlā & n.p. & Douala (capitale économique) \\
\hline
喀麦隆 & Kāmàilóng & n.p. & Cameroun \\
\hline
北京 & Běijīng & n.p. & Pékin \\
\hline
气候 & qìhòu & n. & climat \\
\hline
食堂 & shítáng & n. & cantine, restaurant scolaire \\
\hline
作业 & zuòyè & n. & devoirs \\
\hline
中文 / 汉语 & Zhōngwén / Hànyǔ & n. & chinois (langue) \\
\hline
法文 & Fǎwén & n. & français (langue) \\
\hline
难 & nán & adj. & difficile \\
\hline
严厉 & yánlì & adj. & strict, sévère \\
\hline
微信 & Wēixìn & n.p. & WeChat \\
\hline
方便 & fāngbiàn & adj. & pratique, commode \\
\hline
沟通 & gōutōng & v. & communiquer, échanger \\
\hline
影响 & yǐngxiǎng & v. & influencer, affecter (souvent négatif) \\
\hline
视力 & shìlì & n. & vue, acuité visuelle \\
\hline
球赛 & qiúsài & n. & match de ballon (foot, basket...) \\
\hline
坚持 & jiānchí & v. & persévérer, insister \\
\hline
复习 & fùxí & v. & réviser, réétudier \\
\hline
汉字 & hànzì & n. & caractère chinois \\
\hline
语言环境 & yǔyán huánjìng & n. & environnement linguistique \\
\hline
机会 & jīhuì & n. & occasion, opportunité \\
\hline
练 & liàn & v. & pratiquer, s'entraîner \\
\hline
口语 & kǒuyǔ & n. & langue orale, expression orale \\
\hline
天气 & tiānqì & n. & temps, météo \\
\hline
热 & rè & adj. & chaud (météo) \\
\hline
喝水 & hē shuǐ & v-o & boire de l'eau \\
\hline
价格 & jiàgé & n. & prix \\
\hline
钱 & qián & n. & argent \\
\hline
时间 & shíjiān & n. & temps \\
\hline
所以 & suǒyǐ & conj. & donc, par conséquent \\
\hline
还是 & háishì & adv. & encore, quand même (dans phrase affirmative) \\
\hline
还 & hái & adv. & encore, encore plus \\
\hline
\end{tabularx}
\end{center}

\begin{exercice}[Entraînement : Application et Transformation]
\textbf{Exercice 1 : Compléter avec \cle{跟/和……(不)一样} (+ adjectif si indiqué).}
\begin{enumerate}
    \item Le ndolé de ma mère \_\_\_\_\_\_ celui du restaurant \_\_\_\_\_\_ (bon / 好 hǎo). [Différent en étant bon ? Non, structure \cle{不一样+Adj} = différemment bon. Dire : « Le ndolé de ma mère est différemment bon (meilleur/goût différent) que celui du resto ».]
    \item Mon uniforme \_\_\_\_\_\_ le tien \_\_\_\_\_\_ (pareil).
    \item La vie à Buea \_\_\_\_\_\_ la vie à Douala \_\_\_\_\_\_ (cher / 贵 guì). [Différemment chère].
\end{enumerate}

\textbf{Exercice 2 : Relier les deux idées avec \cle{一方面……，另一方面……}.}
\begin{enumerate}
    \item Idée 1 : Étudier en Chine est fatigant (\cle{累 lèi}). Idée 2 : C'est très efficace (\cle{有效 yǒuxiào}).
    \item Idée 1 : Le foot est amusant (\cle{有意思 yǒu yìsi}). Idée 2 : Ça fatigue (\cle{累人 lèi rén}).
\end{enumerate}

\textbf{Exercice 3 : Compléter avec \cle{如果……，但是/可是……} (+ \cle{还是}/\cle{还}).}
\begin{enumerate}
    \item \_\_\_\_\_\_ 明天 \_\_\_\_\_\_ 下雨 \_\_\_\_\_\_，\_\_\_\_\_\_ 我 \_\_\_\_\_\_ \_\_\_\_\_\_ 去 学校。 (Si pluie demain $\rightarrow$ mais j'irai \textbf{quand même} à l'école).
    \item \_\_\_\_\_\_ 这道 菜 \_\_\_\_\_\_ 很 辣 \_\_\_\_\_\_，\_\_\_\_\_\_ 我 \_\_\_\_\_\_ \_\_\_\_\_\_ 想 吃。 (Si plat épicé $\rightarrow$ pourtant je \textbf{veux encore} manger).
\end{enumerate}

\textbf{Exercice 4 : Expression écrite (60-80 caractères).}
Sujet : « Tu conseilles à ton ami camerounais qui part étudier à Pékin. »
Contraintes :
\begin{itemize}
    \item Compare la météo (Yaoundé vs Pékin) avec \cle{不一样}.
    \item Donne un avantage (ex: métro pratique \cle{地铁方便 dìtiě fāngbiàn}) et un inconvénient (ex: hiver froid \cle{冬天冷 dōngtiān lěng}) de Pékin avec \cle{一方面……，另一方面……}.
    \item Donne un conseil vestimentaire conditionnel avec \cle{如果……，但是……} (ex: si hiver, mais chauffage inside \cle{有暖气 yǒu nuǎnqì}).
\end{itemize}
\end{exercice}

\begin{corrige}[Exercices Leçon 20]
\textbf{Exercice 1 :}
\begin{enumerate}
    \item \zh{我妈妈的炖菜 跟 餐厅的 不一样 好。} (Wǒ māma de dùncài gēn cāntīng de bù yīyàng hǎo.) — *Note : \cle{炖菜} (ragoût) pour ndolé.*
    \item \zh{我的校服 跟 你的 一样。} (Wǒ de xiàofú gēn nǐ de yīyàng.)
    \item \zh{布埃亚的生活 跟 杜拉的生活 不一样 贵。} (Bù'āiyà de shēnghuó gēn Dùlā de shēnghuó bù yīyàng guì.) — \cle{布埃亚} (Bù'āiyà) = Buea.
\end{enumerate}

\textbf{Exercice 2 :}
\begin{enumerate}
    \item \zh{一方面，在中国学习很累；另一方面，很有效。} (Yīfāngmiàn, zài Zhōngguó xuéxí hěn lèi; lìngyī fāngmiàn, hěn yǒuxiào.)
    \item \zh{一方面，踢足球很有意思；另一方面，很累人。} (Yīfāngmiàn, tī zúqiú hěn yǒu yìsi; lìngyī fāngmiàn, hěn lèi rén.)
\end{enumerate}

\textbf{Exercice 3 :}
\begin{enumerate}
    \item \zh{如果 明天 下雨，但是 我 还是 去 学校。} (Rúguǒ míngtiān xiàyǔ, dànshì wǒ \textbf{háishì} qù xuéxiào.)
    \item \zh{如果 这道 菜 很 辣，可是 我 还 想 吃。} (Rúguǒ zhè dào cài hěn là, kěshì wǒ \textbf{hái} xiǎng chī.)
\end{enumerate}

\textbf{Exercice 4 : Proposition de rédaction (≈ 75 caractères) :}
\zh{北京的天气跟雅温得不一样。一方面，地铁很方便；另一方面，冬天很冷。如果你冬天去，但是室内有暖气，要带件大衣。}
\textbf{Pinyin :} Běijīng de tiānqì gēn Yǎwēndé bù yīyàng. Yīfāngmiàn, dìtiě hěn fāngbiàn; lìngyī fāngmiàn, dōngtiān hěn lěng. Rúguǒ nǐ dōngtiān qù, dànshì shìnèi yǒu nuǎnqì, yào dài jiàn dàyī.
\textbf{Traduction :} La météo de Pékin diffère de Yaoundé. D'un côté le métro est pratique ; de l'autre l'hiver est très froid. Si tu y vas en hiver, mais il y a le chauffage à l'intérieur, apporte un manteau.
\end{corrige}

\newpage

\coursec{Exercices}

\courssub{Je vérifie mes connaissances}

\begin{exercice}[QCM : la bonne structure]
Pour chaque phrase, choisis la bonne réponse parmi A, B, C ou D.

1. 我的汉语水平\_\_\_你的\_\_\_一样。\\
A. 跟……不 \quad B. 跟……也 \quad C. 和……很 \quad D. 跟……（rien）

2. Quelle phrase est correcte ?\\
A. 我不跟他一样高。 \quad B. 我跟他不一样高。 \quad C. 我跟他一样不高。 \quad D. 我不一样跟他高。

3. \_\_\_我想去中国留学，\_\_\_我的钱不够。\\
A. 一方面……一方面 \quad B. 一方面……另一方面 \quad C. 如果……但是 \quad D. 因为……所以

4. \_\_\_你每天都学习，\_\_\_你不练习说话，你的口语就不会进步。\\
A. 因为……所以 \quad B. 如果……但是 \quad C. 一方面……另一方面 \quad D. 虽然……可是
\end{exercice}

\begin{exercice}[Vrai ou faux ? Justifie ta réponse]
Dis si les phrases suivantes sont correctes ou incorrectes, puis justifie en citant la règle de grammaire.

1. 雅温得跟北京不一样大。\\
\emph{Yǎwēndé gēn Běijīng bù yīyàng dà.}

2. 一方面我喜欢汉语，二方面我觉得汉字很难。

3. 我的弟弟和我一样高。\\
\emph{Wǒ de dìdi hé wǒ yīyàng gāo.}

4. 如果你学习，所以你成功。
\end{exercice}

\begin{exercice}[Vocabulaire : complète le tableau]
Complète le tableau suivant avec le caractère, le pinyin ou la traduction manquants.

\begin{center}
\begin{tabularx}{\linewidth}{|X|X|X|X|}
\hline
\rowcolor{popblueL} Caractères & Pinyin & Nature & Traduction \\
\hline
一样 & \ldots\ldots\ldots & adverbe/adjectif & \ldots\ldots\ldots \\
\hline
\ldots\ldots\ldots & bǐjiào & verbe/adverbe & \ldots\ldots\ldots \\
\hline
方面 & \ldots\ldots\ldots & nom & \ldots\ldots\ldots \\
\hline
\ldots\ldots\ldots & rúguǒ & conjonction & \ldots\ldots\ldots \\
\hline
可是 & \ldots\ldots\ldots & \ldots\ldots\ldots & mais \\
\hline
\ldots\ldots\ldots & shuǐpíng & nom & \ldots\ldots\ldots \\
\hline
\end{tabularx}
\end{center}
\end{exercice}

\begin{exercice}[Associe les caractères au pinyin]
Relie chaque mot en caractères à son pinyin.

\begin{center}
\begin{tabularx}{\linewidth}{|X|X|}
\hline
\rowcolor{popblueL} Caractères & Pinyin \\
\hline
1. 水平 & a. fāngmiàn \\
\hline
2. 比较 & b. liúxué \\
\hline
3. 方面 & c. shuǐpíng \\
\hline
4. 留学 & d. bǐjiào \\
\hline
5. 进步 & e. jìnbù \\
\hline
\end{tabularx}
\end{center}
\end{exercice}

\courssub{Je m'entraîne}

\begin{exercice}[Traduis en chinois]
Traduis les phrases suivantes en chinois (caractères et pinyin).

1. Mon frère est aussi grand que mon père.

2. D'une part j'aime le chinois, d'autre part je le trouve difficile.

3. Si tu étudies, mais que tu ne pratiques pas, tu oublieras vite.

4. Le ndolé et le riz chinois ne sont pas pareils.

5. Douala n'est pas aussi froide que Pékin en hiver.
\end{exercice}

\begin{exercice}[Transformation de phrases]
1. Réécris la phrase \zh{他比我高。} \emph{(Tā bǐ wǒ gāo.)} avec la structure \zh{……跟……（不）一样}, en adaptant le sens, puis justifie ton choix d'utiliser la forme affirmative ou négative.

2. Réécris la phrase \zh{法语和英语不一样难。} avec la structure \zh{……比……} en gardant le même sens.

3. Transforme les deux phrases \zh{我想去中国。我没有钱。} en une seule phrase avec \zh{一方面……，另一方面……}.
\end{exercice}

\begin{exercice}[Texte à trous]
Complète le texte suivant avec les mots-outils de la liste : \zh{跟} — \zh{一样} — \zh{一方面} — \zh{另一方面} — \zh{如果} — \zh{可是} — \zh{不}.

\begin{textebox}[Ma ville et Shanghai]
我的城市\_\_\_上海\_\_\_一样大。\_\_\_我的城市很安静，\_\_\_上海很热闹。\_\_\_你想学汉语，\_\_\_你不去中国，你也可以在喀麦隆学得很好。上海的生活\_\_\_我们这里的生活很\_\_\_一样吗？我觉得\_\_\_一样。\\
\emph{Wǒ de chéngshì \_\_\_ Shànghǎi \_\_\_ yīyàng dà. \_\_\_ wǒ de chéngshì hěn ānjìng, \_\_\_ Shànghǎi hěn rènao. \_\_\_ nǐ xiǎng xué Hànyǔ, \_\_\_ nǐ bú qù Zhōngguó, nǐ yě kěyǐ zài Kāmàilóng xué de hěn hǎo. Shànghǎi de shēnghuó \_\_\_ wǒmen zhèlǐ de shēnghuó hěn \_\_\_ yīyàng ma? Wǒ juéde \_\_\_ yīyàng.}\\
« Ma ville et Shanghai ne sont pas aussi grandes l'une que l'autre. D'une part ma ville est calme, d'autre part Shanghai est animée. Si tu veux apprendre le chinois, mais que tu ne vas pas en Chine, tu peux aussi bien l'apprendre au Cameroun. La vie à Shanghai est-elle pareille à la vie d'ici ? Je pense que non. »
\end{textebox}
\end{exercice}

\begin{exercice}[Remets les mots dans l'ordre]
Remets les mots suivants dans le bon ordre pour former des phrases correctes, puis traduis-les en français.

1. 一样 / 我 / 跟 / 高 / 不 / 他

2. 一方面 / 喜欢 / 我 / 另一方面 / 汉语 / 难 / 很 / 觉得 / ，

3. 如果 / 可是 / 努力 / 不练习 / 你 / 会 / 忘记 / 你 / ，

4. 足球 / 一样 / 喀麦隆的 / 中国的 / 和 / 有名 / 吗 / ？
\end{exercice}

\courssub{Je me prépare au Bac}

\begin{exercice}[Compréhension de texte et questions de langue]
Lis le texte suivant, puis réponds aux questions.

\begin{textebox}[Deux villes, deux vies]
雅温得跟北京不一样。一方面，雅温得比较安静，人也不太多；另一方面，北京很大，人很多，生活很快。雅温得的天气跟北京的天气也不一样：雅温得不太冷，北京的冬天很冷。如果你想学汉语，但是你不能去中国，没有关系：在喀麦隆也有很多好老师。一方面你可以在学校学习，另一方面你可以跟中国朋友练习说话。两个城市不一样，可是我都喜欢。\\
\emph{Yǎwēndé gēn Běijīng bù yīyàng. Yì fāngmiàn, Yǎwēndé bǐjiào ānjìng, rén yě bú tài duō; lìng yì fāngmiàn, Běijīng hěn dà, rén hěn duō, shēnghuó hěn kuài. Yǎwēndé de tiānqì gēn Běijīng de tiānqì yě bù yīyàng: Yǎwēndé bú tài lěng, Běijīng de dōngtiān hěn lěng. Rúguǒ nǐ xiǎng xué Hànyǔ, dànshì nǐ bù néng qù Zhōngguó, méiyǒu guānxi: zài Kāmàilóng yě yǒu hěn duō hǎo lǎoshī. Yì fāngmiàn nǐ kěyǐ zài xuéxiào xuéxí, lìng yì fāngmiàn nǐ kěyǐ gēn Zhōngguó péngyou liànxí shuōhuà. Liǎng ge chéngshì bù yīyàng, kěshì wǒ dōu xǐhuan.}\\
« Yaoundé n'est pas pareille à Pékin. D'une part, Yaoundé est assez calme et il n'y a pas trop de monde ; d'autre part, Pékin est grande, très peuplée, et la vie y est rapide. Le temps de Yaoundé n'est pas pareil non plus à celui de Pékin : à Yaoundé il ne fait pas très froid, l'hiver à Pékin est très froid. Si tu veux apprendre le chinois mais que tu ne peux pas aller en Chine, ce n'est pas grave : au Cameroun aussi il y a beaucoup de bons professeurs. D'une part tu peux étudier à l'école, d'autre part tu peux t'entraîner à parler avec des amis chinois. Les deux villes sont différentes, mais je les aime toutes les deux. »
\end{textebox}

\textbf{I. Questions de compréhension} (réponds en français)

1. Quelles sont les deux différences entre Yaoundé et Pékin mentionnées dans le texte ?

2. Que propose l'auteur à quelqu'un qui ne peut pas aller en Chine ?

3. Releve dans le texte les trois structures étudiées dans cette leçon et recopie les phrases complètes qui les contiennent.

\textbf{II. Questions de langue}

4. Dans la phrase \zh{雅温得跟北京不一样。}, pourquoi \zh{不} est-il placé avant \zh{一样} et non avant \zh{跟} ? Cite la règle.

5. Transforme la phrase \zh{北京很大，人很多。} en une phrase comparative avec \zh{……比……}.

6. Complète : 北京的天气\_\_\_雅温得的天气\_\_\_一样冷。(utilise \zh{跟} et \zh{不} si nécessaire)
\end{exercice}

\begin{exercice}[Thème et version]
\textbf{I. Thème.} Traduis les phrases suivantes en chinois (caractères et pinyin).

1. Mon frère est aussi grand que mon père.

2. D'une part j'aime le chinois, d'autre part je le trouve difficile.

3. Si tu étudies, mais que tu ne pratiques pas, tu oublieras.

4. Le français n'est pas aussi difficile que le chinois.

\textbf{II. Version.} Traduis le texte suivant en français.

\begin{textebox}[Version]
我的汉语水平跟我的同学不一样。一方面，我的口语比较好；另一方面，我的汉字写得不太好。如果我每天练习，但是我只写不说话，我的口语就不会进步。\\
\emph{Wǒ de Hànyǔ shuǐpíng gēn wǒ de tóngxué bù yīyàng. Yì fāngmiàn, wǒ de kǒuyǔ bǐjiào hǎo; lìng yì fāngmiàn, wǒ de Hànzì xiě de bú tài hǎo. Rúguǒ wǒ měitiān liànxí, dànshì wǒ zhǐ xiě bù shuō, wǒ de kǒuyǔ jiù bú huì jìnbù.}
\end{textebox}
\end{exercice}

\begin{exercice}[Expression écrite : ma ville et une ville chinoise]
\textbf{Sujet.} Tu participes à un échange scolaire entre ton lycée et un lycée chinois. Rédige en chinois un court paragraphe (au moins 6 phrases, environ 60 à 80 caractères) dans lequel tu compares ta ville (Yaoundé, Douala, Buea, etc.) et une ville chinoise de ton choix (Pékin, Shanghai, etc.).

\textbf{Contraintes obligatoires :}
\begin{itemize}
\item utilise au moins une fois la structure \zh{……跟/和……（不）一样} ;
\item utilise au moins une fois la structure \zh{一方面……，另一方面……} ;
\item utilise au moins une fois la structure \zh{如果……，但是/可是……} ;
\item soigne l'ordre des mots et la place de \zh{不} ;
\item écris les caractères correctement, sans pinyin dans la rédaction finale.
\end{itemize}

\textbf{Aide au démarrage :} tu peux commencer par \zh{我的城市跟上海不一样。} puis comparer la taille, le temps qu'il fait, la vie quotidienne et la nourriture.
\end{exercice}

\newpage

\coursec{Corrigés des exercices}

\begin{corrige}[Exercice 1 — QCM : la bonne structure]

\textbf{1. Réponse A :} \zh{我的汉语水平跟你的不一样。} \emph{(Wǒ de Hànyǔ shuǐpíng gēn nǐ de bù yīyàng.)} « Mon niveau de chinois n'est pas le même que le tien. » La négation \zh{不} se place obligatoirement devant \zh{一样}, jamais devant \zh{跟}. B est impossible (\zh{也} ne s'emploie pas ici), C est impossible (\zh{很} ne peut pas précéder \zh{一样} dans cette structure car \zh{一样} n'est pas graduable), D donnerait une phrase affirmative sans sens attendu ici.

\textbf{2. Réponse B :} \zh{我跟他不一样高。} \emph{(Wǒ gēn tā bù yīyàng gāo.)} « Je ne suis pas aussi grand que lui. » Schéma : A + \zh{跟} + B + \zh{不一样} + adjectif. A est fausse (\zh{不} ne peut pas précéder \zh{跟}), C est fausse (\zh{不} ne peut pas s'intercaler entre \zh{一样} et l'adjectif), D est fausse (ordre des mots bouleversé).

\textbf{3. Réponse C :} \zh{如果我想去中国留学，但是我的钱不够。} \emph{(Rúguǒ wǒ xiǎng qù Zhōngguó liúxué, dànshì wǒ de qián bù gòu.)} « Si je veux aller étudier en Chine, mais que je n'ai pas assez d'argent… » La structure \zh{如果……，但是……} exprime une condition assortie d'un obstacle. A et B expriment deux aspects parallèles (\zh{一方面……，另一方面……}), D une cause et une conséquence (\zh{因为……，所以……}) : aucune ne convient au sens.

\textbf{4. Réponse B :} \zh{如果你每天都学习，但是你不练习说话，你的口语就不会进步。} \emph{(Rúguǒ nǐ měitiān dōu xuéxí, dànshì nǐ bù liànxí shuōhuà, nǐ de kǒuyǔ jiù bú huì jìnbù.)} « Si tu étudies chaque jour mais que tu ne t'entraînes pas à parler, ton expression orale ne progressera pas. »
\end{corrige}

\begin{corrige}[Exercice 2 — Vrai ou faux ? Justifie ta réponse]

\textbf{1. Correct.} \zh{雅温得跟北京不一样大。} \emph{(Yǎwēndé gēn Běijīng bù yīyàng dà.)} « Yaoundé n'est pas aussi grande que Pékin. » La structure A + \zh{跟} + B + \zh{不一样} + adjectif est respectée : \zh{不} précède \zh{一样}, et l'adjectif \zh{大} suit.

\textbf{2. Incorrect.} La structure figée est \zh{一方面……，另一方面……} : on ne peut pas dire « 二方面 ». Correction : \zh{一方面我喜欢汉语，另一方面我觉得汉字很难。} \emph{(Yī fāngmiàn wǒ xǐhuan Hànyǔ, lìng yī fāngmiàn wǒ juéde Hànzì hěn nán.)} « D'une part j'aime le chinois, d'autre part je trouve les caractères très difficiles. » Le second membre s'ouvre par \zh{另一方面} (qui intègre l'idée de « l'autre »).

\textbf{3. Correct.} \zh{我的弟弟和我一样高。} \emph{(Wǒ de dìdi hé wǒ yīyàng gāo.)} « Mon petit frère est aussi grand que moi. » Structure affirmative A + \zh{和} + B + \zh{一样} + adjectif, sans adverbe entre \zh{一样} et l'adjectif.

\textbf{4. Incorrect.} En chinois, contrairement au français « si… alors… », on ne combine pas \zh{如果} avec \zh{所以} (qui traduit « donc »). Deux corrections possibles : \zh{如果你学习，你就会成功。} \emph{(Rúguǒ nǐ xuéxí, nǐ jiù huì chénggōng.)} « Si tu étudies, tu réussiras » (conséquence), ou \zh{因为你学习，所以你成功。} \emph{(Yīnwèi nǐ xuéxí, suǒyǐ nǐ chénggōng.)} « Parce que tu étudies, tu réussis » (cause-conséquence).
\end{corrige}

\begin{corrige}[Exercice 3 — Vocabulaire : complète le tableau]

\begin{center}
\begin{tabularx}{\linewidth}{|X|X|X|X|}
\hline
\rowcolor{popblueL} Caractères & Pinyin & Nature & Traduction \\
\hline
一样 & yīyàng & adjectif verbal / adverbe & pareil, de la même façon \\
\hline
比较 & bǐjiào & verbe / adverbe & comparer ; assez, relativement \\
\hline
方面 & fāngmiàn & nom & aspect, côté \\
\hline
如果 & rúguǒ & conjonction & si \\
\hline
可是 & kěshì & conjonction & mais \\
\hline
水平 & shuǐpíng & nom & niveau \\
\hline
\end{tabularx}
\end{center}

\textbf{Points de vigilance :} attention au ton de \zh{比较} \emph{bǐjiào} (3\textsuperscript{e} + 4\textsuperscript{e} ton) et à \zh{水平} \emph{shuǐpíng} (3\textsuperscript{e} + 2\textsuperscript{e} ton). Ne pas confondre \zh{可是} et \zh{但是} : tous deux signifient « mais », \zh{可是} est un peu plus oral, \zh{但是} plus neutre/écrit.
\end{corrige}

\begin{corrige}[Exercice 4 — Associe les caractères au pinyin]

\begin{center}
\begin{tabularx}{\linewidth}{|X|X|X|}
\hline
\rowcolor{popblueL} Caractères & Pinyin & Traduction \\
\hline
1. 水平 & c. shuǐpíng & niveau \\
\hline
2. 比较 & d. bǐjiào & comparer ; assez \\
\hline
3. 方面 & a. fāngmiàn & aspect, côté \\
\hline
4. 留学 & b. liúxué & étudier à l'étranger \\
\hline
5. 进步 & e. jìnbù & progresser ; progrès \\
\hline
\end{tabularx}
\end{center}

Réponses : \textbf{1-c, 2-d, 3-a, 4-b, 5-e}. Astuce mémorielle : \zh{进步} se compose de \zh{进} « avancer » et \zh{步} « pas » : « avancer pas à pas », donc « progresser ».
\end{corrige}

\begin{corrige}[Exercice 5 — Traduis en chinois]

\textbf{1.} \zh{我的哥哥和我的爸爸一样高。} \emph{(Wǒ de gēge hé wǒ de bàba yīyàng gāo.)} Structure affirmative : A + \zh{和} + B + \zh{一样} + adjectif. On accepte aussi \zh{我哥哥跟我爸爸一样高。}

\textbf{2.} \zh{一方面我喜欢汉语，另一方面我觉得汉语很难。} \emph{(Yī fāngmiàn wǒ xǐhuan Hànyǔ, lìng yī fāngmiàn wǒ juéde Hànyǔ hěn nán.)} Le second membre commence par \zh{另一方面}.

\textbf{3.} \zh{如果你学习，但是你不练习，你很快就会忘记。} \emph{(Rúguǒ nǐ xuéxí, dànshì nǐ bù liànxí, nǐ hěn kuài jiù huì wàngjì.)} On accepte \zh{可是} à la place de \zh{但是}.

\textbf{4.} \zh{Ndolé 跟中国的米饭不一样。} \emph{(Ndolé gēn Zhōngguó de mǐfàn bù yīyàng.)} « Le ndolé et le riz chinois ne sont pas pareils. » Pour un plat camerounais, on peut garder le nom « Ndolé » en lettres latines ou écrire \zh{喀麦隆的恩多莱}.

\textbf{5.} \zh{冬天，杜阿拉跟北京不一样冷。} \emph{(Dōngtiān, Dùālā gēn Běijīng bù yīyàng lěng.)} « En hiver, Douala n'est pas aussi froide que Pékin. » Le complément de temps \zh{冬天} se place en tête de phrase ; \zh{不} précède \zh{一样}.
\end{corrige}

\begin{corrige}[Exercice 6 — Transformation de phrases]

\textbf{1.} \zh{他比我高。} \emph{(Tā bǐ wǒ gāo.)} signifie « Il est plus grand que moi », donc moi et lui n'avons pas la même taille : il faut la forme \textbf{négative}. \zh{他跟我（的个子）不一样高。} ou mieux : \zh{我跟他不一样高。} \emph{(Wǒ gēn tā bù yīyàng gāo.)} « Je ne suis pas aussi grand que lui. » Justification : la structure \zh{跟……一样} exprime l'égalité ; comme \zh{比} exprime ici une supériorité, seule la forme négative \zh{不一样} conserve le sens.

\textbf{2.} \zh{法语和英语不一样难。} \emph{(Fǎyǔ hé Yīngyǔ bù yīyàng nán.)} « Le français et l'anglais ne sont pas aussi difficiles l'un que l'autre » : l'un des deux est plus difficile. Avec \zh{比} (en gardant le même sens, deux possibilités) : \zh{法语比英语难。} \emph{(Fǎyǔ bǐ Yīngyǔ nán.)} « Le français est plus difficile que l'anglais », ou l'inverse \zh{英语比法语难。} \emph{(Yīngyǔ bǐ Fǎyǔ nán.)} Les deux sont acceptées car la phrase de départ ne précise pas lequel est plus difficile.

\textbf{3.} \zh{一方面我想去中国，另一方面我没有钱。} \emph{(Yī fāngmiàn wǒ xiǎng qù Zhōngguó, lìng yī fāngmiàn wǒ méiyǒu qián.)} « D'une part je veux aller en Chine, d'autre part je n'ai pas d'argent. » Les deux aspects contradictoires sont mis en parallèle ; la virgule chinoise \zh{，} sépare les deux membres.
\end{corrige}

\begin{corrige}[Exercice 7 — Texte à trous]

Texte complété :

\zh{我的城市\textbf{跟}上海\textbf{不}一样大。\textbf{一方面}我的城市很安静，\textbf{另一方面}上海很热闹。\textbf{如果}你想学汉语，\textbf{可是}你不去中国，你也可以在喀麦隆学得很好。上海的生活\textbf{跟}我们这里的生活很\textbf{不}一样吗？我觉得\textbf{不}一样。}

\emph{Wǒ de chéngshì gēn Shànghǎi bù yīyàng dà. Yī fāngmiàn wǒ de chéngshì hěn ānjìng, lìng yī fāngmiàn Shànghǎi hěn rènao. Rúguǒ nǐ xiǎng xué Hànyǔ, kěshì nǐ bú qù Zhōngguó, nǐ yě kěyǐ zài Kāmàilóng xué de hěn hǎo. Shànghǎi de shēnghuó gēn wǒmen zhèlǐ de shēnghuó hěn bù yīyàng ma? Wǒ juéde bù yīyàng.}

\textbf{Justifications :}
\begin{itemize}
\item Trous 1 et 2 : structure A + \zh{跟} + B + \zh{不一样} + adjectif ; la traduction française (« ne sont pas aussi grandes ») impose la négation.
\item Trous 3 et 4 : deux aspects opposés (calme / animée) : \zh{一方面……，另一方面……}.
\item Trous 5 et 6 : condition + obstacle : \zh{如果……，可是……}.
\item Trous 7, 8 et 9 : question avec \zh{跟……一样吗} ; la réponse « Je pense que non » impose \zh{不一样} à la fin (le trou 8 reçoit donc \zh{不} : \zh{很不一样吗}).
\end{itemize}
\end{corrige}

\begin{corrige}[Exercice 8 — Remets les mots dans l'ordre]

\textbf{1.} \zh{我跟他不一样高。} \emph{(Wǒ gēn tā bù yīyàng gāo.)} « Je ne suis pas aussi grand que lui. » Ordre : sujet + \zh{跟} + B + \zh{不一样} + adjectif.

\textbf{2.} \zh{一方面我喜欢汉语，另一方面我觉得很难。} \emph{(Yī fāngmiàn wǒ xǐhuan Hànyǔ, lìng yī fāngmiàn wǒ juéde hěn nán.)} « D'une part j'aime le chinois, d'autre part je trouve ça très difficile. »

\textbf{3.} \zh{如果你努力，可是你不练习，你会忘记。} \emph{(Rúguǒ nǐ nǔlì, kěshì nǐ bù liànxí, nǐ huì wàngjì.)} « Si tu travailles dur mais que tu ne t'entraînes pas, tu oublieras. »

\textbf{4.} \zh{喀麦隆的足球和中国的足球一样有名吗？} \emph{(Kāmàilóng de zúqiú hé Zhōngguó de zúqiú yīyàng yǒumíng ma?)} « Le football camerounais est-il aussi célèbre que le football chinois ? » Le second \zh{足球} peut être sous-entendu : \zh{喀麦隆的足球和中国的一样有名吗？}
\end{corrige}

\begin{corrige}[Exercice 9 — Compréhension de texte et questions de langue]

\textbf{Barème indicatif (sur 20) :} I.1 : 3 pts ; I.2 : 2 pts ; I.3 : 3 pts ; II.4 : 4 pts ; II.5 : 4 pts ; II.6 : 4 pts.

\textbf{I. Questions de compréhension}

\textbf{1.} Deux différences sont mentionnées : d'une part l'ambiance et la taille (Yaoundé est assez calme et peu peuplée, Pékin est grande, très peuplée, avec un rythme de vie rapide) ; d'autre part le climat (à Yaoundé il ne fait pas très froid, alors que l'hiver à Pékin est très froid).

\textbf{2.} L'auteur propose d'apprendre le chinois au Cameroun : d'une part en étudiant à l'école, d'autre part en s'entraînant à parler avec des amis chinois, car il y a aussi beaucoup de bons professeurs au Cameroun.

\textbf{3.} Les trois structures de la leçon :
\begin{itemize}
\item \zh{……跟……（不）一样} : \zh{雅温得跟北京不一样。} (et \zh{雅温得的天气跟北京的天气也不一样})
\item \zh{一方面……，另一方面……} : \zh{一方面，雅温得比较安静……；另一方面，北京很大……} (et \zh{一方面你可以在学校学习，另一方面你可以跟中国朋友练习说话。})
\item \zh{如果……，但是……} : \zh{如果你想学汉语，但是你不能去中国，没有关系……}
\end{itemize}

\textbf{II. Questions de langue}

\textbf{4.} Dans la structure A + \zh{跟/和} + B + \zh{（不）一样}, \zh{跟} est une préposition (« avec ») qui introduit le terme de comparaison : on ne peut pas la nier. C'est l'adjectif verbal \zh{一样} (« pareil ») qui est nié, donc \zh{不} se place immédiatement devant \zh{一样}. Dire \zh{不跟北京一样} serait une erreur typique de francophone calquant « ne pas être pareil à ».

\textbf{5.} \zh{北京比雅温得大，人也比雅温得多。} \emph{(Běijīng bǐ Yǎwēndé dà, rén yě bǐ Yǎwēndé duō.)} ou plus simplement : \zh{北京比雅温得大。} \emph{(Běijīng bǐ Yǎwēndé dà.)} « Pékin est plus grande que Yaoundé. » Schéma : A + \zh{比} + B + adjectif.

\textbf{6.} \zh{北京的天气跟雅温得的天气不一样冷。} \emph{(Běijīng de tiānqì gēn Yǎwēndé de tiānqì bù yīyàng lěng.)} « Le temps de Pékin n'est pas aussi froid que celui de Yaoundé » (la structure exige \zh{跟} + \zh{不一样} + adjectif).

\textbf{Points de vigilance :} ne jamais placer \zh{不} devant \zh{跟} ; ne pas ajouter \zh{很} entre \zh{一样} et l'adjectif ; dans \zh{一方面……，另一方面……}, le second membre prend obligatoirement \zh{另} (réalisé par \zh{另方面} \rightarrow \zh{另一方面}).
\end{corrige}

\begin{corrige}[Exercice 10 — Thème et version]

\textbf{Barème indicatif (sur 20) :} thème 4 × 2,5 pts = 10 pts ; version 10 pts (exactitude du sens 6 pts, qualité du français 4 pts).

\textbf{I. Thème}

\textbf{1.} \zh{我的哥哥和我的爸爸一样高。} \emph{(Wǒ de gēge hé wǒ de bàba yīyàng gāo.)}

\textbf{2.} \zh{一方面我喜欢汉语，另一方面我觉得汉语很难。} \emph{(Yī fāngmiàn wǒ xǐhuan Hànyǔ, lìng yī fāngmiàn wǒ juéde Hànyǔ hěn nán.)}

\textbf{3.} \zh{如果你学习，但是你不练习，你就会忘记。} \emph{(Rúguǒ nǐ xuéxí, dànshì nǐ bù liànxí, nǐ jiù huì wàngjì.)}

\textbf{4.} \zh{法语跟汉语不一样难。} \emph{(Fǎyǔ gēn Hànyǔ bù yīyàng nán.)} « Le français n'est pas aussi difficile que le chinois. » On accepte aussi la tournure avec \zh{比} : \zh{法语没有汉语那么难。} est à éviter à ce niveau ; privilégier la structure de la leçon.

\textbf{II. Version}

« Mon niveau de chinois n'est pas le même que celui de mon camarade de classe. D'une part, mon expression orale est assez bonne ; d'autre part, je n'écris pas très bien les caractères chinois. Si je m'entraîne tous les jours, mais que je ne fais qu'écrire sans parler, mon expression orale ne progressera pas. »

\textbf{Points de vigilance :} \zh{写得不太好} \emph{(xiě de bù tài hǎo)} se traduit par « ne pas écrire très bien » (verbe + \zh{得} + adjectif : complément de manière) ; \zh{只写不说话} \emph{(zhǐ xiě bù shuōhuà)} = « ne faire qu'écrire sans parler » ; \zh{不会进步} \emph{(bú huì jìnbù)} = « ne progressera pas » (valeur de futur).
\end{corrige}

\begin{corrige}[Exercice 11 — Expression écrite : ma ville et une ville chinoise]

\textbf{Barème indicatif (sur 20) :} respect des trois structures obligatoires 6 pts (2 pts chacune) ; correction grammaticale (ordre des mots, place de \zh{不}) 6 pts ; vocabulaire et caractères 4 pts ; cohérence et longueur (6 phrases minimum) 4 pts.

\textbf{Proposition de rédaction modèle :}

\zh{我的城市是杜阿拉，它跟上海不一样。一方面，杜阿拉很热，一年到头天气都差不多；另一方面，上海有四季，冬天很冷。杜阿拉没有上海那么大，可是人也很热情。如果你想来喀麦隆旅游，但是你怕热，你可以去杜阿拉的海边。杜阿拉的菜跟中国的菜也不一样，我们喜欢吃Ndolé。两个城市不一样，可是我都喜欢。}

\emph{Wǒ de chéngshì shì Dùālā, tā gēn Shànghǎi bù yīyàng. Yī fāngmiàn, Dùālā hěn rè, yì nián dào tóu tiānqì dōu chàbùduō; lìng yī fāngmiàn, Shànghǎi yǒu sì jì, dōngtiān hěn lěng. Dùālā méiyǒu Shànghǎi nàme dà, kěshì rén yě hěn rèqíng. Rúguǒ nǐ xiǎng lái Kāmàilóng lǚyóu, dànshì nǐ pà rè, nǐ kěyǐ qù Dùālā de hǎibiān. Dùālā de cài gēn Zhōngguó de cài yě bù yīyàng, wǒmen xǐhuan chī Ndolé. Liǎng ge chéngshì bù yīyàng, kěshì wǒ dōu xǐhuan.}

« Ma ville est Douala, elle n'est pas pareille à Shanghai. D'une part, Douala est chaude et le temps est à peu près le même toute l'année ; d'autre part, Shanghai a quatre saisons et l'hiver y est très froid. Douala n'est pas aussi grande que Shanghai, mais les gens y sont aussi chaleureux. Si tu veux venir en voyage au Cameroun mais que tu as peur de la chaleur, tu peux aller au bord de la mer à Douala. La cuisine de Douala n'est pas pareille non plus à la cuisine chinoise : nous aimons manger le ndolé. Les deux villes sont différentes, mais je les aime toutes les deux. »

\textbf{Points de vigilance :}
\begin{itemize}
\item Vérifier la place de \zh{不} : toujours devant \zh{一样} (\zh{跟上海不一样}), jamais devant \zh{跟}.
\item Le second membre de la structure parallèle commence par \zh{另一方面}, avec \zh{另} (ton \emph{yī} \rightarrow \emph{yí} devant \emph{miàn} 4\textsuperscript{e} ton).
\item Ne pas écrire \zh{如果……所以……} : avec \zh{如果}, on utilise \zh{但是/可是} (obstacle) ou \zh{就} (conséquence).
\item Pas de pinyin dans la copie finale ; soigner les caractères \zh{热} (clé \zh{灬}, 4 traits bas), \zh{情} (clé \zh{忄}, ordre : gauche avant droite), \zh{旅} (clé \zh{方}, haut avant bas).
\end{itemize}
\end{corrige}

\newpage

\coursec{Fiche bilan}

\courssub{Fiche Bilan — Leçon 20 : Comparer, Nuancer, Hypothétiser}

\begin{popfigure}
\legende{Carte mentale : Leçon 20 — Structures de comparaison, d'argumentation et d'hypothèse}
\begin{center}\tcbox[colback=popink,colframe=popink,coltext=white,arc=9pt,boxsep=4pt,left=6pt,right=6pt]{\bfseries\small Leçon 20}\end{center}
\vspace{-2pt}
\begin{tcbraster}[raster columns=3,raster equal height=rows,raster column skip=6pt,raster row skip=6pt,enhanced,blankest]
\begin{tcolorbox}[enhanced,colback=poporange,colframe=poporange,coltext=white,boxrule=0.9pt,arc=3pt,left=4pt,right=4pt,top=3pt,bottom=3pt,boxsep=1pt,halign=left]\bfseries\scriptsize A \cle{跟/和} B (\cle{不})\cle{一样} (+ Adj)\end{tcolorbox}
\begin{tcolorbox}[enhanced,colback=popgreen,colframe=popgreen,coltext=white,boxrule=0.9pt,arc=3pt,left=4pt,right=4pt,top=3pt,bottom=3pt,boxsep=1pt,halign=left]\bfseries\scriptsize \cle{一方面}… \cle{另一方面}…\end{tcolorbox}
\begin{tcolorbox}[enhanced,colback=poppurple,colframe=poppurple,coltext=white,boxrule=0.9pt,arc=3pt,left=4pt,right=4pt,top=3pt,bottom=3pt,boxsep=1pt,halign=left]\bfseries\scriptsize \cle{如果}… \cle{但是/可是}…\end{tcolorbox}
\begin{tcolorbox}[enhanced,colback=poppink,colframe=poppink,coltext=white,boxrule=0.9pt,arc=3pt,left=4pt,right=4pt,top=3pt,bottom=3pt,boxsep=1pt,halign=left]\bfseries\scriptsize Vocabulaire clé\end{tcolorbox}
\begin{tcolorbox}[enhanced,colback=popgold,colframe=popgold,coltext=white,boxrule=0.9pt,arc=3pt,left=4pt,right=4pt,top=3pt,bottom=3pt,boxsep=1pt,halign=left]\bfseries\scriptsize Pièges francophones\end{tcolorbox}
\begin{tcolorbox}[enhanced,colback=popblue,colframe=popblue,coltext=white,boxrule=0.9pt,arc=3pt,left=4pt,right=4pt,top=3pt,bottom=3pt,boxsep=1pt,halign=left]\bfseries\scriptsize Application écrite/orale\end{tcolorbox}
\begin{tcolorbox}[enhanced,colback=popmuted,colframe=popmuted,coltext=white,boxrule=0.9pt,arc=3pt,left=4pt,right=4pt,top=3pt,bottom=3pt,boxsep=1pt,halign=left]\bfseries\scriptsize Caractères à maîtriser\end{tcolorbox}
\end{tcbraster}

\end{popfigure}

\begin{bilan}

\end{bilan}

\courssub{Lexique clé (HSK 2–3)}
\begin{center}
\begin{tabularx}{\linewidth}{|X|X|X|X|}
\hline
\rowcolor{popblueL} \textbf{Caractères} & \textbf{Pinyin} & \textbf{Nature} & \textbf{Traduction} \\ \hline
比较 & bǐjiào & adv./v. & comparer ; relativement \\
跟 & gēn & prep. & avec (comparaison, action commune) \\
和 & hé & prep./conj. & et ; avec (comparaison) \\
一样 & yīyàng & adj. & identique, pareil \\
不一样 & bù yīyàng & adj. & différent \\
相同 & xiāngtóng & adj. & identique (plus formel) \\
不同 & bùtóng & adj. & différent (plus formel) \\
方面 & fāngmiàn & n. & aspect, côté \\
一方面 & yī fāngmiàn & loc. & d'un côté, en premier lieu \\
另一方面 & lìng fāngmiàn & loc. & d'un autre côté, par ailleurs \\
如果 & rúguǒ & conj. & si, au cas où \\
就 & jiù & adv. & alors, donc (corrélatif de 如果) \\
但是 & dànshì & conj. & mais, cependant \\
可是 & kěshì & conj. & mais, pourtant (plus oral) \\
优点 & yōudiǎn & n. & avantage, point fort \\
缺点 & quēdiǎn & n. & défaut, point faible \\
条件 & tiáojiàn & n. & condition \\
转折 & zhuǎnzhé & n./v. & tournant ; transition (contraste) \\ \hline
\end{tabularx}
\end{center}

\courssub{Règles de grammaire à connaître par cœur}
\begin{center}
\begin{tabularx}{\linewidth}{|>{\raggedright\arraybackslash}X|>{\centering\arraybackslash}X|>{\raggedright\arraybackslash}X|}
\hline
\rowcolor{popblueL} \textbf{Structure / Schéma} & \textbf{Exemple (Caractères + Pinyin)} & \textbf{Traduction / Note} \\ \hline
\textbf{1. Comparaison d'égalité / différence} \\
A + 跟/和 + B + 一样 / 不一样 & 我 \cle{跟} 他 \cle{一样} 高。\newline (Wǒ \cle{gēn} tā \cle{yīyàng} gāo.) & Je suis aussi grand que lui. \\
A + 跟/和 + B + (\cle{不})一样 + \textbf{Adj.} & 这件衣服 \cle{跟} 那件 \cle{不一样} \cle{贵}。\newline (Zhè jiàn yīfu \cle{gēn} nà jiàn \cle{bù yīyàng} \cle{guì}.) & Ce vêtement n'est pas aussi cher que celui-là (il est plus/moins cher). \\
\emph{Note :} \cle{跟} plus oral, \cle{和} plus écrit. L'adjectif se place \textbf{après} \cle{不一样}. & & \\ \hline
\textbf{2. Argumentation équilibrée} \\
\cle{一方面}…, \cle{另一方面}… & \cle{一方面} 学中文 \cle{很有用}，\cle{另一方面} \cle{很难}。\newline (\cle{Yī fāngmiàn} xué Zhōngwén \cle{hěn yǒuyòng}，\cle{lìng fāngmiàn} \cle{hěn nán}.) & D'un côté, apprendre le chinois est très utile, d'un autre côté, c'est très difficile. \\
\emph{Emploi :} Présenter deux aspects contradictoires ou complémentaires d'un même sujet. Pas de verbe entre les deux segments obligatoirement. & & \\ \hline
\textbf{3. Hypothèse + Contraste (Réalité)} \\
\cle{如果}…, \cle{但是/可是}… & \cle{如果} 你 \cle{不} 去，\cle{我} \cle{可是} 要 去。\newline (\cle{Rúguǒ} nǐ \cle{bù} qù, wǒ \cle{kěshì} yào qù.) & Si tu n'y vas pas, moi j'y vais, par contre. \\
\emph{Structure :} \cle{如果} (condition) + Sujet + \cle{就} (souvent omis si contraste fort) + \cle{但是/可是} (réalité opposée). \cle{可是} plus familier/insistant sur le contraste. & & \\ \hline
\end{tabularx}
\end{center}

\courssub{Méthodes express}
\begin{itemize}
    \item \textbf{Comparer A et B :} Identifier les deux termes $\rightarrow$ Choisir \cle{跟} (oral) ou \cle{和} (écrit) $\rightarrow$ Placer \cle{一样} (identique) ou \cle{不一样} + Adj (différence de degré).
    \item \textbf{Rédiger un paragraphe argumentatif :} Ouvrir avec \cle{一方面} + argument 1 $\rightarrow$ Pivoter avec \cle{另一方面} + argument 2 $\rightarrow$ Conclure par une synthèse (\cle{因此} / \cle{总的来说}).
    \item \textbf{Exprimer un choix conditionnel :} Poser l'hypothèse avec \cle{如果} + situation $\rightarrow$ Marquer la conséquence attendue avec \cle{就} (optionnel) $\rightarrow$ Introduire la réalité contraire avec \cle{但是/可是}.
    \item \textbf{Traduction Thème (Fr $\rightarrow$ Ch) :} « Bien que… mais… » $\neq$ \zh{虽然…但是…} (leçon précédente). Ici : « Si… (alors)… mais en fait… » $\rightarrow$ \cle{如果}… \cle{可是}…
\end{itemize}



\begin{aretenir}[Checklist avant le Bac — Leçon 20]
\begin{itemize}
    \item $\square$ Je sais écrire et ordonner les traits de : \zh{比, 较, 跟, 和, 样, 方, 面, 如, 果, 但, 是, 可}.
    \item $\square$ Je distingue \cle{跟} (oral/action commune) et \cle{和} (écrit/liste) dans la comparaison.
    \item $\square$ Je place l'adjectif \textbf{après} \cle{不一样} (ex: \zh{不一样贵} \checkmark / \zh{贵不一样} \textbf{✗}).
    \item $\square$ Je n'utilise \textbf{PAS} \cle{比} (bǐ) avec la structure \cle{跟/和...一样} (confusion HSK 2/3).
    \item $\square$ Je construis correctement : \cle{一方面}…, \cle{另一方面}… (virgule obligatoire, pas de « et » entre les deux).
    \item $\square$ J'utilise \cle{如果}…, \cle{就}… pour la conséquence logique, et \cle{如果}…, \cle{但是/可是}… pour le contraste réel.
    \item $\square$ Je choisis \cle{可是} pour l'oral/insistance sur le refus, \cle{但是} pour l'écrit/neutre.
    \item $\square$ Je sais transformer « A est différent de B » en \zh{A 跟 B 不一样} et « A n'est pas aussi Adj que B » en \zh{A 跟 B 不一样 Adj}.
    \item $\square$ Je peux comparer deux systèmes éducatifs (Cameroun/Chine) avec \cle{一方面}…\cle{另一方面}….
    \item $\square$ Je peux exprimer un projet conditionnel : \zh{如果明天不下雨，我们就去踢球，但是如果下雨，我们就在家看电视。}
    \item $\square$ Je relis mes phrases pour vérifier l'ordre : Sujet + \cle{跟/和} + Objet + \cle{(不)一样} (+ Adj).
    \item $\square$ Je connais le vocabulaire de la fiche (优点, 缺点, 条件, 转折) pour l'expression écrite.
\end{itemize}
\end{aretenir}

\findecours

\end{document}
